%for a more compact document, add the option openany to avoid
%starting all chapters on odd numbered pages
\documentclass[12pt]{cmuthesis}

% This is a template for a CMU thesis.  It is 18 pages without any content :-)
% The source for this is pulled from a variety of sources and people.
% Here's a partial list of people who may or may have not contributed:
%
%        bnoble   = Brian Noble
%        caruana  = Rich Caruana
%        colohan  = Chris Colohan
%        comar    = Cyrus Omar
%        jab      = Justin Boyan
%        josullvn = Joseph O'Sullivan
%        jrs      = Jonathan Shewchuk
%        kosak    = Corey Kosak
%        mjz      = Matt Zekauskas (mattz@cs)
%        pdinda   = Peter Dinda
%        pfr      = Patrick Riley
%        dkoes = David Koes (me)

% My main contribution is putting everything into a single class files and small
% template since I prefer this to some complicated sprawling directory tree with
% makefiles.

% some useful packages
\usepackage{fullpage}
\usepackage{graphicx}
\usepackage{amsmath}
\usepackage{amssymb}
\usepackage{amsthm}
\usepackage{mdwlist}
\usepackage[numbers,sort]{natbib}
\PassOptionsToPackage{hyphens}{url}
\usepackage[backref=page,pageanchor=true,plainpages=false, pdfpagelabels, bookmarks,bookmarksnumbered,
colorlinks=true,
citecolor=citec,
linkcolor=linkc,
urlcolor=urlc,
%pdfborder=0 0 0,  %removes outlines around hyper links in online display
]{hyperref}
\usepackage{subfigure}
\usepackage{xspace}
\usepackage[linesnumbered,ruled]{algorithm2e}
\usepackage{algpseudocode}
\usepackage{microtype}
% "nosort" causes no sort AND no compress of 3+ section references in one;
% it's "sort" to sort but NOT compress; nothing to.. compress but not sort, i guess? not the greatest naming
\usepackage[nameinlink,nosort]{cleveref}

%% custom fonts

%\usepackage{fontspec}
\usepackage{mathspec}
\usepackage{xeCJK}
\usepackage[xetex]{xechangebar}

%\setmainfont{XCharter}
\setallmainfonts(Digits,Latin){XCharter}
%\setmainfont{Utopia}[
%	Path = ./fonts/,
%	UprightFont = UtopiaStd-Regular.otf,
%	BoldFont = UtopiaStd-Semibold.otf,
%	ItalicFont = UtopiaStd-Italic.otf,
%	BoldItalicFont = UtopiaStd-SemiboldIt.otf
%	]

\setCJKmainfont{HeiseiMinStd-W5}[Path = ./fonts/]
%\newfontfamily\arabicfont{Droid Sans Arabic} % doesn't respond to Script = Arabic :(
%\newfontfamily\arabicfont{LateefGR-Regular}[Path = ./fonts/] % doesn't respond to Script = Arabic :(
\newfontfamily\arabicfont{Amiri-Regular}[Path = ./fonts/, Script = Arabic]

% Approximately 1" margins, more space on binding side
%\usepackage[letterpaper,twoside,vscale=.8,hscale=.75,nomarginpar]{geometry}
%for general printing (not binding)
\usepackage[letterpaper,twoside,vscale=.8,hscale=.75,nomarginpar,hmarginratio=1:1]{geometry}

% Provides a draft mark at the top of the document.
\draftstamp{\today}{DRAFT}

% NB. needs to be defined before mainmatter so they work in hyperref
\definecolor{citec}{RGB}{128,0,64}
\definecolor{linkc}{RGB}{0,64,128}
\definecolor{urlc} {RGB}{128,64,0}

% format section references with no space after the \S
% #1 is the ref number; the #2 and #3 here are the hyperlink boundaries
\crefformat{section}{#2\S{}#1#3}
\Crefformat{section}{#2\S{}#1#3}

\begin{document}
\frontmatter

%initialize page style, so contents come out right (see bot) -mjz
\pagestyle{empty}

\title{ %% {\it \huge Thesis Proposal}\\
{\bf Practical Concurrency Testing} \\
\normalsize \vspace{1em}
\begin{tabular}{c}
{\em or: How I Learned to Stop Worrying and Love the Exponential Explosion}
\end{tabular}}
\author{Ben Blum}
\date{October 2018}
\Year{2018}
\trnumber{}

\committee{
Garth Gibson, Chair \\
David A. Eckhardt \\
Brandon Lucia \\
Haryadi Gunawi, University of Chicago
}

\support{}
\disclaimer{}

% copyright notice generated automatically from Year and author.
% permission added if \permission{} given.

%\keywords{landslide terminal, baggage claim, ground transportation, ticketing}
\keywords{concurrency, testing, debugging, verification, model checking, data races, education, transactional memory}

\maketitle

\begin{dedication}
For my family, my teachers, and my students.
\end{dedication}

\pagestyle{plain} % for toc, was empty

%\newcommand\revision[1]{#1\xspace}
\newcommand\revision[1]{\cbstart[0.5em]{}#1\cbend\xspace}
%\newcommand\revisionminor[1]{#1\xspace}
\newcommand\revisionminor[1]{\cbstart{}#1\cbend\xspace}

\begin{abstract}
Concurrent programming presents a challenge to students and experts alike because of the complexity of multithreaded interactions and the difficulty \revisionminor{of reproducing and reasoning} about bugs.
Stateless model checking is a
%concurrency
testing approach which forces a program to interleave its threads in many different ways, checking for bugs each time.
This technique is powerful, in principle capable of finding any nondeterministic bug in finite time,
but suffers from exponential explosion as program size increases.
Checking an exponential number of thread interleavings is not a practical or predictable approach for programmers to find concurrency bugs before their project deadlines.

In this thesis, I develop several new techniques to make stateless model checking more practical for human use.
I have built Landslide, a stateless model checker specializing in undergraduate operating systems class projects.
Landslide extends the traditional model checking algorithm with a new framework for automatically managing
%the exploration of
multiple state spaces according to their estimated completion times,
which I show quickly finds bugs should they exist and also quickly verifies correctness otherwise.
I evaluate Landslide's suitability for inexpert use by presenting the results of many semesters providing it to students in 15-410, CMU's Operating System Design and Implementation class,
and more recently, students in similar classes at the University of Chicago and Penn State University.
%Finally, I will present several new techniques that allow stateless model checking to be practically employed on real-world programs.
%Finally, I will explore broader impact by extending Landslide to test some real-world programs and to be used by students at other universities.
Finally, I extend Landslide with a new concurrency model for hardware transactional memory,
and evaluate several real-world transactional benchmarks to show that stateless model checking can keep up with the developing concurrency demands of real-world programs.
\end{abstract}

\begin{acknowledgments}
\revision{I honor here the many who %made this journey of 7 long years possible.
supported me during these 7 long years and made this journey possible.

\subsubsection{Thesis}

Professor Garth Gibson is largely responsible for shaping me into the researcher I am today.
%Thanks to Garth's guidance,
%this thesis has grown into something that past-me,
%first starting research with the simple vision ``make a better autograder for 410'',
%could not have then imagined.
Despite having primary research interests in other fields,
Garth enthusiastically took on the project,
and always pushed me to explore new problem domains,
to design scientifically thorough experiments,
to explain difficult concepts approachably,
to appreciate related work as charitably as possible,
and to seek guidance from industry people with relevant experience,
despite my often stubborn refusal to leave my own little comfort zone.
So, even when it took me sending 10+ email reminders to get your attention,
truly,
thanks for everything, Garth.

Professor David A. Eckhardt
provided unending guidance on education, writing, career, and life in general, %life, you name it,
helped revise each semester's %iteration of the
recruiting lecture with a keen eye for the student mind,
contributed several clever survey questions, %substantially to the design of the survey,
and supplied constant %enthusiasm and
encouragement about Landslide's value to 15-410, %that Landslide was a boon to 15-410,
not to mention the red-ink-encrusted printed draft of this document I received immediately after my defense.
My debt of thanks to Dave is of such magnitude that I may never hope to pay it back, but perhaps forward instead.
%I can't express my gratitude to Dave strongly enough,
%even if it took until just this week for half of the Landslide-friendly tests to be included in the official P2 test suite.
%even though three of the Landslide-friendly tests never made it upstream into the official P2 test suite.

Ji\v{r}\'{i} \v{S}im\v{s}a, another former student of Garth's,
first sparked my interest in the field,
and then gave invaluable support during my first few years.
Ji\v{r}\'{i}'s mentorship
oriented me in the world of research and helped me find my bearings as a grad student,
from teaching me how DPOR works
to helping revise my paper drafts
to offering the kind of candid career advice you can't get from a professor.
Had I not met Ji\v{r}\'{i} and learned about dBug during grad school applications way back,
none of this would have happened.

{Many have collaborated with me directly on the research itself,
whom I thank here in approximate chapter order.
Michael J. Sullivan and Nathaniel Wesley Filardo were Landslide's first other users
%besides myself
back in its prototype days.
Joshua Wise, Miriam Zimmerman, and Glenn Willen helped revise my explanation of DPOR in \Cref{chap:landslide}.
Professor Brandon Lucia provided valuable revisions to \Cref{chap:quicksand}'s content
that ultimately led to its publication at OOPSLA 2016.
Michael J. Sullivan double-checked the proofs therein.
Professor Timothy Zhu from Penn State University
and Professor Haryadi Gunawi and TA Kevin Zhao from the University of Chicago
graciously allowed me to use their students as guinea pigs
and helped to improve Landslide's
%stability and
robustness for use beyond CMU's walls,
greatly enriching \Cref{chap:education}'s evaluation.
David Blum and Skye Toor provided
invaluable advice towards the statistical analysis
used therein.
%that helped me evaluate the student
%grade distributions therein.
Mario Dehesa-Azuara generously contributed the transactional
% justify last line bc of manual page break
\unskip\parfillskip 0pt \par}

% FIXME - cmon latex pls??
\pagebreak

\noindent
memory data structures
that make up about half of \Cref{chap:tm}'s evaluation suite.
%
Carlo Angiuli, Evan Cavallo, Ziv Scully, Jim McCann,
Matthew Maurer, Stefan Muller, Guillaume Didier,
Sol Boucher, and Michael J. Sullivan
helped me rehearse and polish the snake fight talk.

Most importantly of all, I thank every student who ever used Landslide.
Though the introduction section may say differently,
their smiles and gratitude were the true motivation of this thesis.

\subsubsection{Academia}

The community of 15-410 TAs
%having surrounded me pretty much since starting undergrad,
has nurtured my love of teaching
since well before I took the class myself.
%and well after graduation.
%As skilled as the instructor is,
15-410 could not be the world's best OS class
without %TAs
course staff's %(staves)
unwavering dedication to
thorough grading, the Socratic method,
and empathy to struggling students.
%as much
%from year to year as the crew I have known.
It was a blessing to serve, and later to do research, among you all.
Thanks especially to Stanley Zhang,
%the 2018 TA who took on the mantle of ensuring
the 2018 TA who who worked with me to ensure
15-410 could keep grading projects with Landslide even after I leave.
Thanks to Wind River for generously
providing
%CMU with
Simics educational licenses.
%I am also grateful
Thanks also to Adam Blank %and Miriam Zimmerman % nah, she said it was just adam
for founding and organizing the class Great Practical Ideas,
which teaches oft-neglected programming skills such as version control and debugging,
which I had the opportunity to teach as head instructor one semester.

I thank the Parallel Data Lab, my research group at CMU,
for supporting me at the PDL retreat every fall,
pushing me to refine my presentation skills and to network with strangers (now familiar faces) from industry.
%I also thank those %industry
I also thank the industry attendees
%from each of the PDL Consortium member companies
for their interest %, even in more theoretical work,
%and for the enlightening conversations resulting therefrom.
and enlightening conversations.
Joan Digney and Karen Lindenfelser
make a lot of gears turn behind the scenes at PDL,
and Deb Cavlovich and Catherine Copetas
do likewise for CSD as a whole.
%Thanks to the anonymous OOPSLA reviewer who went the extra mile to read my supplementary proofs
%and urged me to improve the formalism to be more foundational
%(the lack of any stronger criticism a silent vote of confidence for their overall correctness).
Thanks to Nicholas D. Matsakis for the mentorship %and collaboration
during my two summers as a Rust research intern at Mozilla,
and to the rest of the Rust community as well for their %extremely important
work
bringing better concurrent programs into the world. %from a different angle.
And to everyone at large who
%outside my %immediate
%research circle who
%whether at the retreat or otherwise,
insisted how cool it was that my research helps students:
thank you for helping me keep my head up during darker times.

%I am grateful to
The ThursDz Council,
my group of now-former grad student, Shadyside-bound (in spirit if not in person) restaurant-hoppers,
surrounded me with both friendship and mentorship over the years.
In particular,
%in alphabetical for lack of a better order,
Chris Martens,
Jason Reed,
Jim McCann,
Rob Simmons,
Tom Murphy 7,
and
William Lovas:
you each have been role models to me
perhaps more than you know.

Without SIGBOVIK's unique brand of humorous, self-aware, and {\em definitely legitimate} research,
I probably would not have been interested in grad school in the first place,
or perhaps would have grown too frustrated along the way for lack of a cathartic outlet for conference paper woes.
I thank all who've worn the mantle of Harry Bovik over the years,
from the conference's original founders,
to every paper author,
to all future organizers who'll keep those ``mainstream'' conferences on their toes well into the future.

\subsubsection{Life}

My parents, Eve and David Blum, and my grandparents, Margie Granach and Elsie Blum,
provided bottomless emotional support throughout my education.
I appreciate beyond words for the privileged life they provided
that allowed me the opportunity to pursue higher education.

Looking back on my 11 years at CMU,
I am especially grateful for the friendships made early on %in my undergrad years
that continued well past first graduation.
From those who stayed in Pittsburgh %with me
to those who %moved away but
kept in touch,
from my regular tabletop, board, and/or dance gaming partners
to those who %inexhaustibly
welcomed me with company and
a couch to sleep on during travel,
I thank
%-- and in any order other than alphabetical here would be far too cruel --
8 Gianfortoni,
Adam Blank,
Alan Vangpat,
Alex Yuschik,
Anand Subramanian,
Andrew Krieger,
Car Bauer,
Carolyn Sawyer,
Elizabeth Kemp,
Elly Fong-Jones,
Emily Forney,
Emily Leathers,
Emma Cating-Subramanian,
Eric Faust,
Gabe Routh,
Greg Hanneman,
Hazel Vird\'{o},
Jack Ferris,
Jake Lengyel,
Jason Deng,
Josh Keller,
Joshua Wise,
Josiah Boning,
Julia Tuttle,
Kartik Subramanian,
Kellie ``K2'' Medlin,
Laura Abbott,
Maija Mednieks,
Margaret Meyerhofer,
Michael A\texttt{rntz}enius,
Michael J. ``Sully'' Sullivan,
Miriam Zimmerman,
Naomi Saphra,
Nathaniel Wesley Filardo,
Rhett Lauffenburger,
Richard Ha,
Rory Glenn,
Ryan Pearl,
and
Todd Eisenberger
for their wonderful friendship.
I am equally grateful to the new friends I made during my graduate years,
whether through board/card/dance gaming, tea drinking, %and/or
movie watching,
and/or as more permanent Pittsburgh residents that I am extremely torn up about moving away from:
%leaving behind (for now at least):
Alexandra Lee Falk,
Alexis Dyer,
Barbara Jensen,
Brendan McShane,
Brian Saghy,
Carlo Angiuli,
Cassie Orr,
Dan Guzek,
Danny Gratzer,
David Renshaw,
Evan Cavallo,
Gabriele Farina,
Grant Wu,
Jenny Lin,
Karl Schulz,
Kevin Nguyen,
Kristy Gardner,
Lea Albaugh,
Marlena Abraham,
Matt Stanec,
Matthew Maurer,
Sarah R. Allen,
Sol Boucher,
Stefan Muller,
Zachary Waldman,
Ziv Scully,
and the rest of ThursDz already named previously.
And to anyone who may read this list
and feel even slightly dejected that their name was left out:
I owe you a nice dinner. Get in touch.

% I wrote the below myself, but constructed it carefully to not just convey my feelings accurately
% but also produce intelligible output through google translate, for the sake of 日本語のできない読者の方々。
% (The translation is not perfect, like 達成感 means more my personal sense of accomplishment,
% rather than sounding objective (and consequently, kinda pompous) as it does in google's ouput, but it ain't bad.)
% Of course, since google translate is powered by neural nets and is under active research and development,
% I copy the current output here in case it changes in the future.
%
%%%% I would like to express my sincere gratitude to my fellow learners and native speakers
%%%% who have studied Japanese, talked and taught me variously.
%%%% I think that learning a foreign language is the second most successful achievement of my life.
%
% My own (loose) translation would be more like,
% I am forever grateful to the friends I made through Japanese,
% both my companions in language learning and the native speakers I was able to meet and befriend to begin with,
% who studied together with me, carried on conversations with me,
% and taught me a great variety of things.
% I consider learning a foreign language (*nuance of ongoing process rather than completed achievement)
% to be the second greatest sense of accomplishment I have felt in my life.
日本語を
一緒に
勉強したり、
会話したり、
色々教えてくださったりした
仲間の学習者さんたちと
ネイティブスピーカーさんたちへ、
心から感謝を申し上げます。
外国語を学ぶのは
一生の二番
%充実感
達成感
だと考えます。

Lastly,
the Android: Netrunner community has become such an important part of my life
over the last half decade or so.
%for the better part of my time as a grad student.
Netrunner is a deckbuilding strategy card game
that depicts cyberspace hacking battles (completely different from real-life computer security)
alongside dystopian social issues (becoming strangely more familiar with each passing year).
Its characters have a diversity of ethnicity and gender identity %and sexual orientation
unparalleled in any other game I know of,
and the community that has grown around it reflects those values.
I thank especially the other Stimhack moderators,
the members of {\sf \#rainbow-coalition},
and the local Pittsburgh crew
for their support, encouragement,
and mutual exchange of inclusive values in these turbulent yet hopeful times.
%and
%friendships very real and deep, even when forged only in text,
%%\footnote{and Slack emojis},
%that will long outlive the silly card game that brought us together.
%%It'd be an injustice to list only a few names here and omit others, but you know who you are.
}

\end{acknowledgments}

\tableofcontents
\listoffigures
\listoftables

\mainmatter

%% Double space document for easy review:
%\renewcommand{\baselinestretch}{1.66}\normalsize

% The other requirements Catherine has:
%
%  - avoid large margins.  She wants the thesis to use fewer pages,
%    especially if it requires colour printing.
%
%  - The thesis should be formatted for double-sided printing.  This
%    means that all chapters, acknowledgements, table of contents, etc.
%    should start on odd numbered (right facing) pages.
%
%  - You need to use the department standard tech report title page.  I
%    have tried to ensure that the title page here conforms to this
%    standard.
%
%  - Use a nice serif font, such as Times Roman.  Sans serif looks bad.
%
% Other than that, just make it look good...

\newcommand\simics{\textsc{Simics}}
\newcommand{\sect}[1]{\S #1}
\newcommand\hilight[2]{\color{#1}#2\color{black}}
\definecolor{orange}{RGB}{192,96,0}
\definecolor{olivegreen}{RGB}{0,127,0}
\definecolor{brickred}{RGB}{192,0,0}
\definecolor{commentblue}{RGB}{0,0,192}
% colors from sigbovik for tsx/etc
% threads
\definecolor{lavender}{RGB}{102,65,208} % (?? i don't know what happened here anymore but it looks good)
\definecolor{seafoam}{RGB}{56,158,68} % * 2/3
\definecolor{salmon}{RGB}{208,89,74} % * 7/8
\definecolor{goldish}{RGB}{128,96,32}
% eval data
\definecolor{darkcyan}{RGB}{0,128,128}
\definecolor{lime}{RGB}{48,128,0}
\definecolor{darklavender}{RGB}{68,44,138}
\definecolor{brownish}{RGB}{128,96,0}
\definecolor{pinkish}{RGB}{128,34,102} % ff4488
% syntax hilighting
\newcommand\flow[1]{\hilight{brownish}{#1}\xspace}
\newcommand\const[1]{\hilight{brickred}{#1}\xspace}
\newcommand\call[1]{\hilight{pinkish}{#1}}
\newcommand\ctype[1]{\hilight{lime}{#1}\xspace}
\newcommand\ccomment[1]{\hilight{darkcyan}{#1}\xspace}

\newtheorem{lemma}{Lemma}
\newtheorem{theorem}{Theorem}
\newtheorem{definition}{Definition}

\newcommand\llama[1]{\includegraphics[width=#1]{llama.pdf}}
\newcommand\llitem{\item[\raisebox{-0.15em}{\llama{1.2em}}]}

\newcommand{\inspirationallinebreak}{\vspace{0.25em}}
\newcommand{\inspirationalhyphen}{-\hspace{-0.15em}-\xspace}
\newcommand\inspirationalquote[2]{\begin{flushright}
	\vspace{-1em}
	{
	\fontspec[Path=./fonts/]{AlexaStd}
	{\em #1}
	\inspirationallinebreak

	-\hspace{-0.15em}-\hspace{-0.15em}-{#2}
	}
	\vspace{2em}
\end{flushright}}

\chapter{Introduction}
\inspirationalquote{
\begin{tabular}{p{0.7\textwidth}}
Please hold on. This thesis will be departing shortly for the Landslide terminal, baggage claim, ground transportation, and ticketing.
\end{tabular}}
{Pittsburgh International Airport (paraphrased)}

\revision{Pick up any recent concurrency-related research paper, %from the last few decades,
and chances are its very first sentence will tell you
%how notoriously difficult concurrency bugs are compared to sequential ones
that concurrency bugs are notoriously difficult to find, reproduce, diagnose, and/or fix
% ,
%and/or thereafter verify fixed,
%compared to sequential ones,
%on account of the multitude of thread interleavings one must test
\cite{nidhugg,optimal-dpor,coredet,sealing,quicksand,landslide,this-thesis,concuerror,
%cspec, % 2nd sentence in cspec >_>
bpor,
%parrot, % 2nd sentence in parrot
peregrine,
%bugtm, % 1st and 2nd sentence
satcheck,
racerx,
%ifrit, % 1st para, approx
datacollider,
fasttrack,
mcr-tsopso,
mcr,
macemc,
contessa,
%lvish, % "it must become easier"
dthreads,
taxdc,
%learning-from-mistakes, % sentence 3ish
avio,
symbiosis,
chess-icb,
chess,
recordreplaydrs,
kendo, % first WHOLE PARAGRAPH geez dude
%hybriddatarace, % 3rd sentence
parallel-dpor,
%eraser, % 3rd sentence 2nd paragraph
%racefuzzer, % 2nd sentence
%tsan, % 2nd sentence
sully-thesis,
%compcerttso, % 2nd sentence
modist,
maple,
simtester}.
%compared to sequential programming.
%
%Indeed, just as conditional statements and loops introduce additional possible program behaviours
%depending on different input,
%concurrent execution introduces additional possible behaviours even under a fixed input.
%Unlike input nondeterminism,
%concurrency nondeterminism typically manifests at the whims of chance,
%making bug-finding and verification especially challenging.
%
%What is a concurrency bug, and
How could a single programming technique be so troublesome
as to inspire multiple decades of verification research?
Allow me to begin this thesis by answering this question
in layperson's terms,
using two parables.

My favourite way to explain concurrency to a non-technical audience
%(e.g., family members)
is by analogy with unlocking car doors.
In many %makes and
models of cars, the door handle and locking mechanism are intertwined
in such a way that the door cannot unlock while its handle is being pulled~\cite{car-doors}.
Most car users, then,
%are familiar with
have suffered the experience of
trying to open the passenger door
at exactly the same moment as the driver turns the key in the lock:
%not only does the door not open, but
it fails to unlock, and
they must beseech the driver to turn the key again!
If they had pulled the handle any later, it would have opened as normal;
any earlier, and it would have been as though not having the key at all,
but would at least have unlocked properly at the first turn of the key later.
This violation of expected behaviour -- that the key need not be turned twice --
is a {\em bug},
and its dependence on the two travelers' actions {\em interleaving}
%with the driver's actions
in a particular way
makes it a {\em concurrency bug}.
This example involves only three events (pulling the handle, releasing the handle, and turning the key),
but what of a system with hundreds or thousands of times more?
Unfortunately, the number of possible interleavings grows {\em exponentially} with the number of events.

%Exponential explosion is easily understood through the tale of Sissa Ibn D\`{a}hir
The Islamic historian Ibn Khallik\={a}n told perhaps the first cautionary tale of exponential explosion
\cite{ibn-khallikan,khallikān1868ibn}.
Sissa Ibn D\`{a}hir,
credited with the invention of chaturanga (the precursor to modern chess),
was invited by King Shihr\^{a}m to request any reward he desired.
Sissa requested that a grain of wheat be placed on the first square of a chessboard,
two in the second, and so on, doubling the number of grains in each previous square until all 64 be filled.
The king at first laughed, thinking it a pittance,
realizing only later that to fulfill this would require
more than all the wheat in all the cities on Earth.
Modern programmers will recognize this sum as $2^{64}-1$,
the largest value representable by an unsigned 64-bit integer,
or approximately 18 billion billion.
The number of interleavings in a concurrent system grows at a similarly intractable rate
(although not necessarily exactly doubling each time),
with the number of chessboard squares filled with grain corresponding to the number of events in a program's execution.
A concurrency bug, then, is a single poisoned grain of wheat,
which must be rooted out from the stock before feeding the townspeople.
And Sissa himself is the programmer,
whose software is invariably large enough to be far beyond the reach of any exhaustive verification strategy.
}

\section{Motivation}

To take advantage of multiple cores for performance, programmers must write software to execute {\em concurrently} --
using multiple {\em threads} which execute multiple parts of a program's logic simultaneously.
However, when threads access the same shared data, they may interleave in unexpected ways which change the outcome of their execution.
When an unexpected interleaving produces undesirable program behaviour,
for example, by corrupting shared data structures,
we call it a {\em concurrency bug}.
Concurrency bugs are notoriously hard for programmers to find and debug
because the specific thread interleaving required to trigger them arises at random during normal execution,
and often with very low probability.
%Concurrency bugs are notoriously hard to find and reproduce because they only appear in specific thread interleavings, which arise at random during normal program execution.
% TODO(LAYMAN): give example of trying to open car door at same time as friend turns key to unlock it.

Most commonly, a programmer searches for concurrency bugs in her code by running it many times (in parallel, in serial, or both),
hoping that eventually, it will run according to the particular interleaving required to expose a hypothetical bug.
This technique, known as {\em stress testing}, is unreliable,
providing no guarantee of finding the failing interleaving in any finite amount of time.
It also provides no assurance of correctness:
when finished, there is no way of knowing how many distinct thread interleavings were actually tested.
Nevertheless, stress testing remains popular because of how easily a programmer can use it:
she simply wraps her program in a loop, sets it to run overnight, and kills it if her patience runs out before it finds a bug.

{\em Stateless model checking} \cite{verisoft} is an alternative way to test for concurrency bugs,
or to verify their absence,
which provides more reliable coverage, progress, and verification than stress testing.
A stateless model checker tests a program by forcing it to execute a new unique thread interleaving on each iteration of the test,
capturing and controlling the randomness in a finite state space of all possible interleavings.

Unfortunately, the size of these state spaces is exponentially proportional to the size of the tested program.
% TODO(LAYMAN): explain exponential explosion by relating the parable of grains of rice on a chessboard.
For even moderately-sized programs, there may be more possible ways to interleave every thread's every instruction
than particles in the universe.
Accordingly, a programmer who wants her test to make reasonable progress through the state space must choose a subset of ways that her threads could interleave,
focusing on fully testing that subset, while ignoring other possibilities she doesn't \revisionminor{think she cares} about.
However, it is difficult to choose a subset of thread interleavings that will produce a meaningful, yet feasible test.
Until computers can automatically navigate this trade-off in some intelligent way,
programmers will continue to fall back to the random approach of stress testing.

Another problem stateless model checking suffers is that certain types of programs cannot be tested without the programmer putting forth some manual instrumentation effort.
For example, operating system kernels implement their own sources of concurrency and their own synchronization primitives,
so \revisionminor{a checker must} be told how to identify and control the execution of each thread.
Some expert concurrency research wizards may be willing to add manual annotations to their code,
but required manual effort is a serious downside for anyone with a looming deadline,
and especially so for students who are still learning basic concurrency principles.
%We should not expect programmers to add effortful manual annotations to their code,
%or they will abandon our fancy technique to instead simply run stress tests until their deadline tomorrow evening.

\section{Contributions}

This thesis will solve both problems discussed above.
My thesis statement is as follows:

\begin{center}
\begin{tabular}{p{0.8\textwidth}}
	\revisionminor{{\em Combining %both
	theoretically-founded automatic reduction techniques
	and user-informed heuristic ones,
	stateless model checking can sufficiently mitigate exponential explosion
	to be a practical testing technique for inexperienced users and real-world programs alike.}}
\end{tabular}
\end{center}

\vspace{0.75em}

I have built Landslide \cite{landslide}, a stateless model checker for thread libraries and kernels,
and I have developed some techniques for automatically choosing the best thread interleavings to test
and for automatically instrumenting operating system kernels in an educational setting.
This thesis will comprise three major contributions:

\begin{enumerate}
	\item {\bf Meaningful state spaces (\Cref{chap:quicksand}).}
		I will present {\em Iterative Deepening}, a new algorithm for navigating the trade-off in how many preemption points to test at once.
		Iterative Deepening incorporates state space estimation \cite{estimation} to decide on-the-fly whether each state space is worth pursuing, and uses data race analysis \cite{tsan} to find new preemption point candidates based on a program's dynamic behaviour.
		This section will include a large evaluation of the technique, comparing its performance to three \revisionminor{prior-work} approaches across 600+ unique tests.
		I will show that Iterative Deepening of preemption points outperforms prior work in terms \revisionminor{of both} finding bugs quickly and completely verifying correctness when no bug exists.
	\item {\bf Educational use (\Cref{chap:education}).}
		For the past five semesters, I have offered a fully-automated version of Landslide to students in 15-410, CMU's undergraduate Operating System Design and Implementation class \cite{kspec,thrlib}, for use as a debugging aid during the thread library project.
		Recently I have also extended Landslide to handle Pintos kernel projects from other universities \cite{pintos}.
		In the two most recent semesters, I collaborated with Operating Systems course staff at two such schools, the University of Chicago and University of California at Berkeley,
		to provide debugging feedback to their students.

		At all three universities I then collected statistics on the numbers and types of bugs found,
		and surveyed students to understand the human experience.
		This section will present the study's results
		to evaluate the suitability of stateless model checking in an educational setting.
	\item {\bf Hardware Transactional Memory (HTM) (\Cref{chap:tm}).}
		HTM is a relatively new concurrent programming technique \cite{htm-haswell,htm-experience}
		which is not yet addressed by modern model checkers.
		I have extended Landslide's concurrency model to support
		HTM's execution semantics and several advanced features, % abort codes p much
		%both hardware (HTM) and software (STM) variants of TM,
		and tested several ``real-world'' HTM programs and benchmarks.
		This section will discuss the theoretical techniques I used to model the new form of concurrency,
		present associated correctness proofs of my approach,
		and show the verification results.
\end{enumerate}

\section{Organization}

The rest of this dissertation is organized as follows.

\begin{itemize}
	\item {\bf Background:} \Cref{chap:background} will present the requisite background material on concurrent programming, stateless model checking, and the various types of programs targeted by Landslide.
	\item {\bf Landslide:} \Cref{chap:landslide} explains the design and implementation of Landslide
		%the core of my stateless model checking framework,
		and all the special features it's been equipped with over the years.
		It is the foundation upon which all three contributions above build.
	\item {\bf Quicksand:} \Cref{chap:quicksand} presents the Iterative Deepening framework which more intelligently chooses which state spaces to test, corresponding to contribution 1 above.
	\item {\bf Education:} \Cref{chap:education} discusses my evaluation of Landslide
		in CMU's 15-410 class environment using the Pebbles kernel,
		and in the University of Chicago's and Berkeley's OS class environments using the Pintos kernel,
		corresponding to contribution 2 above.
	%\item {\bf Pebbles:} \Cref{chap:410} discusses my evaluation of Landslide in CMU's 15-410 class environment using the Pebbles kernel, corresponding to contribution 2 above.
	%\item {\bf Pintos:} \Cref{chap:410} discusses my evaluation of Landslide in the University of Chicago's and Berkeley's OS class environments using the Pintos kernel, corresponding to contribution 2 above.
	\item {\bf Transactional Memory:} \Cref{chap:tm} presents my extension of Landslide's concurrency model to handle transactional concurrency and the evaluation thereof, corresponding to contribution 3 above.
	\item {\bf Related Work:} \Cref{chap:relatedwork} honors my neighbours and ancestors in research spirit.
	\item {\bf Future Work:} \Cref{chap:warpzone} lists stateless model checking's remaining shortcomings and how new research might address them.
	\item {\bf Conclusion:} \Cref{chap:conclusion} provides some thoughts on the future of the field.
\end{itemize}

\section{Notes on reading this thesis}

\subsubsection{Experience}

I have tried to make this document accessible to readers of all programming experience levels,
although some of the research being theoretical and several layers of abstractions deep,
I cannot promise all easy reading.
\Cref{chap:background}, Background, provides what I hope are friendly concrete examples
to help the reader feel comfortable with each level of intuition that upcoming algorithms will build upon.
These should suffice for the experiments and overall contributions,
if not necessarily the details of each algorithm or soundness proof.
In particular, \Cref{chap:education}, Education,
may be approached with no knowledge of concurrency or model checking,
taking it merely as a study of a magic new debugging tool in the classroom setting.
The more ambitious reader may proceed to the Landslide chapter's algorithm walkthroughs~(\cref{sec:landslide-algs}),
which should equip them to understand every detail herein.
Readers who are here only to skim and skip around should at least be aware of the glossary~(\cref{sec:glossary})
to help clarify any terminology confusion.

\subsubsection{Vision}

Color will be used in figures and graphs to add visual contrast and make the data easier to navigate at a glance,
but only in redundant ways also signaled by symbols.
I have made some effort to choose palettes friendly to color-blindness;
should the reader find contrasts too low anyway,
whether being color-blind or reading a physical copy printed in greyscale,
they may be assured all important distinctions still render in monochrome.
For example, ovals and rectangles typically depict different threads,
and $\dagger$ distinguishes state space estimates from completed verifications.
\revision{Should any vision-impaired reader require fully-textual figure descriptions,
I would be happy to provide them upon request.}

\subsubsection{Language}

Pronoun use will vary between more specific and more ambiguous to convey additional nuance.
The singular ``I'' is associated with my own research contributions,
while the royal ``we'' should be taken to include the reader,
such as when surveying background material or related work,
to which the author and reader share more similar relationships.
The impersonal ``the programmer'' will be referred to as she/her
to challenge readers' unconscious bias about gender in computer science,
as well as to highlight her role as the intended user, separate from the underlying research.
%to distinguish her from the thesis audience concerned with the underlying research,
Gender-neutral pronouns will be used to refer to
individual students who participated in the user studies,
%will be given the more inclusive they/them.
%Gender-neutral pronouns will also be used
as well as
on the author themself.

\subsubsection{Code}

\newcommand\landslidelicense{BSD\xspace}

This document is, in a way, only half the work of the thesis, the other half being Landslide's implementation.
While some readers may prefer to be taught in prose and/or mathematical notation how an algorithm works,
others may find that disorienting and wish to see things in a way a compiler would understand.
\revisionminor{
The Landslide codebase is open-source under the \landslidelicense license,
available at \url{https://github.com/bblum/landslide}.
\Cref{chap:landslide} provides more detail and}
serves as a guide to browsing the repository.
Later chapters will often make parenthetical references to specific files and/or functions
therein which implement a feature under discussion.

% TODO:
% Finally, this thesis also contains a small puzzle-hunt to reward dedicated readers who make it all the way through.

\newpage
\thispagestyle{empty}
\begin{center}
\begin{tabular}{c}
\vspace{12em} \\
\includegraphics[width=0.55\textwidth]{how-i-learned.png}
\end{tabular}
\end{center}

\chapter{Background}
\label{chap:background}

This chapter will introduce the necessary background material on concurrency, stateless model checking, data-race analysis, and the relevant undergraduate operating systems classes.

\section{Concurrency Bugs}

\subsection{The Basics}

Modern software often turns to multithreading to improve performance.
In a multithreaded program, multiple execution units (or {\em threads}) execute the same or different sections of code simultaneously.

{\bf Simultaneity.}
This simultaneity of threads is achieved either by executing each one on a separate CPU, or by interleaving them nondeterministically (as controlled by clock interrupts) on the same CPU.
Because clock interrupts can occur at any instruction\footnote{
	With some exceptions in kernel-level programming, which we discuss later.
},
we consider single-CPU multithreading to be simultaneous at the granularity of individual instructions.
%
Likewise, when multiple CPUs access the same memory,
hardware protocols generally ensure that the events of a single instruction are executed atomically from the perspective of all CPUs.
Although there are some exceptions --
unlocked memory-to-memory instructions,
unaligned writes \cite{unaligned-writes},
and weak memory consistency models \cite{memory-consistency-models} --
we model multicore concurrency the same way as above,
deferring these exceptions beyond the scope of this work.
We refer to an execution trace depicting the global sequence of events as a {\em thread interleaving} or {\em schedule}.

{\bf Shared state.}
% TODO cite more languages / type theories
When a programming language offers multithreaded parallelism but forbids access to any shared state between threads \cite{rust-language},
the simultaneity of threads is largely irrelevant to the program's behaviour.
However, ``thread-unsafe'' languages such as C, C++, Java, and so on remain popular,
in which threads may access global or heap-allocated variables and data structures with no enforced access discipline.
The behaviour of such programs is then subject to the manner in which these accesses interleave.

{\bf Identifying bugs.}
Even if a program's behaviour is nondeterministic, that does not necessarily mean it has a bug.
After all, many programs use random number generation to intentionally generate different outputs.
We say a {\em concurrency bug} occurs when one or more of a program's nondeterministic behaviours is both {\em unanticipated} and {\em undesired}.
Most often, a concurrency novice who programs with shared state will consider the possible interleavings where one thread's access sequence occurs entirely before the other's, but neglect to consider intermediate outcomes in which the threads' access sequences are interleaved.
%
Consider the program in Figure \ref{fig:concurrency-bug}: Any output between 2 and 2000 is possible\footnote{
	Exercise for the reader: Show why 2 is a possible output, but 1 is not!
},
but whether this constitutes a bug is a matter of perspective.
Was the program written to count to 2000, or was it written to compute a randomized distribution?
In this thesis, we make no attempt to reason about the ``intent'' of programs,
so we further restrict {\em concurrency bug} to denote a program behaviour which is mechanically identifiable,
according to subjective notion of what programs behaviours are always bad.
%
Bug conditions include assertion failures,
memory access errors (i.e., segmentation fault or bus error),
heap errors (i.e., use-after-free or overflow),
deadlocks,
and infinite loops (which must be identified heuristically \cite{entscheidungsproblem}).

\begin{figure}[t]
	\begin{tabular}{cc}
		\begin{tabular}{p{0.45\textwidth}p{0.5\textwidth}}
			{\footnotesize
			\begin{tabular}{l}
				\texttt{int x;} \\
				\texttt{void count() \{} \\
				\texttt{~~~~for (int i = 0; i < 1000;} \\
				\texttt{~~~~~~~~~i++)} \\
				\texttt{~~~~~~~~x++;} \\
				\texttt{\}} \\
				\texttt{void main() \{} \\
				\texttt{~~~~tid1 = thr\_create(count);} \\
				\texttt{~~~~tid2 = thr\_create(count);} \\
				\texttt{~~~~thr\_join(tid1);} \\
				\texttt{~~~~thr\_join(tid2);} \\
				\texttt{~~~~printf("\%d\textbackslash{}n", x);} \\
				\texttt{\}} \\
			\end{tabular}
			}
			&
			{\footnotesize \begin{tabular}{ll}
				{\bf \normalsize Thread 1} & {\bf \normalsize Thread 2} \\
				\hline
				\texttt{load tmp <- x;} & \\
				& \texttt{load tmp <- x;} \\
				& \texttt{add tmp <- 1;} \\
				& \texttt{store x <- tmp;} \\
				\texttt{add tmp <- 1;} & \\
				\texttt{store x <- tmp;} & \\
			\end{tabular}
			}
			\\
			\\
			(a) Source listing for a multithreaded program which might count to 2000.
			&
			(b) Example interleaving of the compiled assembly for (a),
			in which 2 concurrent iterations of the loop yield 1 net increment of {\tt x}.
		\end{tabular}
	\end{tabular}
	\caption{Example concurrent program in which simultaneous accesses to shared state may interleave to produce unexpected results.}
	\label{fig:concurrency-bug}
\end{figure}

\subsection{Concurrency Primitives}

% TODO

%%%%%%%%%%%%%%%%%%%%%%%%%%%%%%%%%%%%%%%%%%%%%%%%%%%%%%%%%%%%%%%%%%%%%%%%%%%%%%%%

\section{Stateless Model Checking}

Model checking \cite{verisoft} is a testing technique for systematically exploring the possible thread interleavings of a concurrent program.
A model checker executes the program repeatedly, each time according to a new thread interleaving, until the state space (or the CPU budget) is exhausted.
During each execution, it forces threads to execute serially, thereby confining the program's nondeterminism to controlled thread switches.
Some model checkers explicitly store the set of visited program states as a means of identifying equivalent interleavings \cite{spin}.
This approach is called {\em stateful} model checking.
In this thesis, I focus on {\em stateless} model checking,
which instead analyzes the sequence of execution events to avoid a prohibitive memory footprint.
Henceforth I will abbreviate ``stateless model checking'' simply as ``model checking'' for brevity.

{\bf Static versus dynamic analysis.}
Model checking is a {\em dynamic} program analysis, meaning that it observes the operations and accesses performed by the program as its code is executed.
In contrast, {\em static} program analyses check certain properties at the source code level.
Static analyses are ideal for ensuring certain standards of code quality, which often correlates with correctness,
but cannot decide for certain whether a given program will fail during execution without actually running the code \cite{incompleteness}.
Static analyses face the challenge of {\em false alarms} (or {\em false positives}):
code patterns which look suspicious but are actually correct.
A debugging tool which reports too many false alarms will dissuade developers from using it \cite{racerx}.
Dynamic analysis, our approach, identifies program behaviours that are definitely wrong,
so each bug report is accompanied by concrete evidence of the violation.
Assertions, segfaults, use-after-free of heap memory, and deadlock are examples of such failures we check for,
although a checker may also include arbitrary program-specific predicates.

{\bf Preemption points.}
During execution, a model checker identifies a subset of the program's operations as ``interesting'', i.e.,
points at which interrupting the current thread to run a different one is likely to produce different behaviour.
These points, called {\em preemption points}, may be identified by any combination of human intuition and machine analysis.
Typical preemption points include the boundaries of synchronization APIs (e.g., {\tt mutex\_lock}) or accesses to shared variables.
Considering that at each preemption point multiple threads exist as options to run next, the set of possible ways to execute the program can be viewed as a tree.
% TODO: fig tree
I show an example execution tree in Figure~\ref{fig:tree}.
The number of preemption points in each execution define the depth of this tree, and the number of threads available to run define the branching factor.
Hence, in a program with $n$ preemption points and $k$ threads available to run at each, the state space size is $n^k$.
Addressing the scaling problem in this exponential relation is the central research problem for all model checkers.

{\bf Reduction techniques.}
Dynamic Partial Order Reduction \cite{dpor} (henceforth, DPOR) is the premier algorithm for mitigating the exponential explosion that arises as program size increases.
It identifies equivalent execution sequences according to Mazurkiewicz trace theory \cite{mazurkiewicz},
and tests at least one execution from each equivalence class.
Intuitively, if two thread transitions between preemption points do not conflict on any shared resource access, reordering them produces an equivalent interleaving, i.e., the same program behaviour.
Figure~\ref{fig:dpor} shows a part of an execution tree in which the operations by threads 1 and 2 are independent, and hence only one of the subtrees need be tested.
Over the years, researchers have developed many enhancements to DPOR, such as Optimal DPOR \cite{optimal-dpor}, parallelizable DPOR \cite{parallel-dpor}, SAT-directed model checking \cite{satcheck}, Maximal Causality Reduction \cite{mcr}, and DPOR for relaxed memory architectures \cite{tsopso}.

\begin{figure}[t]
	\begin{center}
	\includegraphics[width=0.4\textwidth]{dpor.pdf}
	\end{center}
	\caption{DPOR identifies independent transitions by different threads which can commute without affecting program behaviour. Here, if the transitions marked 1 and 2 have no shared memory conflicts, the states marked with the red arrow are guaranteed identical. Hence, only one of the subtrees need be explored.}
	\label{fig:dpor}
\end{figure}

{\bf Search heuristics.}
However, even though DPOR can prune an exponential number of redundant interleavings, the state space size is still exponential in the number of {\em dependent} (conflicting) interleavings.
Developers will always want to test larger and larger programs, so no matter the quality of our reduction algorithm,
we must accept that some tests will be too large to be fully tested in a reasonable time.
Hence, recent model checking research has turned to heuristic techniques for achieving further reduction,
optimizing the search to try to uncover bugs faster (should they exist)
at the expense of possibly missing other bugs,
or missing the chance to complete a full verification.
Iterative Context Bounding \cite{chess-icb} is a popular such technique which heuristically reorders the search to prioritize interleavings with fewer preemptions first.
This heuristic is based on the insight that most bugs require few preemptions to uncover, so interleavings with a number of preemptions that exceeds a certain bound will be de-prioritized, only tested until after all the fewer-preemption interleavings are completed.
Preemption sealing \cite{sealing} is another heuristic strategy which restricts the scope of the search by limiting the model checker to use only preemption points arising from certain functions in the source code.
This allows developers to vastly reduce state space size by identifying which program modules are already trusted,
although it requires some human intuition to correctly mark those boundaries.

%%%%%%%%%%%%%%%%%%%%%%%%%%%%%%%%%%%%%%%%%%%%%%%%%%%%%%%%%%%%%%%%%%%%%%%%%%%%%%%%

\section{Data Race Analsyis}

\begin{figure}[t]
        \small
	\begin{center}
\begin{tabular}{c}
\begin{tabular}{rll}
        & \multicolumn{2}{c}{\texttt{int x = 0; bool y = false; mutex\_t mx;}} \\
        \\
        & {\bf Thread 1} & {\bf Thread 2} \\
        1 & \texttt{\hilight{brickred}{x++;}~// A1} & \\
        2 & \texttt{mutex\_lock(\&mx);} & \\
        3 & \texttt{mutex\_unlock(\&mx);} & \\
        4 & & \texttt{mutex\_lock(\&mx);} \\
        5 & & \texttt{mutex\_unlock(\&mx);} \\
        6 & & \texttt{\hilight{brickred}{x++;}~// A2} \\
\end{tabular}
\\
\\
	{\normalsize (a) True potential data race.}
\\
\\
\begin{tabular}{rll}
        %& \multicolumn{2}{c}{\texttt{int x = 0; bool y = false; mutex\_t mx;}} \\
        & {\bf Thread 1} & {\bf Thread 2} \\
        1 & \texttt{\hilight{brickred}{x++;}~// B1} & \\
        2 & \texttt{mutex\_lock(\&mx);} & \\
        3 & \texttt{y = true;} & \\
        4 & \texttt{mutex\_unlock(\&mx);} & \\
        5 & & \texttt{mutex\_lock(\&mx);} \\
        6 & & \texttt{bool tmp = y;} \\
        7 & & \texttt{mutex\_unlock(\&mx);} \\
        8 & & \texttt{if (tmp) \hilight{brickred}{x++;}~// B2} \\
        %8 & & \texttt{if (tmp)} \\
        %9 & & \texttt{~~~~\hilight{brickred}{x++;}~// B2} \\
\end{tabular}
\\
\\
{\normalsize (b) No data race in any interleaving.}
\end{tabular}
	\end{center}
\caption{{Data-race analyses may be prone to either {\em false negatives} or {\em false positives}.
Applying Happens-Before to program (a) will miss the potential race possible between A1/A2 in an alternate interleaving,
while using Limited Happens-Before on (b) will produce a false alarm on B1/B2.}}
\label{fig:hb-example}
\end{figure}

Data race analysis \cite{eraser} identifies pairs of unsynchronized memory accesses between threads.
Two instructions are said to race if:
\begin{enumerate}
	\item they both access the same memory address,
	\item at least one is a write,
	\item the threads do not hold the same lock,
	\item and no synchronization enforces an order on the thread transitions (the {\em Happens-Before} relation, described below).
\end{enumerate}
%In Figure~\ref{fig:example}, lines 3 and 5 each race with 2 and 6, and line 6 races with 8.
In Figure~\ref{fig:hb-example}, the pairs of lines marked with comments (A1 and A2, B1 and B2) race.

A data race analysis may be either {\em static} (inspecting source code) \cite{racerx} or {\em dynamic} (tracking individual accesses arising at run-time) \cite{tsan}.
This paper focuses exclusively on dynamic analysis,
so although our example refers to numbered source lines for ease of explanation,
in practice we are actually classifying the individual memory access events corresponding to those lines during execution.
Actually, each {\tt x++} statement likely compiles to two separate load and store instructions, so each of those two instructions from each of the two marked source lines pairwise will race (except for the two loads, which are both reads).

{\bf Variants of Happens-Before.}
        Most prior work focuses on {\em Happens-Before} \cite{lamport-clocks} as the order relation between accesses.
\cite{predictive-dr} and \cite{hybriddatarace} identify a problem with this approach:
it cannot identify access pairs separated by an unrelated lock operation which could race in an alternate interleaving,
as shown in the example program in Figure~\ref{fig:hb-example}(a).
We call such unreported access pairs {\em false negatives}.
\cite{hybriddatarace} introduces the {\em Limited Happens-Before} relation,
which will report such potential races
by considering only blocking operations like {\tt cond\_wait} to enforce the order.
However, consider the similar program in Figure~\ref{fig:hb-example}(b),
in which the access pair ceases to exist in the alternate interleaving.
Limited Happens-Before will report all potential races, avoiding false negatives \cite{tsan},
but at the cost of necessarily reporting some such {\em false positives}.

%Finally, the {\em Causally-Precedes} relation \cite{predictive-dr} %strikes a middle ground,
%extends Happens-Before to additionally report a subset of potential races while soundly avoiding false positives.
%It tracks conflicting accesses in intervening
%critical sections to determine whether lock events are unrelated to a potential race.
%Causally-Precedes will identify the potential race in Figure~\ref{fig:hb-example}(a), as the two critical sections do not conflict,
%although it can still miss true potential races in other cases.

In Chapter~\ref{chap:quicksand}, we use the Limited Happens-Before relation for our analysis.
The justification for this is that, while stand-alone data-race analyses must avoid inundating the user with false alarms \cite{racerx},
my work incorporates data-race analysis in an internal feedback loop, and reports only directly observed failures to the user.
Hence, I accept some overhead from false positives for the sake of more thorough testing.

%%%%%%%%%%%%%%%%%%%%%%%%%%%%%%%%%%%%%%%%%%%%%%%%%%%%%%%%%%%%%%%%%%%%%%%%%%%%%%%%

\section{Education}

In this thesis I will tackle Pebbles and Pintos, two different system architectures used in educational operating systems courses.
This section describes the projects which students implement and which Landslide tests.

\subsection{Pebbles}

In the course of a semester, students work on five programming assignments;
the first two are individual, and the remaining three are the products of two-person teams.
I will focus on the third and fourth of these, the thread library and kernel,
called ``P2'' and ``P3'' respectively (the project numbers start at 0).
The other three (a stack-crawling backtrace utility, a bare-metal game with device drivers, and a small extension to the P3 kernel) are not of concern in this thesis.
Both P2 and P3 are built using the {\em Pebbles} system call specification, outlined in Table~\ref{tab:syscalls}

\begin{table}
        \center
        \begin{tabular}{|l|p{0.75\textwidth}|}
                \hline
                \bf System call name & \bf Summary \\
                \hline
                \multicolumn{2}{c}{\em Lifecycle management} \\
                \hline
                \texttt{fork} & Duplicates the invoking task, including all memory regions. \\
                \texttt{thread\_fork} & Creates a new thread in the current task.\\
                \texttt{exec} & Replaces the program currently running in the invoking task with a new one specified. \\
                \texttt{set\_status} & Records the exit status of the current task. \\
                \texttt{vanish} & Terminates execution of the calling thread. \\
                \texttt{wait} & Blocks execution until another task terminates, and collects its exit status.\\
                \texttt{task\_vanish}* & Causes all threads of a task to \texttt{vanish}. \\
                \hline
                \multicolumn{2}{c}{\em Thread management} \\
                \hline
                \texttt{gettid} & Returns the ID of the invoking thread. \\
                \texttt{yield} & Defers execution to a specified thread. \\
                \texttt{deschedule} & Blocks execution of the invoking thread. \\
                \texttt{make\_runnable} & Wakes up another \texttt{deschedule}d thread. \\
                \texttt{get\_ticks} & Gets the number of timer ticks since bootup. \\
                \texttt{sleep} & Blocks a thread for a given number of ticks. \\
                \texttt{swexn} & Registers a user-space function as a software exception handler.\\
                \hline
                \multicolumn{2}{c}{\em Memory management} \\
                \hline
                \texttt{new\_pages} & Allocates a specified region of memory. \\
                \texttt{remove\_pages} & Deallocates same. \\
                \hline
                \multicolumn{2}{c}{\em Console I/O} \\
                \hline
                \texttt{getchar}* & Reads one character from keyboard input. \\
                \texttt{readline} & Reads the next line from keyboard input. \\
                \texttt{print} & Prints a given memory buffer to the console. \\
                \texttt{set\_term\_color} & Sets the color for future console output. \\
                \texttt{set\_cursor\_pos} & Sets the console cursor location. \\
                \texttt{get\_cursor\_pos} & Retrieves the console cursor location. \\
                \hline
                \multicolumn{2}{c}{\em Miscellaneous} \\
                \hline
                \texttt{ls} & Loads a given buffer with the names of files stored in the RAM disk ``file system.'' \\
                \texttt{halt} & Ceases execution of the operating system. \\
                \texttt{misbehave}* & Selects among several thread-scheduling policies. \\
                \hline
        \end{tabular}
        \caption{The Pebbles specifcation defines 25 system calls. Students are not required to implement ones marked with an asterisk (*), though the reference kernel provides them. }
        \label{tab:syscalls}
\end{table}

{\bf P2.}
% TODO

{\bf P3.}
% TODO

\subsection{Pintos}
% TODO

\chapter{Landslide}
\inspirationalquote{Somewhere is the promise of an uncharted trail, with 700 branching limbs and 700 ways to fail.}
{ThouShaltNot, "Cardinal Directions"}

Landslide is a model checker implemented as a plug-in module for x86 full-system simulators.
The program to be tested runs in a simulated environment,
and Landslide uses its access to the simulator's internal state to inspect and manipulate the memory and thread scheduling of the program as it executes.
As of this thesis's writing, Landslide supports the use of two possible simulators:

\begin{itemize}
	\item {\bf Simics} \cite{simics}, a proprietary simulator licensed commercially by Wind River, used at CMU in 15-410 to run Pebbles thread libraries and kernels, and
	\item {\bf Bochs} \cite{bochs}, an open-source (LGPL) simulator used at the University of Chicago, Berkeley, and other schools to run Pintos kernels.
\end{itemize}

% TODO: update this linkerino (repository name)
% TODO: add some backup linkeroos
The Bochs port of Landslide is likewise open-source and available at \url{https://github.com/bblum/bochs}.
% TODO: refresh
The HEAD commit at the time of writing is 8984021.
The Simics port is available upon request.

This chapter will discuss Landslide's outer and inner workings for both aspiring users and aspiring developers in all gory detail. Strap in.

\section{User interface}

This section describes the features of Landslide the average student user should expect to interact with.
% TODO: make more specific section reference
Separate user guides also exist, described in Chapter~\ref{chap:410}.

\subsection{Setup}

Two setup scripts are provided, one for each supported kernel architecture: {\tt p2-setup.sh} and {\tt pintos-setup.sh}.
The user should supply the directory containing their project implementation.
The latter script also supports arguments specifying which of the Pintos projects to target.
For example:
\begin{itemize}
	\item {\tt ./p2-setup.sh /path/to/my/p2}
	\item {\tt ./pintos-setup.sh /path/to/my/threads} (2nd argument defaults to ``{\tt threads}'')
	\item {\tt ./pintos-setup.sh /path/to/my/userprog userprog}
\end{itemize}

These scripts accomplish the following setup tasks (among other trivialities):
\begin{itemize}
	\item Copy the user's code into {\tt pebsim/p2-basecode/} or {\tt pebsim/pintos/},
		which contain a pre-annotated Pebbles reference kernel binary or pre-annotated Pintos basecode, respectively.
	\item Build the code in its new location.
	\item Run the instrumentation script on the resulting binary to let Landslide know where all the important functions are
		(see \sect{\ref{sec:landslide-glue}}).
\end{itemize}

\subsection{Running Landslide through Quicksand}
\label{sec:landslide-quicksand-options}

The preferred method of invoking Landslide is through Quicksand, the Iterative Deepening wrapper program which has all of Chapter~\ref{chap:quicksand} to itself.
% TODO: add quicksand symlink
This is done via the {\tt ./landslide} script in the top-level directory, which:
\begin{itemize}
	\item Checks if the user needs to run {\tt *-setup.sh} again, in case their source code was more recently updated than the existing annotated build (a common mistake),
	\item Passes its arguments through to {\tt id/landslide-id}, the Quicksand binary,
		and
	\item (If during the student user study,) compresses the resulting log files,
		creates a snapshot tarball of them and the current version of the user's code,
		and sends it to me for nefarious research purposes.
\end{itemize}

\subsubsection{Command-line argments}

The following command line arguments are recommended for the common user.

\begin{itemize}
	\item {\tt -p PROGRAM}: the name of the test case to invoke
	\item {\tt -t TIME}: wall-clock time limit, in seconds; or suffixed with one of {\tt ydhms} for years, days, hours, minutes, or seconds respectively (default 1h)
	\item {\tt -c CPUS}: maximum number of Landslide instances to run in parallel (defaults to half the number of system CPUs)
	\item {\tt -i INTERVAL}: interval of time between printing progress reports (default 10s)
	\item {\tt -d TRACEDIR}: directory for resulting bug traces (default current directory)
	\item {\tt -v}: verbose mode (issues output for each executed interleaving by each instance of landslide, makes progress reports more detailed, etc)
	\item {\tt -l}: leave Landslide log files from completed state spaces even when no bug was found (deleted automatically by default)
	\item {\tt -h}: print help text and exit immediately
\end{itemize}

The following ``secret'' arguments also exist, primarily for my own use in running experiments or debugging.

\begin{itemize}
	\item {\tt -C}: enable ``control experiment'' mode, i.e., run only 1 instance of Landslide, with all (non-data-race) preemption points enabled in advance
		% TODO: put a section reference here
	\item {\tt -I}: enable Iterative Context Bounding (requires {\tt -C}, although future work may relax this restriction);
		this generally causes bugs to be found faster should they exist, but degrades completion time
		(\sect{\ref{sec:landslide-icb}})
	\item {\tt -0}: enable Preempt-Everywhere mode (\sect{\ref{sec:quicksand-eval}}, requires {\tt -C})
	\item {\tt -H}: use Limited Happens-Before for data-race analysis (\sect{\ref{sec:background-hb}})
		(default for Pebbles kernelspace mode)
	\item {\tt -V}: use vector-clock-based Pure Happens-Before for data-race analysis (\sect{\ref{sec:background-hb}})
		(default for P2 userspace and Pintos modes)
	\item {\tt -X}: support transactional memory (Chapter~\ref{chap:tm})
		% TODO: make a more specific section reference
	\item {\tt -A}: support multiple abort codes during transaction failure (Chapter~\ref{chap:tm});
		required for testing programs which behave differently under different abort circumstances,
		but impacts the state space size
	\item {\tt -P}: support Pintos architecture (enabled automatically when {\tt pintos-setup.sh} is run)
	\item {\tt -4}: support Pebbles architecture (enabled automatically when {\tt p2-setup.sh} is run)
		% TODO: put a section reference
	\item {\tt -e ETAFACTOR}: configure heuristic state space ETA deferring factor (described in detail in {\tt id/option.c})
	\item {\tt -E ETATHRESH}: configure heuristic threshold of state space progress for judging ETA stability (described in detail in {\tt id/option.c})
\end{itemize}

Quicksand will automatically generate configuration files and invoke Landslide according to the process described in the next section.

\subsection{Running Landslide directly}

Rather than letting Quicksand juggle multiple instances of Landslide,
the user may run a single instance directly, optionally configuring the preemption points by hand.
This is recommended only for enthusiastic users who are annotating their own kernel.

The script {\tt pebsim/landslide} invokes Landslide thus.
It should be run from within the {\tt pebsim/} directory.
When supplied no arguments, it reads configuration options from {\tt pebsim/config.landslide}
(a bash script expected to define certain variables as described in \sect{\ref{sec:landslide-glue}}).
The user may optionally specify a file containing additional config directives
% , such as custom preemption points,
as an argument.\footnote{
Quicksand actually supplies two such files as arguments: one ``static'' config file and one ''dynamic'' config file.
The former contains options which require recompiling Landslide (e.g., whether or not to use ICB is controlled by an {\tt \#ifdef} in Landslide's code),
while the latter contains options which Landslide interprets at runtime (e.g., which preemption points to use).
The static options do not change between Landslide instances in a single Quicksand run,
avoiding long Landslide start-up times.
}
Such supported options are as follows.

\subsubsection{Dynamic configuration options}
First, the following options may be changed without triggering a recompile of Landslide.
They are implemented as bash functions defined in {\tt pebsim/build.sh}.

\begin{itemize}
	\item {\tt within\_function FUNC} - adds {\tt FUNC} to a whitelist of functions required to appear in the current stack trace before identifying a preemption point (see \sect{\ref{sec:landslide-pps}})
	\item {\tt without\_function FUNC} - as above, but a blacklist instead of a whitelist
	\item {\tt within\_user\_function FUNC} - as two above but finds the function in the userspace test program rather than the kernel code.
	\item {\tt without\_user\_function FUNC} - difference to two above same as stated one above.
	\item {\tt data\_race ADDR TID LAST\_CALL CURRENT\_SYSCALL} - specifies a data-race preemption point.
		\begin{itemize}
			\item {\tt ADDR} shall be the code address (in hex) of the racing address,
			{\em before} the execution of which a preemption will be issued.
			\item {\tt TID} indicates a thread ID required to be running for this data race.
				To specify data race PPs across all threads at once, set {\tt FILTER\_DRS\_BY\_TID=0} (see next section).
			\item {\tt LAST\_CALL} indicates a code address required to be the site of the last {\tt call} instruction executed
				(similar to specifying a stack trace, but using a full stack trace here degrades performance too much),
				or 0 to not use this feature.
				From personal experience I found this option rather useless and recommend always supplying 0.
				% TODO: fix this section ref
				For further discussion see \sect{\ref{sec:quicksand-pps}}.
			\item {\tt CURRENT\_SYSCALL} indicates the system call number if a user-space data race comes from within a kernel system call which accesses user memory (Pebbles only).
				Usually 0 (i.e., not in kernel code) but {\tt deschedule}'s system call number is common as well.
		\end{itemize}
	\item {\tt input\_pipe FILENAME} - FIFO file used for receiving messages from Quicksand (e.g. to suspend or resume execution).
		Requires {\tt id\_magic} option to be set (next section below).
		The odds that a human user will find spiritual enlightenment through using this option by hand are infinitesimal.
	\item {\tt output\_pipe FILENAME} - as above but for sending messages.
\end{itemize}

\subsubsection{Static configuration options}

Next, configuration options which affect an {\tt \#ifdef} in Landslide and will trigger a recompile upon changing.
%These span a wide variety of features, sorted below in subsections by roughly how interesting I think they are.
Unless otherwise specified these are boolean flags (1 or 0) and the example value shown indicates the default used if unspecified.

{\bf Search algorithm options}
\begin{itemize}
	\item {\tt ICB=0} - enable Iterative Context Bounding (\sect{\ref{sec:landslide-icb}});
		corresponds to {\tt -I} in \sect{\ref{sec:landslide-quicksand-options}}.
	\item {\tt PREEMPT\_EVERYWHERE=0} - enable Preempt-Everywhere mode (\sect{\ref{sec:quicksand-eval}});
		corresponds to {\tt -0} in \sect{\ref{sec:landslide-quicksand-options}}.
	\item {\tt EXPLORE\_BACKWARDS=0} - configure whether, at each newly encountered preemption point,
		to allow the current thread to run first then later upon backtracking to preempt (0),
		or to issue preemptions first and then try continuing the current thread later (1).
		0 tends to produce shorter preemption traces while 1 tends to find bugs faster (\cite{landslide} \S{}8.7.1).
		Not compatible with ICB.
\end{itemize}

{\bf Memory analysis options}
\begin{itemize}
	\item {\tt PURE\_HAPPENS\_BEFORE=1} - select Pure Happens-Before (1) or Limited Happens-Before (2) (\sect{\ref{sec:background-hb}});
		corresponds to {\tt -V}/{\tt -H} in \sect{\ref{sec:landslide-quicksand-options}}.
	\item {\tt FILTER\_DRS\_BY\_TID=1} - configures whether to use the {\tt TID} parameter of {\tt data\_race} described above.
	\item {\tt FILTER\_DRS\_BY\_LAST\_CALL=0} - configures whether to use the {\tt LAST\_CALL} parameter of {\tt data\_race} described above.
	\item {\tt ALLOW\_LOCK\_HANDOFF=0} - configures lockset tracking to permit or disallow a lock taken by one thread to be released by another thread.%
		\footnote{If enabled, accesses performed by the second thread before unlocking will not be considered protected by that lock,
			as Landslide cannot infer what prior event abstractly represented the lock's ownership changing,
			leading to spurious data race reports.
			This could be solved in future work with a new annotation.}
	\item {\tt ALLOW\_REENTRANT\_MALLOC\_FREE=0} - allow two threads to be in {\tt malloc}, {\tt free}, or so on simultaneously without declaring it a bug.%
		\footnote{Used in Pintos, where those functions lock/unlock the heap mutex themselves rather than relying on a wrapper function to do so before invoking them.}
\end{itemize}

{\bf Interface options}
\begin{itemize}
	\item {\tt TEST\_CASE} - configure the name of the test program to run (mandatory; no default)
	\item {\tt VERBOSE=0} - enable more verbose output
	\item {\tt BREAK\_ON\_BUG=0} - configure whether to exit the simulator or drop into a debug prompt when a bug is found. Simics only and not compatible with Quicksand.
	\item {\tt DONT\_EXPLORE=0} - if enabled, Landslide will not perform stateless model checking but rather will execute the default thread interleaving then exit (useful for manual inspection of preemption points).
	\item {\tt PRINT\_DATA\_RACES=0} - as it says on the tin (for stand-alone use; will message them to Quicksand regardless).
	\item {\tt TABULAR\_TRACE=1} - configure whether to emit bug reports to the console (0) or to an HTML trace file (1)
\end{itemize}

\subsection{Kernel annotations}

The educational experiments in this thesis focus on projects which students implement on top of provided kernel basecode which Landslide already ``understands''.
Such understanding is conferred via the annotations described in this section.
For P2 and Pintos students I supply these annotations behind the scenes,
but a CMU 15-410 student who wishes to use Landslide on their kernel project shall need to brave forth hereupon.

\subsubsection{config.landslide annotations}

The following annotations are specified in {\tt pebsim/config.landslide} akin to the static configuration options described above.
These specify the names of kernel functions, global variables, default values, and so on
which are required to accurately track the kernel's scheduler state:
{\tt CONTEXT\_SWITCH},
{\tt EXEC},
{\tt FIRST\_TID},
{\tt IDLE\_TID},
{\tt INIT\_TID},
{\tt MEMSET},
{\tt PAGE\_FAULT\_WRAPPER},
{\tt PREEMPT\_ENABLE\_FLAG},
{\tt PREEMPT\_ENABLE\_VALUE},
{\tt READLINE},
{\tt SFREE},
{\tt SHELL\_TID},
{\tt SPURIOUS\_INTERRUPT\_WRAPPER},
{\tt THREAD\_KILLED\_ARG\_VAL},
{\tt THREAD\_KILLED\_FUNC},
{\tt TIMER\_WRAPPER},
{\tt VM\_USER\_COPY},
{\tt VM\_USER\_COPY\_TAIL},
{\tt YIELD}.

The following less self-explanatory secret options also exist:
\begin{itemize}
	\item {\tt PINTOS\_KERNEL=0} - configure Landslide for Pebbles (0) or Pintos (1) kernel architecture. Normally set automatically by the setup scripts.
	\item {\tt TESTING\_USERSPACE=1} - configure Landslide whether to test (i.e., focus preemption points, memory analysis, etc. on) the userspace or kernelspace code.
	\item {\tt CURRENT\_THREAD\_LIVES\_ON\_RQ=0} - Landslide infers the list of runnable threads from the {\tt on\_rq()} and {\tt off\_rq()} annotations (described below).
		Some kernels\footnote{most, actually} remove the current thread from their runqueue,
		such that the abstract set of all runnable threads is actually the runqueue plus the current thread rather than just the runqueue.
		Other kernels\footnote{the author's own student kernel from long ago}
		leave the current thread on the runqueue,
		removing it only when it's descheduling and should actually be considered blocked.
		Set this option to 0 to support the former kernel type or 1 to support the latter.
	\item {\tt PATHOS\_SYSCALL\_IRET\_DISTANCE} - indicate how much stack space is used by the reference kernel's system call wrappers (used for cross-kernel-to-userspace stack traces)
	\item {\tt PDE\_PTE\_POISON} - indicate a poison value used in the page tables to indicate absent VM mappings
	\item {\tt BUG\_ON\_THREADS\_WEDGED=1} - set to 0 to disable deadlock detection but instead let the kernel keep receiving system interrupts when all threads appear blocked.%
		\footnote{once used in the bad old days; now recommended for debugging use only}
	\item {\tt TIMER\_WRAPPER\_DISPATCH} - used to manually indicate a label before the end of the timer interrupt assembly wrapper, in case the {\tt iret} instruction couldn't be found automatically (see {\tt pebsim/definegen.sh}).
\end{itemize}

\subsubsection{In-kernel code annotations}

% TODO tell_landslide

\subsection{Bug reports}

\section{Architecture}

\subsection{Scheduler}
\subsection{Memory analysis}
\subsection{Machine state manipulation}
\subsection{State space traversal}
% time travel, explore, arbiter
\subsection{Bug-finding output}
% symtable stuff goes here
% stack tracing goes here
% foundabug goes here
\subsection{Pebbles-specific features}
\subsection{Pintos-specific features}
\subsection{Handy scripts}
\label{sec:landslide-glue}

\section{Algorithms of note}

\subsection{Kernel annotations}
\subsection{Bug identification}
\subsection{Data race analysis}
\subsection{Preemption point identification}
\label{sec:landslide-pps}
% TODO: talk about how w/in func and w/out func are processed
\subsection{State space estimation}
\subsection{Dynamic Partial Order Reduction}
\label{sec:landslide-dpor}
\subsection{Iterative Context Bounding}
\label{sec:landslide-icb}
\subsection{Heuristic loop, deadlock, and synchronization detection}

\chapter{Quicksand}
\label{chap:quicksand}
\inspirationalquote{
% original
% {\footnotesize 嫌な事も悲しい事もあったけど、守りたい物だって沢山この世界にはあったから…} \\
% trans. slightly paraphrased from crunchyroll subtitle; removed "..."s, added past tense for あった
There were awful, sad things in this world.
But there were a lot of things worth protecting, too.
}
{Kaname Madoka, Mahou Shoujo Madoka{\raisebox{0.1em}{$\scriptstyle \bigstar$}}Magica}
%\inspirationalquote{Always, somewhere, someone is fighting for you.
%As long as you remember her, you are not alone.}
%{Mahou Shoujo Madoka{\raisebox{0.1em}{$\scriptstyle \bigstar$}}Magica}

% TODO

\section{Preemption points} % ??? -- label needed from reference to LAST_CALL in landslide.tex
\label{sec:quicksand-pps}

\section{Evaluation}
\label{sec:quicksand-eval}

\chapter{Education}
\label{chap:education}

\inspirationalquote{Knowing the students might one day
%find a way to
fix their concurrency bugs...
it fills you with determination.}{Undertale (paraphrased)}

Concurrency is taught in as many different ways as there are
systems programming classes at universities which teach the subject.
Yet one thing they all have in common is presenting the concurrency bug
as some elusive menace,
against which humanity's best weapon is mere random stress testing.
This chapter will prove stateless model checking's mettle as a better alternative in the educational theatre.

While the previous chapter demonstrated Landslide's bug-finding power
compared to prior MC techniques in a controlled environment,
whether it offers pedagogical merit in the hands of students and/or TAs is a separate question.
And while my MS thesis~\cite{landslide} showed that students
could annotate P3 Pebbles kernels and thence use Landslide to debug them,
the annotations alone required 2 hours of effort on average per user,
meaning the only students who could benefit were the ones already succeeding enough to have such free time.
Since then, I have extended Landslide with a fully-automatic instrumentation process
for Pebbles thread libraries (P2s) (\sect{\ref{sec:education-pebbles-instrumentation}})
and Pintos kernels (\sect{\ref{sec:education-pintos-instrumentation}})
to improve its accessibility.

I have run several user studies in the Operating Systems classes
at Carnegie Mellon University (CMU), University of Chicago (U. Chicago), and Penn State University (PSU),
wherein students get to use Landslide to find and diagnose their own bugs during the semester.
At CMU, I analyzed logs and code snapshots taken as students used Landslide during P2
(\sect{\ref{sec:education-eval-bugfinding}}),
as well as the grades ultimately assigned after students who either did or didn't use Landslide submitted their projects
(\sect{\ref{sec:education-eval-grades}}).
At CMU and PSU, I surveyed students on their experience after submitting their Landslide-debugged P2s
(\sect{\ref{sec:education-eval-survey}}).
At U. Chicago, I collaborated with a TA to check submitted Pintos kernels,
then returned any resulting bug reports to students
(\sect{\ref{sec:education-eval-bugfinding}})
and likewise surveyed them on the quality of the diagnostic output
(\sect{\ref{sec:education-eval-survey}}).

%%%%%%%%%%%%%%%%%%%%%%%%%%%%%%%%%%%%%%%%%%%%%%%%%%%%%%%%%%%%%%%%%%%%%%%%%%%%%%%%

\section{Pebbles}

This section presents the user studies done in
%CMU's 15-410 and PSU's \psuos classes,
%in semesters Fall 2015 to Spring 2018 and in Spring 2018 alone, respectively,
%taught by David Eckhardt and Timothy Zhu, respectively.
CMU's 15-410 in semesters Fall 2015 to Spring 2018,
taught by David Eckhardt,
and in PSU's \psuos in Spring 2018,
taught by Timothy Zhu.
In both cases the instructors assisted to introduce me during the guest lecture
%(see below)
and to distribute the recruiting emails;
TAs were not involved.
The in-house user study has CMU IRB approval under study number STUDY2016\_00000425,
and the external user study under STUDY2017\_00000429.

% Possible experiment questions
% compare e.g. use after free bug reporce from OOPSLA data set, to P2 grade files, see who has thread exit uafs
% is landslide better at finding thread uafs than TAs
% same Q for other stuff.. (expect paraguay answer to be "no", explain why)

\subsection{Recruiting}
\label{sec:education-pebbles-recruiting}

Since the Spring 2015 semester I have given a guest lecture in 15-410
to recruit students to participate in the user study.
The 50-minute lecture is given 1 week into the 2.5-week-long P2 project,
approximately when the students should be getting child threads running in {\tt thr\_create()}
and experiencing concurrency bugs for the first time.
It introduces the research subject abstractly
using an example ``Paradise Lost'' bug from a previous lecture \cite{paradise-lost},
explains how Landslide works concretely,
shows a short demo of effortlessly using Landslide to find the example bug,
and provides the necessary IRB legalese about the risks and benefits of participation.
The most recent lecture slides are available on the course website at
\url{http://www.cs.cmu.edu/~410-s18/lectures/L14_Landslide.pdf},
and all semesters' editions at
\url{https://github.com/bblum/talks/tree/master/landslide-lecture}.

The PSU version of the lecture
%was given in Spring 2018, and
is available at
\url{http://www.contrib.andrew.cmu.edu/~bblum/psu-lecture.pdf}
as well as under the github link above.
Being a 70-minute lecture slot rather than 50, I extended the demo to
both find and (attempt to) verify a fix for two bugs:
one a simple data race and the other the more complicated Paradise Lost bug as above.
After finding each bug, I demonstrated using Landslide on a fixed version of the code
to show how it proves the test case correct by completing all state spaces,
or (in the case of Paradise Lost) suffers an exponentially-exploding state space.
Not that I scientifically measured it or anything,
but this extended demo seemed to help students more clearly understand Landslide's intended workflow,
at the cost of about 10-15 extra minutes of lecture time.

At both schools students then signed up using a Google form I emailed them,
which upon completion linked them to the Landslide user guide,
which is available at
\url{http://www.contrib.andrew.cmu.edu/~bblum/landslide-guide-p2.pdf}
(CMU version)
and
\url{http://www.contrib.andrew.cmu.edu/~bblum/landslide-guide-psu.pdf}
(PSU version)
and
\url{https://github.com/bblum/talks/tree/master/irb}
(both versions).

During the last week of P2 at CMU, I held several ``Landslide clinic'' sessions
(basically office hours, but given a different name to remind students
to limit themselves to questions a normal TA couldn't answer),
where students could receive in-person technical and/or moral support.
Collecting study information during these sessions was not included in the IRB protocol.
In the PSU study, I had returned to Pittsburgh shortly after giving the lecture,
so technical support for PSU students was limited to email correspondence.

\subsection{Automatic instrumentation}
\label{sec:education-pebbles-instrumentation}

As described in \sect{\ref{sec:landslide-setup}},
all setup from the user's point of view is handled through the {\tt p2-setup.sh} script%
\footnote{PSU's version is called {\tt psu-setup.sh};
in this section {\tt p2-setup.sh} refers to both unless otherwise noted.
}.
It, its helper scripts (\sect{\ref{sec:landslide-glue}}),
and the {\tt landslide} script itself contain several checks to prevent
studence
from accidentally misusing Landslide in ways that could produce mysterious crashes, false bug reports, and so on
(the need for each one, as the reader might imagine, discovered through bitter experience).
These include:
\begin{itemize}
	\item {\tt p2-setup.sh} checks if the directory argument correctly points at the top-level P2 basecode directory
		rather than any subdirectories such as {\tt user/libthread/}.
	\item {\tt check-need-p2-setup-again.sh} checks if any source files in the original P2 source directory
		(the argument supplied to {\tt p2-setup.sh}),
		in case the student hoped to fix some bug and verify their fix but forgot to re-run the setup script.
	\item {\tt landslide} checks the supplied test name matches one of the endorsed Landslide-friendly tests
		(students love trying to run Landslide with {\tt racer}, {\tt largetest},
		or even the string {\tt OPTIONS}).
	\item {\tt landslide} checks if any other instance of itself is simultaneously running in the same directory,
		and if so, refuses to do so and advises the student
		to {\tt git clone} the repository afresh for simultaneous use%
		\footnote{This is ironically implemented with a non-atomic lock file
		and should really be using {\tt flock} instead.
		}.
\end{itemize}
\vspace{1em}

\noindent Landslide also includes several P2-specific instrumentations and features to cope with various student irregularities:
\begin{itemize}
	\item Quicksand emits different combinations of {\tt within\_function}/{\tt without\_function} directives
		for Landslide depending on the name of the test.
		For example, for {\tt paradise\_\allowbreak{}lost} Landslide will not preempt in a function named {\tt critical\_section()},
		which the test case uses to protect an internal counter used to detect the bug;
		and it will not preempt in any of the {\tt thr\_*()} thread library API functions
		for tests intended to target just the concurrency primitives.
		In future work this could be improved as annotations to be placed inside the test case code itself.
	\item Landslide finds ad-hoc synchronization patterns,
		such as {\tt while (!flag) yield()} or {\tt while (xchg(...)) continue},
		which students often open-code rather than using the prescribed synchronization API,
		and treats them as synchronization points as described in \sect{\ref{sec:landslide-blocking}}.
	\item Landslide finds ``too suspicious'' spinwait-loops in mutex implementations
		which are neither yield- nor {\tt xchg}-loops (as described above),
		which would ordinarily be classified as infinite loop bugs,
		and reports them with a suggestive message ({\tt undesirable\_\allowbreak{}loop\_html()} in {\tt landslide.c})
		referring the student to the appropriate lecture material
		\cite{synchronization-2}.
	\item The {\tt landslide} wrapper script logs the time and command-line options of invocation
		and captures a snapshot of the student code and results of the test and saves them to AFS
		(CMU's network file system)
		after each run.
		% HURDLE_VIOLATION -- no, this is for p3 (context switcher) features only
\end{itemize}

\subsection{Test cases}
\label{sec:education-pebbles-tests}

Landslide ships with several ``approved'' test cases,
i.e., programs copied from, derived from, or at least vaguely resembling
the tests distributed with P2,
which I curated to produce concurrent behaviour suitable for stateless model checking.
Some tests are crafted to target specific bugs which,
from personal experience as a TA, are common in many student submissions;
others are crafted to exercise generally concurrency-heavy code paths and uncover any number of unforeseen problems.
Many use some of the features/annotations described in \sect{\ref{sec:landslide-testcases}}.

%\subsubsection{Test case list}

The following tests were released to CMU students:
\begin{itemize}
	\item {\tt broadcast\_test}:
		Tests the {\tt cond\_broadcast()} signalling path with a single waiter.
	\item {\tt mutex\_test}:
		Tests student mutexes under 2 threads with 2 iterations
		(the 2nd iteration serves to expose problems with {\tt mutex\_unlock()} as well as {\tt mutex\_lock()}).
		This test uses the {\tt TESTING\_MUTEXES}
		described in \sect{\ref{sec:landslide-staticconfig}}
		to enable data-race preemption points within the mutex implementation.
	\item {\tt paradise\_lost}:
		Written for the sake of the Landslide lecture demo
		(\sect{\ref{sec:education-pebbles-recruiting}}).
		Tests for the Paradise Lost bug by attempting to break mutual exclusion.
	\item {\tt paraguay}:
		Copied directly from the P2 test suite;
		tests for proper handling of seemingly ``spurious'' wakeups in {\tt cond\_wait()}.
		Written by Michael Sullivan.
	\item {\tt rwlock\_downgrade\_read\_test}:
		Copied directly from the P2 test suite;
		tests for mutually-exclusive and deadlock-free {\tt rwlock\_downgrade()}.
		Written by me (as a TA).
	\item {\tt thr\_exit\_join}:
		Copied directly from the P2 test suite;
		tests for a variety of problems between {\tt thr\_exit()} and {\tt thr\_join()},
		but especially for memory issues pertaining to stack deallocation.
\end{itemize}
\vspace{1em}

\noindent The following tests were released to PSU students, in addition to the ones above:
\begin{itemize}
	\item {\tt atomic\_compare\_swap}:
		Tests the {\tt cmpxchg} assembly function for being properly atomic.
		Uses the {\tt magic\_*} global variables described below, and invokes {\tt vanish()} directly,
		to avoid requiring the student to implement {\tt thr\_join()}/{\tt thr\_exit()}
		before being able to run this test.
	\item {\tt atomic\_exchange}:
		As above for {\tt xchg}.
	\item {\tt atomic\_fetch\_add}:
		As above for {\tt xadd}.
	\item {\tt atomic\_fetch\_sub}:
		As above for {\tt xadd}.
	\item {\tt broadcast\_two\_waiters}:
		As {\tt broadcast\_test}, but uses two waiting threads to ensure both get signalled.
\end{itemize}
\vspace{1em}

\noindent The tests can all be viewed at \url{https://github.com/bblum/landslide/tree/master/pebsim/p2-basecode/410user/progs}.

\subsection{Survey}
\label{sec:education-survey-pebbles}

Starting in Fall 2017, I sought to gauge the students' personal opinions on their experience with Landslide,
in addition to simply counting
from the automatic snapshots
how many bugs were found.
Shortly after the P2 submission deadline,
I asked participants to answer several survey questions, reproduced below.
%distributed via email as a Google Doc.

\begin{enumerate}
	\item How many bugs did Landslide help you find in your code? (Please indicate a number.)
	\item How many of the bugs you found with Landslide do you believe you fixed before submitting your project? (You may answer ``all'', ``none'', or a number.)
	\item How many of the bugs you found with Landslide did you verify you had fixed by running Landslide again to make sure the bug was gone? (You may answer ``all'', ``none'', or a number.)
	\item In addition to the bugs Landslide found, did it report anything that you believe was NOT a bug? For example, Landslide printed an execution trace that was actually impossible, or Landlside reported a bug about some behaviour that was actually allowed by the P2 specification. (If so, please describe.)
	\item I found Landslide's debugging output easy to understand.
		(Multiple choice from strongly disagree to strongly agree.)
	\item It's easier to diagnose the root cause of a bug with Landslide than with a stress test (e.g. juggle).
		(Multiple choice from strongly disagree to strongly agree; plus ``Not sure'' and ``Easier for some bugs but harder for others'')
	\item I felt the time I saved by having Landslide to help debug was worth the time it took me to learn how to use Landslide.
		(Multiple choice from strongly disagree to strongly agree.)
	\item I feel that by using Landslide I learned to understand concurrency better.
		(Multiple choice from strongly disagree to strongly agree.)
	\item Suppose after you submitted your project, % "P2" for cmu, "project" for pintos
		we gave you 100 CPU-hours on the cloud provider of your choice to test it. Then we extended the project deadline by a day for you to use the results to fix bugs and get partial credit. How would you divide that CPU time between the staff-provided stress tests and Landslide?
		(Multiple choice: 0/10/.../100 CPU-hours on Landslide, 100/90/.../0 CPU-hours on stress tests.)
	\item If I found out next semester that a friend of mine (or a student in my degree program) were taking OS, I would recommend that they should probably invest some time during the project % "P2" for cmu, as above
		to learn Landslide and try to find bugs with it.
		(Multiple choice from strongly disagree to strongly agree.)
	\item Regarding the previous question, why or why not?
\suspend{enumerate}
\vspace{1em}

\noindent The following questions were served only on the CMU version of the survey.
\resume{enumerate}
	\item Did you answer this survey together with your partner, or on your own while they were busy? (If you both have time for it, please try to submit one survey together.)
		(Multiple choice: together or alone)
	\item Your andrew ID
	\item Your partner's andrew ID (if any)
\suspend{enumerate}
\vspace{1em}

\noindent The following questions were served only on the PSU version of the survey.
\resume{enumerate}
	\item Any feedback on how Landslide's user interface could be improved / made easier to use or understand? (setup process, messages printed while running, or the execution trace / stack traces emitted after a bug is found?)
	\item Your PSU username
\end{enumerate}

%%%%%%%%%%%%%%%%%%%%%%%%%%%%%%%%%%%%%%%%%%%%%%%%%%%%%%%%%%%%%%%%%%%%%%%%%%%%%%%%

\section{Pintos}

This section presents the user study done in U. Chicago's \uchos class in the Fall 2017 semester,
taught by Haryadi Gunawi.
Kevin Zhao, the TA, assisted to run Landslide on student submissions
and to distribute recruiting materials and testing results.
The study has CMU IRB approval under study number STUDY2017\_00000429.

\subsection{Recruiting}

For this study students were recruited remotely via email.
After each of the {\em threads} and {\em userprog} project deadlines (\sect{\ref{sec:overview-pintos}}),
\uchos staff sent students an email inviting them to volunteer to receive Landslide's bug reports,
disclaiming that it did not represent part of the official grading process but could help improve their future submissions.

\subsection{Automatic instrumentation}
\label{sec:education-pintos-instrumentation}

As described in \sect{\ref{sec:landslide-setup}},
all setup from the user's point of view is handled through the {\tt pintos-setup.sh} script.
It
and its helper {\tt pebsim/pintos/import-pintos.sh}
perform most of the same sanity checks as listed in \sect{\ref{sec:education-pebbles-instrumentation}},
then applies the patch {\tt annotate-\allowbreak{}pintos.patch}
(plus several more hacks in the script itself)
to insert the {\tt tell\_landslide()} annotations (\sect{\ref{sec:tell-landslide}})
into the student's kernel code.
The following tricks were developed after trial-and-error on student submissions from the same semester,
and serve to make sure the annotations apply consistently
to (almost) all variations of commonly-submitted code.

\begin{itemize}
	\item Finds the declaration of {\tt ready\_list}, the scheduler runqueue declared by the basecode,
		and detects if the student has modified to be an array of lists rather than a single one.
		If so, defines the length of that array in a macro to be used by {\tt is\_runqueue()}
		(part of the patch described below).
		Either way defines a function {\tt get\_rq\_addr()} to return the address of the (first) list.
	\item Changes the basecode's definition of {\tt TIME\_SLICE} from 4 to 1 (units of timer ticks)
		so Landslide's timer injection will properly drive the context switcher.
	\item Inserts {\tt tell\_landslide\_forking()} into {\tt thread.c}
		(using {\tt sed} rather than the patch, described below,
		because it must go in a function which students have to implement,
		which is likely to disturb the context and make a patch fail).
	\item Adds the new {\tt priority-donate-multiple} test.
	\item Applies the {\tt annotate-pintos.patch} patch to the imported student implementation, which:
	\begin{itemize}
		\item Adds {\tt tell\_landslide\_thread\_on\_rq()}
			and {\tt tell\_landslide\_thread\_off\_rq()}
			annotations
			to {\tt list\_insert()} and {\tt list\_remove()} respectively
			(in {\tt lib/kernel/\allowbreak{}list.c}, which the students don't modify),
			which
			check whether the argument list
			is the scheduler runqueue
			using a helper function {\tt is\_runqueue},
			which in turn uses {\tt get\_rq\_addr()} and {\tt READY\_LIST\_LENGTH} described above.
		\item Modifies the existing {\tt priority-sema} and {\tt alarm-simultaneous} tests to be more Landslide-friendly.
		\item Inserts the {\tt tell\_landslide\_sched\_init\_done()},
			{\tt tell\_landslide\_vanishing()},
			and {\tt tell\_landslide\_thread\_switch()}
			annotations in the appropriate places
			(which the students generally do not modify).
	\end{itemize}
	\item Detects if the student has renamed the {\tt elem} field of the TCB struct,
		and if so renames its use in {\tt is\_runqueue()} (described above) correspondingly.
	\item Detects if the student has renamed the {\tt cur} (currently running thread) variable
		in the context switcher, and if so renames it back.
\end{itemize}

\subsection{Test cases}
\label{sec:education-pintos-tests}

Like the P2 tests, the Pintos test cases are either hand-picked from the provided unit tests,
with an eye for which will produce interesting concurrent behaviour,
or created using a TA's intuition for the most likely student bugs.
The following tests are approved to be Landslide-friendly:

\begin{itemize}
	\item {\tt priority-sema}:
		Modified to be Landslide-friendly from the basecode, for {\em threads}.
		Creates two child threads to wait on a semaphore and signals them.
		Replaces threads with different priorities
		(originally chosen to produce deterministic output which the test checked for)
		with threads of the same priority.
	\item {\tt alarm-simultaneous}:
		Modified to be Landslide-friendly from the basecode, for {\em threads}.
		Creates two child threads which each invoke {\tt timer\_sleep()} for a different amount of time.
		Number of (threads,iterations) reduced from (3,5) to (2,1).
	\item {\tt wait-simple}:
		Unmodified from the basecode's version, for {\em userprog}.
		Userspace process {\tt exec()}s a child process, which immediately exits, and {\tt wait()}s on it.
	\item {\tt wait-twice}:
		Unmodified from the basecode's version, for {\em userprog}.
		Slightly more complicated version of {\tt wait-simple},
		intended to expose failure-path bugs if the former finds no common-case ones.
	\item {\tt priority-donate-multiple}:
		Written by Kevin Zhao, TA at U. Chicago, for {\em threads}.
		Tests for a priority donation race during {\tt lock\_release()}
		in which a thread holding a lock can accidentally keep a contending thread's donated priority
		after finishing releasing it.
\end{itemize}
\vspace{1em}

The (unpatched versions of) the first four tests are available at
\url{https://github.com/Berkeley-CS162/group0/tree/master/pintos/src/tests}.
The fifth test is available at
\url{https://github.com/bblum/landslide/blob/master/pebsim/pintos/priority-donate-multiple.c}.

\subsection{Survey}
\label{sec:education-survey-pintos}

Similar to the survey for Pebbles projects
\sect{\ref{sec:education-survey-pebbles}},
I surveyed the Pintos user study participants for their opinions.
Because of the different nature of the user study, of course,
the questions here focus more on the debugging experience than on using Landslide directly.

\begin{enumerate}
	\item How many Landslide bug reports did you receive from course staff? (Please indicate a number.)
	\item Among those bug reports, how many were you able to diagnose and recognize the root cause in your code? (You may answer ``all'', ``none'', or a number.)
	\item Among those bug reports, how many described a behaviour that you believe was NOT a bug? For example, Landslide printed an execution trace that was actually impossible, or Landslide reported a bug about some behaviour that was actually allowed by the Pintos specification. (You may answer ``all'', ``none'', or a number.)
	\item About how much time did you spend interpreting Landslide's debugging output? (Please indicate a number of minutes, or a range if uncertain, e.g. ``30-60 minutes''.)
	\item I found Landslide's debugging output easy to understand.
		(Multiple choice from strongly disagree to strongly agree.)
	\item It's easier to diagnose the root cause of a bug with Landslide than with a stress test (for example {\tt exec-multiple}).
		(Multiple choice from strongly disagree to strongly agree; plus ``Not sure'' and ``Easier for some bugs but harder for others'')
	\item I feel that by interpreting Landslide's debugging output I learned to understand concurrency better.
		(Multiple choice from strongly disagree to strongly agree.)
	\item These kinds of concurrency bugs are important to fix, even though they don't count against my grade.
		(Multiple choice from strongly disagree to strongly agree.)
	\item Suppose after you submitted your pintos, we gave you 100 CPU-hours on the cloud provider of your choice to test it. Then we extended the project deadline by a day for you to use the results to fix bugs and get partial credit. How would you divide that CPU time between the staff-provided stress tests and Landslide?
		(Multiple choice: 0/10/.../100 CPU-hours on Landslide, 100/90/.../0 CPU-hours on stress tests.)
	\item If course staff were to allow students to resubmit updated code after reviewing Landslide bug reports to receive partial credit for each bug that had been fixed, it would be worth my time to try that (even if I could be spending that time working on the next project instead).
		(Multiple choice from strongly disagree to strongly agree.)
	\item If a friend of mine took OS next semester, I would recommend that they should sign up to receive Landslide bug reports during projects in the future.
		(Multiple choice from strongly disagree to strongly agree.)
	\item Regarding the previous question, why or why not?
	\item Your name.
	\item Your project partner's name (if applicable)
\end{enumerate}

%%%%%%%%%%%%%%%%%%%%%%%%%%%%%%%%%%%%%%%%%%%%%%%%%%%%%%%%%%%%%%%%%%%%%%%%%%%%%%%%

\section{Evaluation}

I pose the following evaluation questions.

\begin{itemize}
	\item How many bugs, and of what severities, does Landslide help students find and fix before submitting their code?
	\item Does Landslide use result in higher quality submissions, whether directly for P2, or for subsequent projects as well?
	\item Do students feel the experience is worthwhile, as compared to stress testing?
	\item How well does Landslide apply to operating systems projects outside of CMU?
\end{itemize}
\vspace{1em}

The data set is comprised of Landslide's automatically-generated usage snapshots
from CMU students from semesters Spring 2015 to Spring 2018 inclusive,
% and PSU \psuos from Spring 2018, % no snapshots from psu
CMU's official project grades from same,
CMU students' survey responses and submitted project code from Fall 2017 and Spring 2018
PSU students' survey responses and submitted project code from Spring 2018,
and U. Chicago students' survey responses from Fall 2017.
Table~\ref{tab:photo-of-ze-studence} shows the student participation rate across semesters
% TODO update for psu
for the thread library study%
\footnote{
Participation at CMU was determined by receiving any automatic usage snapshots;
participation at PSU was determined, for lack of snapshots, by sign-up form submissions, which might over-count slightly.
}.
% and PSU

\begin{table}[h]
	\begin{center}
		% ugh this looks so stupid
		% but better than latex's output looking stupid
		\begin{tabular}{c|crcrc|r}
			& \multicolumn{5}{c|}{\bf used Landslide?} & \\
			\bf semester & & \bf yes & & \bf no & & \bf total \\
			\hline
			CMU S'15	& &   8	& &  18	& &  26 \\
			CMU F'15	& &  22	& &   8	& &  30 \\
			CMU S'16	& &  18	& &  16	& &  34 \\
			CMU F'16	& &  12	& &   6	& &  18 \\
			CMU S'17	& &  19	& &  13	& &  32 \\
			CMU F'17	& &  14	& &   5	& &  19 \\
			CMU S'18	& &  10	& &   8	& &  18 \\
			% NB. CMU's participance are counted by how many ppl i got at least one usage tarball for,
			% while PSU's are by how many sign-up forms i got.
			% the latter may be overcounting, because a student may sign up and then never use it;
			% but i don't have tarballs for PSU so this is the best you're getting.
			% TODO
			% PSU s18	& &  38
			\hline
			\bf total	& & 103	& &  74	& & 177 \\
		\end{tabular}
	\end{center}
	\caption{Student participation rates across semesters,
		among students who submitted thread library implementations.}
	\label{tab:photo-of-ze-studence}
\end{table}

\subsection{Bug-finding}
\label{sec:education-eval-bugfinding}

Firstly, I sought to prove Landslide does as advertised:
finds concurrency bugs of severity, subtlety, and difficulty
consistent with the lessons an advanced operating systems class should hope to teach its students,
and provides diagnostic output that helps students understand and solve them.
% without necessarily spoiling the answer to the puzzle ... not scientifically evaluatable,
% but at least the trace format is de0u-approved.

At CMU, I configured Landslide to save a usage snapshot every time a student ran Landslide,
including command-line options issued, the current version of their project code,
a log of Landslide's output,
and the preemption traces for any bugs found.
These I analyzed by hand to determine how many distinct bugs were reported
(as multiple preemption traces may refer to the same bug; inspecting their code if necessary),
how many were deterministic
(i.e., Landslide reported them on the very first interleaving tested, with no need for artificial preemptions)
and how many were concurrency bugs,
whether any were false positives,
and whether the student was able to fix them thereafter
(determined if a subsequent snapshot showed them re-running the test and verifying the bug's absence;
an approximate measure at best, but far easier to implement than checking all the bugs by hand).

At U. Chicago,
course staff ran Landslide on student submissions behind-the-scenes
after each project deadline had passed,
and returned its preemption traces to the study volunteers.
I collaborated with the course staff
to confirm whether each trace represented a real bug and not a false positive
before distributing them to the students.
On account of the small sample size (4 volunteering groups),
we decided on this approach,
rather than to study how students would cope with false positives themselves,
to optimize for student happiness over scientific rigor.

As PSU's version of the study was planned on relatively short notice,
and without immediate access to a shared, yet confidential, academic file system (such as CMU's AFS),
I was unable to examine its students' objective usage data or bug reports for this section.
Results from PSU are presented instead in \sect{\ref{sec:education-eval-survey}}.

\subsubsection{CMU}
\label{sec:education-eval-bugs-cmu}

In the recruiting lectures for every semester after the first,
I included tallies of Landslide's bug-finding achievements to date
as an extra advertisement to tempt students.
These included each of the measures I counted by hand as described above,
as well as aggregate totals of how many groups participated,
how many groups found bugs,
how many found at least one concurrency bug,
and how many were able to fix and subsequently verify any or all thereof.
Among the concurrency bugs, a wide variety of types were found:
deadlocks, use-after-frees, segfaults, infinite loops, assertion failure, and unit test failure.
Table~\ref{tab:cmubugs} shows these statistics for each semester.
% may also be an overestimate if they only hacked around the bug, but that's less likely
Considering the method for deciding whether students fixed each bug, described above,
note that the tallies thereof may be undercounting
on account of the possibility some students may have skipped the verification step despite truly fixing them.

\begin{table}[t]
	% TODO
	\begin{center}
		\small
	\begin{tabular}{r||r|r|r|r|r|r|r||r}
			& \bf S'15 & \bf F'15 & \bf S'16 & \bf F'16 & \bf S'17 & \bf F'17 & \bf S'18 & \bf Total \\
			\hline
			%			S15	F15	S16	F16	S17	F17	S18
			\bf Participating groups& 8	& 22	& 18	& 12	& 19	& 14	& 10	& 103	\\
			% nb not counting FPs
			Groups w/any bugs found	& 5	& 21	& 12	& 7	& 14	& 10	& 6	& 75	\\
			Groups w/any bugs fixed	& 4	& 18	& 9	& 6	& 9	& 7	& 5	& 58	\\
			Groups w/all (>0) bugs fixed
						& 2	& 11	& 3	& 3	& 7	& 5	& 4	& 35	\\
			\hline
			% the 2 for f15 were pretty messy; i had to go back and get em by hand by grepping for NOTE
			% but this semester also had a LOT of seemingly true race NOTE mx lock inf loop,
			% so i counted only the det ones, of which there were 2 (both also in "fixed" category)
			% otherwise 1 race recursive mx and 6 race mutex loops and 1 race cvar loop were left as races.

			% this 1 for f17 was salvaged from results-raw; the grupo who had it it was their only bug
			% and it was an invl heap that i think came from multiple landslides
			% (and i hadn't listed them as "fixing" it; they were listed as det 1 race 0 fixed 0)
			% so i'm -1ing for "det bugs found" and "groups w any bugs found" to compensate
			Known false positives	& 1	& 2	& 6	& 5	& 3	& 1	& 0	& 18 	\\
			\hline
			Deteterministic bugs found ($\le$)
						& 5	& 20	& 18	& 18	& 30	& 15	& 9	& 115	\\
			Concurrency bugs found	& 5	& 56	& 19	& 15	& 13	& 9	& 5	& 122	\\
			\bf Total bugs found	& 10	& 76	& 37	& 33	& 43	& 24	& 14	& 237	\\
			\hline
			Deteterministic bugs fixed ($\ge$/$\le$)
						& 3	& 19	& 9	& 15	& 22	& 12	& 7	& 87	\\
			Concurrency bugs fixed ($\ge$)
						& 1	& 38	& 11	& 10	& 12	& 4	& 4	& 80	\\
			\bf Total bugs fixed	& 4	& 57	& 20	& 25	& 34	& 16	& 11	& 167	\\
	\end{tabular}
	\end{center}
	\caption{Landslide's bug-finding performance across all semesters of the CMU 15-410 study.
	$\le$ marks possible overcounts on account of unidentified false positives;
	$\ge$ marks possible undercounts on account of students not necessarily re-running to verify bugfixes.}
	\label{tab:cmubugs}
\end{table}

This table also shows the number of false positives, as self-reported by Landslide
as described in \sect{\ref{sec:education-pebbles-instrumentation}}.
% nah where actually doing this, where making it hapen.
%which I started counting since S'16.
Please note that this approach %to counting false positives
is limited by Landslide's ability to classify them as ``suspicious'', even if not definitely bugs.
These represent technical obstacles that either
prevented Landslide from being able to analyze the synchronization in play
(e.g., recursive {\tt mutex\_lock()} invocations)
or were deemed too pedagogically important to allow the student to proceed without fixing
(e.g., busy spin-wait synchronization loops)%
\footnote{
Landslide was also configured to issue a ``bug report''
if {\tt MAGIC\_BREAK}, a Simics debugging trap, was ever invoked,
with specific instructions to remove it;
I consider these closer to compilation errors than bug reports and so did not count them
even among false positives.
}.
Other false positives arising from bugs in Landslide itself
required more individual effort to confirm;
to classify these I relied on students to help report them
during the Landslide clinics and survey,
and I report on them in \sect{\ref{sec:education-eval-survey-falsepositives}}.
Nevertheless,
I am aware of no case in which Landslide would output a false positive {\em nondeterministic} bug%
% this is really shaky
% not too shaky for me to feel uncomfortable publishing it
% but too shaky for me to feel comfortable leaving it unsaid too...
\footnote{I did observe Landslide's ``suspicious loop'' message
for mutex or condition variable internal synchronization
on subsequent interleavings,
but at this point the line between style warning and true bug blurs,
so I count these as true concurrency bugs for simplicity's sake.}
% i *think*.......
all unexpected false positives encountered to date were output on the first interleaving tested.
Though unscientific,
this provides some reassurance that Table~\ref{tab:cmubugs}'s count of concurrency bugs is accurate,
even if the deterministic tally may overcount slightly.

Overall, Landslide helped students find and fix a lot of bugs.
Among the 103 participating groups across all semesters,
roughly three-quarters received bug reports,
slightly more than half were able to verify their fix for at least one,
and one-third overall were able to fix and verify {\em all} such bugs.
Even avoiding the deterministic bug series for their possible overcounting,
it's fair to conclude that Landslide caused at least two-thirds
of all concurrency bugs it found to get fixed before project handin.
Although lacking a control experiment to make a more scientific comparison of these tallies,
I tentatively conclude the students were well rewarded for their time.
The questions of whether it was a {\em better} use of time than conventional stress testing,
or whether the bug-finding resulted in submissions that course staff also, independently, judged better than before,
are beyond the scope of just tallying bug reports and will instead be addressed in the upcoming sections.

\subsubsection{University of Chicago}
\label{sec:education-eval-bugs-uc}

Among the 4 groups volunteering to receive Landslide bug reports,
Landslide found 7 distinct bugs across the two projects tested,
which I confirmed by manual inspection of the students' code.
3 among these were deterministic; 4 were concurrency bugs.

\begin{enumerate}
	\item {\tt priority-donate-multiple} found 2 deterministic bugs and 1 concurrency bug in {\tt threads} projects.
		It found the targeted priority-leak bug
		described in \sect{\ref{sec:education-pintos-tests}} in one group,
		and exposed a {\tt NULL} pointer dereference in another
		when attempting to donate priority to an already-exited thread.
		The latter bug was observed deterministically.
		A third group neglected to implement priority donation during {\tt lock\_release()} at all,
		which was found deterministically.
	\item {\tt wait-simple} found 1 deterministic bug and 3 concurrency bugs in {\tt userprog} projects.
		It found a deadlock in two groups' implementations,
		% for one of these groups 2 traces were emitted, same stack trace (i think) just w diff pps
		wherein {\tt process\_wait()}'s synchronization assumed
		it would always be called before {\tt process\_exit()}.
		In both cases the former's {\tt cond\_wait()} call would block forever if reordered after the latter.
		This test also found several use-after-frees for one of those two groups,
		whose {\tt process\_exit()} didn't guard against a parent process exiting and deallocating its memory first;
		this root cause manifested in heap errors in 3 separate locations.
		Lastly, the heap checking alone found a deterministic use-after-free in a third group's {\tt exec()}.
\end{enumerate}

In cases where Landslide emitted multiple preemption traces for the same bug,
whether manifested the same way through different combinations of preemption points
or with different stack traces entirely arising from the same underlying cause,
course staff sent them all to the student,
including a disclaimer along the lines of ``some of these may indicate the same bug''.
%
The other 3 tests ({\tt priority-sema}, {\tt alarm-simultaneous}, and {\tt wait-twice}) found no bugs among the 4 groups.
One of the 4 groups had no bugs found among any of the 5 tests.

%%%% i think it'd be an irb violation even to publish this at all
%%%% and the potential false positive issue makes it not scientific enough anyway
% Course staff also ran Landslide with wait-simple on submissions from even non-volunteers;
% in total finding bugs in 11 of 18 groups' implementations.
% However, ... couldn't check for false positives because of non-volunteers ...

% XXX: did I actually fix this in the annotation scripts, or did I just hackhackhack to suppress it?
% my email to kevin just says, "I suppressed the false positive bug from wait-simple in GROUPNAME and got a real bug next".
Finally, one false positive was found while testing {\tt wait-simple}.
Landslide mistakenly reported that {\tt free()} had reentered {\tt malloc()}
(normally a heap corruption bug)
due to a technical discrepancy between when Pebbles and Pintos kernels take and release their heap allocation lock
with regard to the rest of the {\tt malloc()} API.
After discarding the bug report and suppressing this false positive,
Landslide moved on to find the true deterministic use-after-free described above.

Overall, I consider all of the 7 bugs to be ``severe'':
they are all either correctness violations (priority mis-donation)
or stability issues (crash, deadlock, or data corruption).
In the upcoming survey section, although only one of these groups did the survey,
they enthusiastically reported finding the bug reports very helpful (\sect{\ref{sec:education-survey-response-data}}).
Despite the small sample size, I consider this
a positive result that Landslide can be a useful concurrency programming aid
even at universities with different class projects than CMU's 15-410.
As for 1 false positive among 8 total reports,
it is difficult to decide what rate of mistakes is acceptable \cite{racerx},
but the upcoming survey results (\sect{\ref{sec:education-eval-survey-falsepositives}})
suggest that students are generally at least patient enough
to ask for help and make progress when technical support is available.

\subsection{Student submission quality}
\label{sec:education-eval-grades}

I evaluated Landslide's impact on student submission quality in two ways.
First, I analyzed CMU 15-410 students' overall project grades between Landslide users and non-users
to see if Landslide helped them submit overall better implementations.
Second, I studied
several % how many? 3
individual concurrency bugs that I thought Landslide was likely to detect,
all already well-known by 15-410 course staff,
to see if using Landslide correlated with submitting projects absent those bugs.
% results..?

\subsubsection{Impact on grades}

Figure~\ref{fig:photo-of-ze-studence} shows the distribution of project grades in CMU 15-410
between Fall 2013 and Spring 2018,
grouped by whether or not students used Landslide during P2.
%Landslide users are, of course, the experimental group.
The control group is further distinguished by whether Landslide was offered and the student delined,
or whether the study had not yet begun that semester.
Intuition suggests students in the former group would be more likely
to be struggling too much to have free time to volunteer in the first place,
so any comparison between them and Landslide users is vulnerable to selection bias.
The latter control group, in which students did not have a choice, mitigates such bias,
although itself may be vulnerable to other confounding factors such as grading criteria varying across semesters.
%
In addition to P2, I also show the subsequent P3 (kernel project) grades,
likewise broken down by who used Landslide previously during P2%
\footnote{The P3 distributions are slightly smaller than those from P2,
as some students dropped the class in between.
A few new project groups also tend to appear during P3 as partnerships break up or re-form,
but as these cannot be meaningfully classified as Landslide users or non-users,
I omit them here.
}.
Should Landslide use be correlated with higher grades on the later, more difficult concurrency project,
one explanation could be that Landslide helped students
internalize new concurrency programming and/or debugging skills
that would help them write better kernels even without Landslide's aid.

% colors for these graphs in case i need them again later
% #720000
% #FF8080
% #804700
% #FFC680
% #807700
% #FFF780
% #004111
% #ABFFC0
% #002680
% #ABC3FF
% #3B0144
% #F4ABFF
\begin{figure}[p]
	\begin{center}
	\begin{tabular}{c}
		\includegraphics[width=0.96\textwidth]{photo-of-ze-studence.pdf}
		\\
		(a) Visualized as a box plot.
		\\
		\\
		\includegraphics[width=0.96\textwidth]{photo-of-ze-studence-cdf.pdf}
		\\
		(b) Visualized as an EDF.
	\end{tabular}
	\end{center}
	\caption{Distribution of project grades among Landslide users and non-users.
	Study participants (S'15 to S'18) are the experimental group;
	non-participants (S'15 to S'18)
	and students from semesters predating the study (F'13-F14) are the control groups.
	Compare Landslide users to non-users for within-semester differences,
	or users to pre-study students to mitigate selection bias.}
	\label{fig:photo-of-ze-studence}
\end{figure}

%We see in Figure~\ref{fig:photo-of-ze-studence} that, indeed,
{\bf Results.}
Overall,
Landslide users did indeed receive
slightly better grades overall on their P2 submissions than non-users from the same semesters,
although comparing to the second control group lessens that difference somewhat.
The difference is also less pronounced in both comparisons among P3 grades
(unsurprisingly, as learning lasting lessons about concurrency should be harder than fixing bugs case-by-case),
although the experimental group still maintains a tiny lead.

{\bf Statistical significance.}
Table~\ref{tab:table-of-ze-studence} presents more detailed statistics from these six distributions.
To assess whether the differences are significant,
% discrete == I used "version 2" of the AD metric that R outputs.
I use the $k$-sample discrete Anderson-Darling test \cite{anderson-darling}
provided by the R statistical programming environment \cite{r-lang,r-ksamples},
comparing each control group to its respective Landslide users group.
Anderson-Darling uses a weighted sum-of-squares distance metric to compare two empirical distribution functions (EDFs),
testing if their samples are likely to arise from the same underlying, though unspecified, distribution%
\footnote{
I choose Anderson-Darling over the simpler Kolmogorov-Smirnov,
which computes only the maximum instantaneous distance between EDFs,
for its better sensitivity both to repeated deviations and to tail differences \cite{beware-kolmogorov-smirnov}.
}.
I deem the difference in grade distributions significant when $p < 0.05$%
% FIXME: is this legit?
% should you be using the combined/grouped AD test instead?
\footnote{I choose not to correct for multiplicity among these 4 comparsons like I do in the next section
because each tests a different hypothesis.
},
although as in all $p$-value calculations,
this rejects only the null hypothesis,
not alternative hypotheses (such as ``more skilled students were more likely to sign up to begin with'').
%
Ultimately, only the same-semester comparisons were significant.
However, the distribution difference being smaller in P3 than in P2
% TODO: is there a good way to treat this more formally?
further suggests, albeit informally,
that the larger impact in P2 is not attributable {\em only} to selection bias,
or else a similar difference should have been observed in the P3 distributions.
Another possible explanation could simply be that the P3 grading critera result in less grade variance overall,
but this pattern is at least consistent with the ``Landslide helps students submit better P2s'' hypothesis.

% P2 vs opt out
% ad.test(c(42,50,51,55,55,55,56,57,57,58,62,63,65,65,65,66,66,67,67,67,67,68,68,68,69,69,69,70,70,70,70,71,71,71,71,71,72,73,73,73,73,73,74,74,74,74,74,75,75,75,75,75,76,76,76,76,77,77,77,77,78,78,78,78,79,79,79,79,79,79,80,80,80,80,81,81,81,82,82,82,82,83,84,84,84,84,85,85,86,86,86,87,87,87,88,89,90,91,92,93,93,93,95),c(38,42,47,50,52,53,55,55,56,56,57,57,58,58,58,59,60,60,61,62,62,63,65,66,67,68,68,68,69,70,70,70,71,71,71,71,72,72,72,73,73,73,74,74,76,76,76,77,77,78,78,78,78,78,79,79,79,80,82,82,83,83,83,84,84,84,85,86,86,88,90,95,96,96))
%               AD  T.AD  asympt. P-value
% version 1: 2.611 2.139          0.04263
% version 2: 2.760 2.334          0.03576

% P2 vs pre study
% ad.test(c(42,50,51,55,55,55,56,57,57,58,62,63,65,65,65,66,66,67,67,67,67,68,68,68,69,69,69,70,70,70,70,71,71,71,71,71,72,73,73,73,73,73,74,74,74,74,74,75,75,75,75,75,76,76,76,76,77,77,77,77,78,78,78,78,79,79,79,79,79,79,80,80,80,80,81,81,81,82,82,82,82,83,84,84,84,84,85,85,86,86,86,87,87,87,88,89,90,91,92,93,93,93,95),c(43,51,56,58,59,59,60,60,60,60,61,61,61,61,61,62,63,65,65,67,68,69,69,69,69,69,70,70,71,71,72,72,72,73,73,73,73,73,74,74,74,74,75,75,75,75,76,76,77,77,77,77,78,78,78,78,79,79,79,80,80,80,80,80,80,81,81,82,82,82,82,83,83,83,83,85,85,85,86,86,86,88,88,90,92,93))
%                AD    T.AD  asympt. P-value
% version 1: 0.6248 -0.4977           0.6275
% version 2: 0.5790 -0.5581           0.6706

% P3 vs opt out
% ad.test(c(65,68,72,72,73,75,75,75,75,75,76,76,76,76,77,77,77,78,78,78,78,78,78,79,79,80,80,80,80,80,80,80,81,82,82,82,82,82,82,83,83,83,83,83,83,84,84,84,84,85,85,85,85,85,85,85,85,85,85,86,86,86,86,86,86,86,86,86,87,87,87,88,88,88,88,88,89,89,89,89,90,90,90,90,90,91,91,91,91,92,92,93,93,93,94,95,95,95),c(45,58,66,70,72,72,72,72,74,74,74,75,75,75,77,77,77,78,78,78,79,79,79,79,79,79,80,80,81,81,81,81,81,82,82,83,83,83,83,83,84,84,85,85,86,86,86,86,86,86,87,87,89,89,89,89,90,91,91,91,93,93,93))
%               AD  T.AD  asympt. P-value
% version 1: 2.645 2.186          0.04085
% version 2: 2.800 2.398          0.03376

% P3 vs pre-study
% ad.test(c(65,68,72,72,73,75,75,75,75,75,76,76,76,76,77,77,77,78,78,78,78,78,78,79,79,80,80,80,80,80,80,80,81,82,82,82,82,82,82,83,83,83,83,83,83,84,84,84,84,85,85,85,85,85,85,85,85,85,85,86,86,86,86,86,86,86,86,86,87,87,87,88,88,88,88,88,89,89,89,89,90,90,90,90,90,91,91,91,91,92,92,93,93,93,94,95,95,95),c(60,62,63,64,65,68,69,72,72,74,75,75,75,75,76,76,77,77,77,77,78,78,78,78,79,80,80,80,81,81,81,81,81,81,81,81,81,82,82,82,82,82,83,83,83,84,84,84,84,84,84,84,85,86,86,87,87,88,88,88,88,88,88,88,89,89,89,89,89,90,90,90,91,92,92,92,93,94,95))
%               AD   T.AD  asympt. P-value
% version 1: 1.552 0.7322           0.1633
% version 2: 1.640 0.8473           0.1454

% TODO: should i report the T.AD test statistic? i don't even know what it means lol but it's what the p value like, applies to..
\begin{figure}[t]
	\begin{center}
		\small
	\begin{tabular}{cr|r||r|r|r|r||cc}
		\multicolumn{3}{c||}{Population} & \multicolumn{4}{c||}{Distribution} & \multicolumn{2}{c}{Significance} \\
		project	& group		& $N$	& min	& mean	& max	& stddev & $p$ vs users	& cutoff\\
		\hline
			& users		& 103	& 42	& 74.8	& 95	& 10.3	& --	& --	\\
		P2	& non-users	& 74	& 38	& 70.9	& 96	& 12.5	& 0.036	& 0.05	\\
			& pre-study	& 86	& 43	& 73.7	& 93	& 9.7	& 0.671	& 0.05	\\
			\hline
			& users		& 98	& 65	& 83.6	& 95	& 6.2	& --	& --	\\
		P3	& non-users	& 63	& 45	& 80.7	& 93	& 8.2	& 0.034	& 0.05	\\
			& pre-study	& 79	& 60	& 81.6	& 95	& 7.6	& 0.145	& 0.05	\\
	\end{tabular}
	\end{center}
	\caption{Detailed statistics from student grade distributions.}
	\label{tab:table-of-ze-studence}
\end{figure}

\subsubsection{Common bugs}

Based on the results from \sect{\ref{sec:education-eval-bugfinding}},
I selected 3 bugs to study in depth to ascertain
whether Landslide played an instrumental difference
in helping the students ultimately submit respecitvely correct implementations.
To avoid bias of picking too obscure and/or trivial bugs
that Landslide alone might find but even course staff would not expect students to solve,
I chose only bugs which had substantial penalties in the grading rubric
(guided, as well, by my own intuition as a former TA).
While checking for a given bug's presence by manual inspection,
I blinded myself to whether each group had been a Landslide user or not.
After unblinding, I then also classified students who used Landslide in general,
but whose usage snapshots showed they did not run the test case in question,
as non-users.
Table~\ref{tab:eval-common-bugs} presents the results.

% TODO check figure placement
\begin{table}[t]
	\begin{center}
		\small
		% note, smth like fisher's method which "combines independent studies into a meta p value"
		% cannot be applied for 2 reasons - 1, these are not independent since the same studence were
		% running landslide on diff test cases; 2, the hypotheses tested are different for each test
	\begin{tabular}{r||cc|cc|c|cc}
		& \multicolumn{2}{c|}{Landslide users} & \multicolumn{2}{c|}{Non-users} & & \multicolumn{2}{c}{Significance} \\
		Bug name & \#correct & \#buggy & \#correct & \#buggy & $\Delta$\%correct & $p$ & cutoff \\
		\hline
		% calculated in R with (rest done similarly)
		% fisher.test(rbind(c(11,10),c(0,16)))$p.value --> 0.0006143355
		Exit UAF  & 11	& 10	& 0	& 16	& $+52.4\%$	& 0.00061	& 0.0125 \\
		% I spent more hours than I should have looking at this landslide-user-submitted-bug group.
		% Can't name the group here, but if you're me and have the analysis file,
		% it's the one from f17 with 3 mutex_test logs in their "u" snapshot.
		% It seems like a mysterious older version of mutex.c was present between their "a" version
		% (which gets fully verified very quickly btw) and their "u" version (for which landslide finds
		% the bug in just a few minutes, no way they could be running the same version; plus the reported
		% data races and so on in the id-log don't line up to the "u" mutex.c source code).
		% That version was correct, then they updated it (introducing the bug) and ran a bunch of other tests,
		% but didn't rerun mutex_test to check their update.
		Mutex     & 17	& 1	& 19	& 0	& $-05.6\%$	& 0.49 & 0.0125 \\
		Paraguay  & 20	& 1	& 14	& 2	& $+07.7\%$	& 0.57	& 0.0125 \\
		% nb this is any downgrade related ding not just dgr tested ones
		% nb 1 is msising bc didn't impl rwlock at all
		% nice p value
		Downgrade & 10	& 3	& 19	& 4	& $-05.7\%$	& 0.69	& 0.0125 \\ % TODO verify these buggers by eye
		\hline
		Total     & 58	& 15	& 52	& 22	& $+09.2\%$	& 0.26 & 0.05 \\
	\end{tabular}
	\end{center}
	\caption{
		Correlation of student Landslide use with solving certain bugs in their final submission during Fall 2017 and Spring 2018 semesters.
		%$\Delta P$ stands for the increased probability that a student submits a correct P2
		%(with respect to each bug)
		%given that they use Landslide, versus not using Landslide; i.e.,
		%$P(\mathsf{correct}|\mathsf{landslide})-
		%P(\mathsf{correct}|\mathsf{non\text{-}user})$.
		Note that the totals in the bottom row double-count students,
		which makes sense only if you believe the incidences of each bug in a given submission are independent.
		}
	\label{tab:eval-common-bugs}
\end{table}

Each of the four bugs studied is described in detail as follows.

\begin{enumerate}
\item
{\bf Thread exit/join.}
%During {\tt thr\_exit()}, a thread must arrange for its own stack memory allocation to be freed.
After a thread finishes exiting, the memory allocated for its stack must be reclaimed for other uses.
A common student pitfall is to allow a {\tt thr\_join()}ing thread to free said stack space
while {\tt thr\_exit()} is still executing userspace C code (which typically accesses the stack),
or even for {\tt thr\_exit()} to free it itself.
In the former case, threads must interleave specifically to exhibit a use-after-free;
in the latter case, the use-after-free will be deterministic (i.e., present in all interleavings).
Other than Landslide, however,
the students have no Valgrind-like heap debugging tool which would report a bug immediately upon any illegal heap access.
This means that a subsequent {\tt thr\_create()} invocation would need to race to recycle the memory for a new thread stack
and conflict with the old exiting thread
before any problem could be detected.
The Landslide-friendly test {\tt thr\_exit\_join} is most likely to expose this bug.
%, and
%the official rubric penalties for it are
%{\tt :exit\_stack\_minor -2} and
%{\tt :exit\_stack\_race -4} (for the former case),
%and
%{\tt :uses\_freed\_stack -4} (for the latter).
Whether each student's submitted P2 solved or suffered from this bug
was determined by me personally inspecting their code.
%

\item
{\bf Mutex.}
The Landslide-friendly {\tt mutex\_test} checks for the possibility of two contending threads
accessing a mutex-protected critical section simultaneously.
It includes one thread repeating once the lock, unlock cycle
so as to check for unlock/lock races as well as lock/lock interactions,
and, as mentioned previously (\sect{\ref{sec:education-pebbles-tests}}),
checks for data race access pairs inside the mutex implementation itself.
As students have free rein to design their mutex internals,
this refers to the general class of mutex bugs
in which any number of things can go wrong, depending on the implementation,
leading to mutual exclusion (or otherwise assertion) failure.
Whether each student's submitted P2 solved or suffered from such bugs
was determined by checking their grade files for any mutex-related penalties (assessed by the TAs),
then me double-checking their implementations by hand to confirm.

Further investigating the one group who submitted a buggy mutex,
I found that while they had run {\tt mutex\_test} in Landslide
(and even found and fixed a separate deterministic bug already),
Landslide found no bug in what was presumably a correct implementation,
then they updated % oops
their code, introducing the bug, without testing it again thereafter.
% TODO: some reference to future work about focused testing on small code diffs? talked with some industry ppl about this sometimes

\item
{\bf Paraguay.}
Named after Ivan Jager, the Paraguayan 15-410 TA from 2004-2006 who originally discovered it,
this refers to a condition-variable bug in which a thread which sequentially sleeps on two different condition variables
can spuriously wake up early from the second {\tt cond\_wait()}.
The precise reasoning why many na\"ive student implementations are susceptible to this bug,
as well as the 3 major ways of fixing it,
are one of 15-410 staff's closely-guarded secrets (\sect{\ref{sec:410-secrecy}}).
Suffice it to say that the {\tt paraguay} test invokes a custom Pathos {\tt misbehave} mode
which biases thread scheduling towards the interleaving required to exhibit the bug
(Landslide, of course, replaces this {\tt misbehav}iour with model checking).
% TODO: i feel tihs sentence should go somewhere else, like after analysis, not in description
Hence, despite being a subtle concurrency bug,
the official course test suite is likely to expose it,
so comparing how many Landslide users and non-users submitted this bug in particular
speaks more to the impact of Landslide's preemption traces
in helping students diagnose a bug that a ``stress'' test could already find.
%
{\tt paraguay} is itself the Landslide-friendly test for its eponymous bug.
%, and
%the official rubric penalty for it is
%{\tt :paraguay -2}.
Whether each student's submitted P2 solved or suffered from this bug
was determined by me personally inspecting their code.

\item
{\bf R/W-lock downgrade.}
In addition to the standard R/W-lock interface,
P2 requires students to implement {\tt rwlock\_downgrade()}:
called with the lock held in write mode,
the caller adopts the reader role instead,
allowing other waiting reader threads to proceed simultaneously,
all while allowing no waiting writers to access in between.
%
The Landslide-friendly test {\tt rwlock\_downgrade\_read\_test}
checks that readers are allowed simultaneous access after a downgrade with no possibility for deadlock.
Rather than one specific bug, this refers to any of several failures that can arise during a downgrade.
%, and
%the official rubric penalties related to this bug are
%{\tt :downgrade\_no\_signal -1},
%%{\tt :unsafe\_rwlock\_downgrade -2}, % don't think this is right; not relevant anyway (overlaps)
%{\tt :downgrade\_partial\_signal -1}, % not sure if this means smth that dgr would expose
%{\tt :downgrade\_\allowbreak{}lies -3}, and
%{\tt :downgrade\_deadlock -4}.
Whether each student's submitted P2 solved or suffered from these bugs
was determined by checking their grade files for any downgrade-related penalties
(which were assessed by the TAs as a result of their manual inspection, rather than mine).
\end{enumerate}

{\bf Statistical significance.}
% two-tailed btw, the better (more conservative) way,
% tbh one tailed might even make sense here, but it doesn't change anything wrt the cutoffs anyway, so
The $p$ values
in Table~\ref{tab:eval-common-bugs}
are calculated in R \cite{r-lang} using Fisher's exact test \cite{fishers-exact-test},
treating each row as an independent 2x2 contingency table.
I divide 0.05, the standard significance cutoff,
by the number of bugs to conservatively account for multiplicity \cite{xkcd-jellybeans}
(not that it matters with these $p$ values).

The latter three bugs' occurences are thoroughly uncorrelated with Landslide use,
the {\tt thr\_exit()} bug standing alone with extremely high significance:
not a single student who did not use Landslide these semesters
submitted a correct implementation.
This difference is easily explainable:
the class-provided unit tests already do a good enough job catching the other three bugs
that students are able to find and fix them before submission regardless of what testing tool they used.
(The high correct-to-buggy ratio corroborates this.)
Nevertheless, that does not mean Landslide is pointless for these bugs;
it may well have helped the students reproduce them more reliably and/or diagnose them faster.
This is merely a null result for submission quality, not necessarily a negative one,
and the next section will attempt to evaluate such quality-of-life improvements instead.

On the other hand, because the {\tt thr\_exit()} bug typically manifests as a use-after-free,
the official tests must stress the thread library until the memory in question gets re-allocated
in just the right interleaving pattern to cause corruption that leads to a visible crash.
Landslide, on the other hand, detects this error immediately upon access (\sect{\ref{sec:landslide-valgrind-mode}}),
with no need for such complex corruption conditions.
Among the 11 Landslide users who submitted correct {\tt thr\_exit()}s in this regard,
it is difficult to say whether the heap checking alone, together with unit or stress tests,
would have been sufficient,
or whether heap checking and model checking were both necessary in concert to help them.
To isolate the impact of model checking,
I simulated the students having access to a stand-alone heap checker
by checking only the first thread interleaving with Landslide
% TODO - more specific section reference - quicksand det/nondet drpps eval
(similarly to \sect{\ref{sec:quicksand-eval}})
and assuming the students would find and fix any such ``deterministic'' use-after-free bugs before the deadline.
Re-classifying these into the ``correct'' group,
the 11-10-0-16 distribution becomes
% fisher.test(rbind(c(11,10),c(3,13)))$p.value -->  0.0475343
11-10-3-13,
with a new $p$ value of 0.048:
still positively correlated, but no longer significant under the multiplicity-corrected cutoff.

\subsection{Survey responses}
\label{sec:education-eval-survey}

Analyzing only the raw technical data of how users interacted with Landslide can paint only part of the picture.
For one, offering students better testing and debugging tools
may not necessarily find strictly more bugs or help students submit more correct implementations
than with stress testing;
it may instead find the same bugs faster,
affording the students more free time apart from the project,
a quality-of-life improvement normally invisible to graders.
For two, Landslide's automatic snapshots, being captured at the time of issuing each preemption trace,
necessarily miss the student's subsequent experience interpreting them.
The surveys listed in sections \sect{\ref{sec:education-survey-pebbles}} and \sect{\ref{sec:education-survey-pintos}}
serve to probe these more qualitative aspects of the Landslide experience.

\begin{figure}[p]
	\begin{center}
		\begin{tabular}{cc}
			\includegraphics[width=0.42\textwidth]{survey1.pdf} &
			\includegraphics[width=0.42\textwidth]{survey2.pdf} \\
			\includegraphics[width=0.42\textwidth]{survey3.pdf} &
			\includegraphics[width=0.42\textwidth]{survey4.pdf} \\
			\includegraphics[width=0.42\textwidth]{survey5.pdf} &
			\includegraphics[width=0.42\textwidth]{survey6.pdf} \\
			\includegraphics[width=0.42\textwidth]{survey7.pdf} &
			\includegraphics[width=0.42\textwidth]{survey8.pdf} \\
			\includegraphics[width=0.42\textwidth]{survey9.pdf} &
			\includegraphics[width=0.42\textwidth]{survey10.pdf} \\
		\end{tabular}
	\end{center}
	\caption{Student survey responses.
	SD/D/N/A/SA stands for
	strongly disagree/\allowbreak{}disagree/\allowbreak{}neutral/\allowbreak{}agree/\allowbreak{}strongly agree.}
	\label{fig:survey}
\end{figure}

\subsubsection{Response data}
\label{sec:education-survey-response-data}

% probably just better to state this in one line of prose
%\begin{table}[h]
%	\begin{center}
%		\begin{tabular}{r|crcrc|r}
%			& \multicolumn{5}{c|}{\bf Did survey?} & \\
%			\bf Semester & & \bf yes & & \bf no & & \bf Total \\
%			\hline
%			CMU F'17	& &  7 & &  7 & & 14 \\
%			U. Chicago F'17	& &  1 & &  3 & &  4 \\
%			CMU S'18	& &  5 & &  5 & & 10 \\
%			PSU S'18	& & 11 & & 27 & & 38 \\
%		\end{tabular}
%	\end{center}
%	\caption{Survey response rate among students who used Landslide. (The survey was first introduced in F'17.)}
%	\label{tab:photo-of-ze-surveys}
%\end{table}

The response distributions for each of the surveys' multiple choice questions
are shown in Figure~\ref{fig:survey}.
In total, 28 students (or pairs thereof) answered the survey:
12 pairs from CMU, 15 individuals from PSU, and 1 pair from U. Chicago.
The first four questions/graphs focus on concrete debugging results,
and the latter six on the students' subjective opinions.
%
Note that two of the questions
(3 and 7 from \sect{\ref{sec:education-survey-pebbles}})
were not asked on U. Chicago's version of the survey,
so their corresponding graphs show only CMU and PSU response data.
Likewise, U. Chicago's questions 4, 8, and 10
(see \sect{\ref{sec:education-survey-pintos}}),
which were not asked on CMU's and PSU's surveys,
having only one respondent,
are not pictured;
the answers thereto were ``15-20 minutes'', ``Agree'', and ``Agree'', respectively.
Also, this respondent's answer of 6 on the ``How many bugs'' question
appears to indicate the total number of preemption trace files course staff sent them;
%despite the attached ``some of these may indicate the same bug'' disclaimer;
upon further investigation,
the 6 traces seem to represent 2 distinct bugs among them (\sect{\ref{sec:education-eval-bugs-uc}}).

The survey responses were very positive:
% TODO: corroborate this with snapshot data?
students reported being able both to diagnose and to verify as fixed the vast majority of Landslide's reported bugs,
compared Landslide favorably to stress testing,
and found the experience worthwhile and worth recommending.
Regarding the 100 CPU-hours question in particular,
students could approach answering it in several different ways:
do the three students who answered 20/80 think stress testing is four times as good as Landslide at finding bugs,
or that stress testing requires four times as long to reach the same point of diminishing returns?
On the other extreme, one student said they would spend all 100 CPU-hours on Landslide,
which probably speaks more to their enthusiasm than to a careful attempt to maximize expected number of bugs found
% TODO: put a section reference to future work resource exhaustion / malloc failure injection discussion
(even the author themself recognizes stress testing's advantages for certain types of bugs, such as resource exhaustion).
Nevertheless, the responses' overall bias to spending at least half the CPU-time on Landslide
shows clearly that the students found the experience worthwhile.

Three questions with open-form answers bear further discussion:
what kind of false positives Landslide reported, % oopsative mood
reasons they found it worth recommending to a friend,
and suggestions for improving the interface
(PSU only\footnote{I thought to ask this question too late for CMU's surveys.}). % ososugichatte

\subsubsection{False positives}
\label{sec:education-eval-survey-falsepositives}

Even though 71\% of students reported receiving no false positive bug reports,
the nature of Landslide's bug-detection algorithms is such
that it should ideally never report any correct behaviour as wrong,
so I consider those 29\% that did report such the most negative result among the survey responses.
They described their false positives, and I either make excuses or own up, as appropriate, as follows.
Reports from CMU students (Simics version):
\begin{enumerate}
	\item Landslide complained of a nonexistent {\tt MAGIC\_BREAK} (Simics debugging function),
		despite the student's code never invoking it.
		This arose because of technical confusion between the test program's and the shell's address spaces,
		and was subsequently fixed in commit 3f24d67 (Simics repository only).
	\item Landslide reported ``some weird errors'' and/or mysteriously crashed
		when multiple instances were run from the same directory.
		(Multiple students suffering this failure mode contacted me for support,
		though only one reported it on the survey;
		in some cases, I recall said weird errors manifesting as false positive invalid heap access reports.)
		Simultaneous Landslides can clobber certain auto-generated header and/or temporary files,
		scrambling its instrumentation and leading to chaotic behaviour.
		I introduced a guard against this in commit 977e8fb (Simics repository) and e49b5df (Bochs repository).
	\item One student reported ``Data races which I believe they are not''.
		Looking at their usage snapshots to corroborate this,
		these appear to be true, yet benign, data races
		(in several cases corresponding to {\tt mutex\_test}'s verification of the mutex's internal memory accesses
		(see \sect{\ref{sec:education-pebbles-tests}}));
		i.e., expected behaviour rather than false positive bug reports.
		Student confusion about these could be alleviated
		by improving Landslide's user interface messages when printing data race information. % uuuuuuu its already so
	\item One student complained of a bug report that, while truly a bug, showed wrong filenames and line numbers in stack traces,
		and (quite naturally) suggested that accurate stack traces would make debugging easier.
		Checking against their usage snapshots,
		this seems to be an issue with Simics's symtable's handling of assembly functions.
		The Bochs version handles these correctly.
\suspend{enumerate}
Reports from PSU students (Bochs version):
\resume{enumerate}
	\item Landslide kept crashing for one student while running {\tt paraguay}.
		This is likely due to exceeding Linux's process limit and/or exhausting the class's official VM's memory,
		which can arise when many data race candidates cause Quicksand to spawn many Landslide jobs,
		and (as on {\tt paraguay}) each with very large state spaces that must defer.
		I had been working on mitigating this problem just before beginning the PSU study;
		properly addressing it would involve improving Quicksand's memory-exhaustion detection code
		and/or making Quicksand at all aware of the process limit to begin with.
		(Note that this is not strictly a false positive bug report, just a Landslide crash.)
	\item Students who initialized child threads with a base pointer value of {\tt 0xffffffff} observed Landslide crashes,
		as it attempted to stack trace through that address and access memory that wrapped around the address space.
		(Several students contacted me about this via email, and I issued a prompt fix;
		one student later reported it on the survey.)
		Commits 0573e34 (Bochs repository) and 654f459 (Simics repository) fixed this bug.
		(Like the above, not actually a false positive bug report.)
	\item Landslide issued false invalid heap access reports to one student
		who had been using {\tt new\_pages} to allocate thread stacks ``very close'' to the {\tt malloc} heap.
		They reported this during the study and I fixed it promptly for them in commit 4a26da7 (Bochs repository only),
		then reminded me of it again in the survey.
	\item Finally, one student reported simply, ``Race condition''.
		Without the same usage snapshots to consult as I'd have for a CMU student,
		or more self-reported detail, I regrettably can offer no comment.
\end{enumerate}

\noindent
Overall, most of the issues reported as answers to the ``false positive'' survey question
were merely Landslide crashes or user interface confusion.
Those few truly erroneous bug reports (items 1, 2, and 7),
while certainly guilty of burdening % oopsative mood
students with worry over whether the problem is their own code or in Landslide itself,
are at least not too discouraging for two reasons.
Firstly, in all cases the students were able to recognize the false positive quickly enough to ask for help,
and I was able to deploy a fix and let them proceed before the project came due.
Secondly, each such problem, now having been fixed, will befall no future student again --
not to assert Landslide is completely bug-free now,
but it at least grows more and more stable with each passing semester.

\subsubsection{Reasons worthwhile}

After the question asking ``Would you recommend a friend taking OS next semester to use Landslide?'',
I asked the students for open-form reasons why or why not.
As the former question's answers were exclusively positive (only 1 student even answering ``no opinion''),
this question's answers turned out to be mostly praise.
3 students declined to answer this question;
I reproduce the rest, paraphrasing for clarity and brevity,%
\footnote{Note that most students used the term ``race condition'' rather than ``concurrency bug'',
as taught in CMU and PSU lecture material;
I replaced these while paraphrasing in accordance with \sect{\ref{sec:glossary}}.}
here.
Answers from CMU:

\begin{enumerate}
	\item Easy to use and found nontrivial bugs.
	\item Useful to find some bugs, but other types of bugs (such as memory corruption) were impossible to find with Landslide.
	\item Helped verify our atomic primitives were correct.
	\item Would recommend, but takes too much time (unclear, but seems to be referring to execution time rather than setup/usage).
	\item Tests are automated and can be left running for a long time.
		However, faced uptime issues with CMU's linux servers (which reboot every night);
		wished for a more reliable execution environment.
	\item It's pretty helpful. % lol
	\item Seems more reliable than stress tests.
	\item Found a bug not found by stress tests.
	\item Finds concurrency bugs with little effort that may be undiscovered otherwise.
		(Also provided some interface feedback here, since I didn't ask CMU students a separate question for such;
		see next section.)
	\item Easy to learn and a simple way to test our P2.
\suspend{enumerate}
Answers from PSU:
\resume{enumerate}
	\item Although didn't find any concurrency bugs for me, gives me more confidence about my code.
	\item Helpful, just didn't have a lot of time to use it.
	\item Does not make concurrency debugging easy, but definitely makes it easier.
	\item Helpful to find bugs you weren't previously aware of.
		Makes more sense to use an expressly designed tool rather than (unit/stress) tests.
	\item Helpful and easy to use.
	\item Saves time overall, can run long tests overnight.
	\item Helps to find concurrency bugs and their root causes better than stress tests.
	\item It finds the concurrency bugs you need to fix for full credit.
	\item Helpful for finding some uncommon bugs I hadn't found or wasn't looking for.
	\item Found bugs I kept overlooking, which may have taken many hours to find otherwise.
	\item Helpful for the difficult step from code being ``finished'' (scare-quotes theirs; presumably meaning ``feature-complete'')
		to getting rid of all concurrency bugs.
	\item Easy to use, setup taking no more than 5-10 minutes, and allowing being run overnight.
	\item Very efficient at concurrency testing.
		Stress test crash reports do not necessarily point to the root cause, due to memory corruption for example;
		plus bugs may not show up every time due to nondeterminism.
	\item Helped find a bug I wouldn't have found otherwise.
		Did not show an interleaving directly (i.e., did not issue a bug report),
		but reported a data race that turned out to be a concurrency bug upon inspection.
\suspend{enumerate}
Answer from U. Chicago:
\resume{enumerate}
	\item Found several subtle, legitimate bugs we wouldn't have easily caught otherwise, but made sense once revealed.
		Fixing them took little time but allowed us to proceed confidently on the next project.
		Often wished for Landslide to have been available to use during the next project as well.
\end{enumerate}

\noindent
Overall, students most commonly praised Landslide's ease of use,
its ability to find bugs that elude stress testing,
and the confidence instilled from verifying bugs had been fixed.

\subsubsection{Interface suggestions}

Lastly, I asked the PSU students for any feedback they might have on making Landslide's interface easier to use or understand.

\begin{enumerate}
	\item Requested for preemption traces to be more clear about the meaning of each stack trace in each table cell,
		and complained of inaccurate line numbers
		(likely referring to how the current behaviour indicates the line of code {\em after} a function call,
		corresponding to the {\tt call} instruction's pushed return address,
		rather than the function call itself). (This answer from CMU; see above.)
	\item Including a manual or tutorial would be helpful (presumably beyond the user guide's instructions, such as recapping the procedure shown in the lecture demo which wasn't written down anywhere).
	\item Don't print warnings about line length exceeding 80 characters (inconsistency between 15-410 and \psuos compilation options).
	\item ``It takes too long. But I guess that's impossible to fix.'' (Well, it's an open research problem to fix!)
	\item Preemption traces should explicitly indicate where in the interleaving the bug occurred.
		(Root cause identification is its own research area, but more detail is certainly possible.)
	\item Improve explanation of data races (in the user guide, perhaps).
	\item Happy with it as-is (3 students)
	\item No response (7 students)
\end{enumerate}

Though I did not ask this question on the CMU survey,
my experience handling student questions in person
suggests CMU students
also mostly wish for
better explanations of data races and
more detail and clarity in the preemption traces.
Though I present the formal definition of data races in the lecture
(\sect{\ref{sec:education-pebbles-recruiting}})
and refer back to it in Landslide's documentation and output,
showing a concrete example
in future iterations of the user guide
would go a long way to illustrate the abstract concepts.
The preemption traces could be improved by making it clearer that the stack trace in each cell of the table
represents executing the thread in question from wherever it previously left off (or its inception)
all the way until it reaches that stack trace, then preempting it to run another thread.
They could also easily report more diagnostic information;
for example, showing the sets of memory conflicts between each thread,
annotating the type of each preemption point (yield, mutex, or data race),
and/or indicating the adversarial memory access for each data race preemption point.

\subsection{Other universities}

Comparing the survey response distributions between CMU versus PSU and U. Chicago (Figure~\ref{fig:survey}),
CMU students tended to verify their bugfixes more often by re-running Landslide
and found the preemption traces easier to understand,
while PSU and U. Chicago students generally compared Landslide more favourably to stress testing
(while CMU students comparatively preferred the more nuanced answer that it depends on the type of bug),
and reported more often that helped them understand concurrency better.
%Considering that CMU's 15-410 demands substantial prerequisite concurrency experience
%(15-213 \cite{sigcse01:CSaPP} or equivalent)
%whereas PSU's \psuos is more towards the average between CMU's 15-213 and 15-410
%(tempering the difficulty of P2 to make it more accessible (\sect{\ref{sec:overview-psu}})),
Considering the higher demand CMU's 15-410 makes for prerequisite concurrency experience
compared to PSU's relative tempering of P2's difficulty to make it more accessible (\sect{\ref{sec:pebbles}}),
% TODO: represent u c in this sentents?
these trends seem to correlate with the different levels of preparation each course's students had,
showing that students of various skill levels
can each benefit from the experience in different ways.


% TODO

%%%%%%%%%%%%%%%%%%%%%%%%%%%%%%%%%%%%%%%%%%%%%%%%%%%%%%%%%%%%%%%%%%%%%%%%%%%%%%%%

\section{Discussion}

Human subjects research is inherently messy.
Each of the previous section's approaches to evaluating Landslide's educational value
was accompanied by some drawback which prevented it from being perfectly objective science,
but many of them presented tentatively positive results nonetheless.
Landslide helped many students find and fix many bugs (\sect{\ref{sec:education-eval-bugfinding}}),
but making a direct comparison to stress testing, the prior state of the art, is not straightforward.
%and not statistically significant,
Immediate improvement in students' project grades was observed,
(\sect{\ref{sec:education-eval-grades}}),
although statistical significance was lost when attempting to account for selection bias;
and impact on grades alone is a very narrow measure of pedagogical value anyway.
Landslide's debugging power was also found to be statistically significant
for the {\tt thr\_exit\_join} bug in particular.
Students responded overwhelmingly positively in the survey (\sect{\ref{sec:education-eval-survey}}),
although it is easy to imagine students being equally happy with a ``debugging tool'' that just tells them all the answers.

Nevertheless, I believe each of these partial results taken together
paint an overall picture of success:
students fixed their own bugs {\em and} were happy about it;
students were able to ask for help rather than be deterred by inevitable technical difficulties;
students provided intelligent feedback suggesting they truly understood the debugging process;
% well not the student who said "traces should point to where the bug is" but eh can't win em all c.c
et cetera.
%
The rest of this section will discuss the study's limitations and provide some perspective for the future.

\subsection{Bias}

% TODO - u can verify the bug just by taking survey graph #1 and comparing to landslide snapshots (sep. by groups who answered survey)
% anything else?

As long as an educational user study is run on a volunteer basis,
one cannot completely avoid selection bias:
those with enough free time to participate are more likely to be the most capable students already,
who are least in need of the extra debugging help.
This was especially evident in the annotation-required P3 study from \cite{landslide},
in which only 5 groups volunteered (15\% of the 34 total who submitted P3 that semester),
among which 2 had enough after just the annotation phase and did not continue to do any in-depth testing.
In contrast, since switching to the automatically-annotatable P2,
the participation rate rose to 58\% (Table~\ref{tab:photo-of-ze-studence})
among all P2-submitting groups.
Reaching over half the class could already be seen as a major step in mitigating said selection bias.

Beyond participation, I anticipated several forms of bias affecting the survey responses.
First of all, I suspected that students might feel compelled to be overly polite in their answers
(whether consciously or subconsciously).
I attempted to counteract this by concluding my survey link emails with,
{\em Please try to answer honestly rather than flatteringly--if any part
of the experience was bad for you, I want to hear about it to make
Landslide better in the future!}
It's also possible that survey respondents % respondence
were more likely to be those who enjoyed Landslide the most,
meaning I might not hear as much negative feedback as I should.
% trying to isolate this bias would be pretty fraught...
% possible flaws in table: students who found bugs might have had either a good (bug!) or bad (false positive) experience
% could studence who didn't find bugs also have had a good experience, like enjoying the verification guarantee
% think of how to say it presentably...

% TODO: is this ok to report without IRB approval? ask dave
I took no special measures to compensate for bias in gender, race, or being non-natively English-speaking
during recruitment or the survey.
Surveying students to measure any differences in these demographics between study participants and the overall class population
would have required mandatory survey participation,
and in turn, a more rigorous IRB approval process.
According only to my memory of students who attended the Landslide clinics
(\sect{\ref{sec:education-pebbles-recruiting}}),
the racial diversity was roughly representative of the class at large,
and the proportion of women I perceived was in fact somewhat higher than the overall gender ratio.
More scientific analysis of such statistics was deemed beyond the study's scope for now.

\subsection{Retrospect}

More than just trying to draw firm conclusions from
the opinions of students who are just learning concurrency to begin with,
student feedback in turn guided the constant development
of Landslide, and the experimental design itself,
as the semesters went by.
In this section I will fantasize about how I might have run more perfect experiments
granted the impossible wish of knowing then what I know now.

\subsubsection{Pebbles}

The survey was introduced into the study only in time for two semesters' worth of student responses,
after several iterations of collecting only usage and bug report snapshots.
Apart from the obvious improvement of having been surveying students from the beginning,
the following questions could have improved the survey.
\begin{enumerate}
	\item {\em Did you have any technical difficulties with Landslide that I had to intervene on,
		whether in person or over email?}
		(Some students reported this in the ``false positives'' question,
		although fewer than I helped overall, so others must have not mentioned it.)
		Comparing answers to this question across subsequent semesters would give a sense
		of how much Landslide's stability was improving over time
		and whether it was mature enough for unsupervised use in the future.
	\item {\em How well do you feel you understood the research challenges explained in the lecture?}
		and,
		{\em How well do you feel a user should need to understand same in order to benefit from Landslide's bug reports?}
		(Answers on a scale from ``Not at all; Landslide is a magic black box to me''
		to ``I'm ready to work on research in this field myself''.)
		These questions would help fine-tune the lecture material and user guide to maximize student comfort,
		and potentially also corroborate the claim that Landslide is accessible even to novice users.
	\item {\em What additional debugging information would you want displayed on the preemption traces?}
		Knowing now that interpreting preemption traces was a sticking point for many users,
		I would hope to identify the most wished-for features to know what to prioritize improving.
		This could also assess students' understanding of what kinds of information
		would or would not be reasonable for Landslide to record and report.
	\item {\em For each bug Landslide found in your code, how trivial or severe do you feel it was?}
		This would help get a sense of how the students regarded Landslide
		on a spectrum between annoying style check or as a life-saver,
		and potentially suggest options to make Landslide suppress certain types of bug reports.
		For example, it currently reports spin-wait loops in {\tt mutex\_lock()} as bugs
		with a special message referring students to relevant lecture slides,
		but it's possible refusing to test any code beyond until that bug was fixed
		might have made Landslide less useful overall.
\end{enumerate}

Some students (between 0 and 2 per semester)
emailed me during P3
to ask if they could test their kernels with Landslide just like their thread libraries.
I answered each by explaining that it would take more effort on their part,
and if they were still interested,
guided them through the annotation process on a case-by-case basis.
This process was not included in the IRB-approved study protocol,
so I collected no results from them.
If I had planned in advance, I could have supported this ``bonus stage'' officially,
and further surveyed the brave volunteers about how P3 Landslide could be made more generally accessible.

Finally, to evaluate whether the experience of using Landslide
left the students with any lasting lessons learned,
a follow-up survey could have been given one or two years later.
Such a survey would ask, for example,
{\em Have you encountered any debugging problems since finishing OS that made you wish for a tool like Landslide?}
and
{\em Do you feel the way you think about testing, debugging, and program correctness
has been influenced in any way by using Landslide?}.

\subsubsection{Pintos}

While part of the point of this experimental design was to evaluate Landslide as a grading tool in the hands of TAs,
I would be remiss to mention that I also feared the automatic annotation process would not be as robust as the P2 version.
Indeed, while helping Kevin get oriented with using Landslide,
I implemented
several
fixes/improvements to the setup scripts
as I found student kernels that failed to automatically annotate
(for example, those with {\tt ready\_list} changed to an array,
as described in \sect{\ref{sec:education-pintos-instrumentation}}).
Had I given Landslide directly to students that semester,
the students themselves would have had to email me for tech support.

I attribute the low participation rate of Pintos students to P2 students to two major factors:
one, not incentivizing the students to directly improve their grades
(instead offering only the vague promise of a ``learning experience'' debugging their code only after handin),
and two, not traveling to the university to introduce the research topic in an in-person lecture
(leaving the students potentially confused about what advantage, if any, was offered over stress testing).
Hypothetically, I could have achieved greater user study participation
either by offering extra credit to students
or by offering an autograder-like interface for students to receive bug reports before their deadlines instead of after
(either way requiring a more rigorous IRB review process).

\subsection{Future educational use}

%{\bf Pebbles.}
Now that I'm done collecting student usage data to publish as research results,
and no longer bound by the IRB's requirement that Landslide not be used to assign grades
lest it be seen as coercion to participate,
Professor Eckhardt and I have been discussing options to deploy it as an official part of 15-410's curriculum.
This section has already clearly shown students are capable of debugging with it on their own time,
and I believe it well-automated enough to supplement Fritz (the existing stress testing infrastructure)
in the class's grading process as well.
TAs could also, at their option, use Landslide by hand to confirm any bugs encountered during manual inspection
and/or write new Landslide-friendly tests to expose bugs not yet targeted by the 6 tests offered here.
Over the study's seven semesters,
I believe the stability of Landslide's instrumentation process has improved enough
to require little to no ongoing technical support anymore,
although Landslide-specific office hours may still prove helpful.
I am also willing to continue giving the guest lecture as long as proximity and the curriculum allow.
Future problems to address include grading bias
(i.e., students submitting blindly-hacked code that just barely passes Landslide,
even if not necessarily correct,
thereby gaming the autograder),
and improving usability to reach even the most struggling students
(i.e., those past 42\% who submitted P2s without participating in the study).

%{\bf Pintos.}
Regarding non-research use in Pintos classes,
Landslide can now handle a considerably wider variety of student implementation quirks
on account of the fixes from this time (\sect{\ref{sec:education-pintos-instrumentation}}).
In its original shape
(before the F'17 semester,
having only enough instrumentation necessary for the Pintoses used in \sect{\ref{sec:quicksand-eval}}'s experiments),
Landslide was already able to automatically instrument 18 out of the 21 {\tt threads} project submissions at U. Chicago.
I was able to quickly deploy a fix to make the annotation scripts handle the other 3 cases,
although such technical support is not something any TA would be able to do.
%
In its current shape I would recommend it for TA use grading,
but not necessarily directly to students without someone familiar with the codebase on immediate hand for tech support.
However, I also believe Landslide's success in these user studies,
provided me present to handle technical issues,
serves as testament for stateless model checking in general in the educational theatre.
While Pintos's kernel-level environment presents a unique challenge for concurrency testing,
other, more readily automatic model checkers for user-space programs,
such as dBug \cite{dbug-ssv},
could easily be used on other thread-library-like programming projects at any university.


\chapter{Transactions}
\label{chap:tm}

%\inspirationalquote{This paper presents earth-shaking work that deserves to be read by the entire computer science community.
%To that end, have you considered putting it in a thesis?}
%{Anonymous SIGBOVIK reviewer}

\inspirationalquote{
\begin{tabular}{p{0.47\textwidth}}
One can only build a sand castle where the sand is wet. \\
But where the sand is wet, the tide comes. \\
Yet we still build sand castles.
\end{tabular}}
{Yuri, Doki Doki Literature Club}

Transactional memory \revisionminor{(TM)~\cite{transactional-memory}} is a concurrency primitive
by which programmers may attempt an arbitrary sequence of shared memory accesses,
which will either be committed (made visible to other processors/threads) atomically,
such that no intermediate state modification is ever visible,
or, in the case of a conflict which would prevent such,
discarded with an error code returned to allow the programmer to \revisionminor{invoke} a synchronized backup path.
%
Transactional memory specifications typically have three API functions, abstractly speaking:
\begin{itemize}
	\item  {\sf \revisionminor{begin}} begins a transaction,
		staging any subsequent shared memory accesses in some temporary thread-/processor-local storage,
		and checking for conflicts with the accesses of any other threads or CPUs.
		\revisionminor{When the transaction is started, {\sf begin} returns a success code.}
		If the transaction is unsuccessful, as described below,
		{\sf begin} instead returns an error code
		indicating the programmer should fall back to some other, possibly slower, synchronization method.
	\item {\sf \revisionminor{end}} ends a transaction,
		attempting to commit all staged accesses to the shared memory atomically with respect to
		reads or writes from other concurrently-executing code.
		If any of those accesses conflict (i.e., read/write or write/write)
		with any other access to the same memory since the transaction started,
		they are instead discarded and execution state reverts to the {\sf begin} with an error code as described above.
	\item {\sf \revisionminor{abort}} explicitly aborts a transaction,
		regardless of any memory conflicts,
		discarding changes and reverting execution as described above.
		Some implementations allow an arbitrary abort code to be specified
		which will appear \revisionminor{as} {\sf begin}'s error code.
\end{itemize}

\subsubsection{Implementations}

% do they track changes in pthread-style TLS? or?
Software TM implementations (STM) typically function as a library,
tracking staged memory accesses in local memory,
and aborting whenever a conflict is detected between two transactions' tracked accesses \cite{stm-pldi06,mcrt-stm,swisstm}.
Hardware TM implementations (HTM) use processor-level hardware support,
which stages changes in per-CPU cache lines,
and may abort for STM's reason above \cite{intel-tsx-overview},
%\cite{htm-experience},
or additionally whenever a conflict is detected between one transaction's traced access and {\em any} other memory access,
or whenever a conflict occurs on the same cache line, not necessarily the same address,
or in case of any system interrupt or cache overflow.
Intel's commercial implementation of HTM has a rocky history of hardware bugs \cite{intel-disable-tsx,intel-tsx-update},
attesting to the feature's complexity and the need for formal verification on both sides of its API.

\subsubsection{Terminology}

The world of transactional memory is home to several more confusing acronyms similar to ``HTM''.
Transactional Synchronization Extensions (TSX)
refers to Intel's implementation of HTM on Haswell and more recent microarchitectures
\cite{htm-haswell}.
Restricted Transactional Memory (RTM)
refers to the {\tt xbegin}, {\tt xend}, and {\tt xabort} subset of TSX instructions,
which of course correspond to {\sf begin}, {\sf end}, and {\sf abort} listed above,
as well as {\tt xtest}, an instruction which returns whether or not the CPU is executing transactionally.
GCC and Clang expose these as C/C++ intrinsics named {\tt \_xbegin()}, {\tt \_xend()}, {\tt \_xabort()}, and {\tt \_xtest()}
\cite{htm-gcc}.
Hardware Lock Elision (HLE)
refers to the {\tt xacquire} and {\tt xrelease} subset of TSX instructions,
which extend the traditional interface to offer a slightly higher-level way to access the CPU feature,
optimized for simplicity for locking-like synchronization patterns
% other possible things to cite here:
% https://lwn.net/Articles/534758/
% https://software.intel.com/en-us/node/683688
\cite{speculative-lock-elision,lock-elision}.
In this thesis I focus on RTM, the more general (i.e., expressive (i.e., bug-prone)) interface,
and among all these acronyms restrict myself to ``HTM''
(when referring to transactional memory as a concurrency primitive in the abstract)
and ``TSX''
(when referring to Intel's implementation and/or GCC's intrinsics interface).
The non-pedantic reader may treat these as interchangeable.

\revision{Transactional atomicity comes in two flavours,
{\em strong} and {\em weak atomicity}~\cite{htm-subtleties}.
%
HTM implementations use the processor's cache coherence protocol to detect conflicts,
and hence any conflicting memory access from another CPU will cause an abort,
regardless of whether that access was itself transactional.
Hence, HTM transactions are {\em strongly atomic}, i.e., appear indivisible to all other code.
%
STM implementations rely on the programmer to use its interface to access any shared memory
she intends to protect with its transactions,
which cannot detect conflicts with non-transactional accesses that bypass it.
Hence, STM transactions are {\em weakly atomic}, i.e.,
its transactions will abort only from conflicts with other transactions,
while other code may interleave with them freely.
%
The manner in which non-transactional code may interleave with STM transactions
further depends on whether the STM uses
{\em eager versioning},
wherein transactional updates are immediately visible to non-transactional code,
or
{\em lazy versioning},
wherein such updates are visible only after commit
\cite[\S{}2]{stm-weak-atomicity}.
Unless explicitly noted as weak atomicity,
this chapter assumes strong atomicity,
and ``STM semantics'' refers to a strongly-atomic HTM program emulating STM with retry loops
(see \cref{sec:txn-abort-codes}).
Treatment of weak atomicity in
\cref{sec:tm-weak-atomicity} and \cref{sec:tm-impl-weak-atomicity}
assumes eager versioning.}

% or "Contribution"
\subsubsection{Motivation}

\revision{For the most part,
existing stateless model checkers
model interrupt-driven thread switching as the only source of nondeterminism in a program's execution.
In other words, assuming a fixed test input,%
\footnote{Symbolic execution \cite{symbolic}, a neighbouring research area,
addresses input nondeterminism; see \cref{related:symbolic}.}
the only way multiple behaviours can arise is by scheduling different threads at each preemption point.
Only recently have checkers emerged that extend this model by
incorporating other nondeterminism sources:
store buffer reordering arising from relaxed memory models
\cite{cdschecker,tsopso,rcmc}
and event-driven application programming models
\cite{r4,eventracer},
both of which are especially relevant in mobile programming.
Under TM, {\sf begin}'s return value is also nondeterministic,
and moreover may depend on execution events from the future in a parallel universe
(i.e., subsequent memory conflicts detected and reverted),
and no model checker yet exists which can model this.
In this chapter I will develop a concurrency model to incorporate abort nondeterminism
(\cref{sec:tm-design}),
including two equivalence proofs to help cope with the consequent state space explosion
(\cref{sec:tm-design-failure} and \cref{sec:tm-design-locks}),
extend the proofs to account for weak atomicity (\cref{sec:tm-weak-atomicity}),
extend Landslide to support testing TM programs under a variety of execution semantics
(\cref{sec:tm-implementation}),
and demonstrate its practicality
%\cref{sec:tm-eval}
both by finding new bugs
(\cref{sec:tm-eval-bugs})
and by verifying correctness
(\cref{sec:tm-verif})
in a variety of microbenchmarks and real-world HTM programs.}

\revisionminor{The following sections will make heavy reference
to the example TSX program from \Cref{fig:htm-example} (\cref{sec:overview-tm})
while defining the new concurrency model,
so I reproduce it here in \Cref{fig:htm-example-reproduced} for the reader's convenience.}

\begin{figure}[h]
	\begin{center}
		\begin{tabular}{cl}
			1 & \texttt{\flow{if} ((status = \call{\_xbegin}()) == \const{\_XBEGIN\_STARTED}) \{} \\
			2 & \texttt{~~~~x++;} \\
			3 & \texttt{~~~~\call{\_xend}();} \\
			4 & \texttt{\} \flow{else} \{} \\
			5 & \texttt{~~~~\call{mutex\_lock}(\&m);} \\
			6 & \texttt{~~~~x++;} \\
			7 & \texttt{~~~~\call{mutex\_unlock}(\&m);} \\
			8 & \texttt{\}} \\
		\end{tabular}
	\end{center}
	\caption{Example TSX program.}
	\label{fig:htm-example-reproduced}
\end{figure}

%%%%%%%%%%%%%%%%%%%%%%%%%%%%%%%%%%%%%%%%%%%%%%%%%%%%%%%%%%%%%%%%%%%%%%%%%%%%%%%%

\section{Concurrency model}
\label{sec:tm-design}

While up to now Landslide's tested programs' concurrency has been limited to timer-driven thread scheduling,
HTM presents a fundamentally new dimension of nondeterminism,
namely the hardware's ability to revert execution sequences
and the delayed visibility of changes to other threads.
In order to efficiently test HTM programs in Landslide,
in this section
I develop a simpler concurrency model and offer a proof of equivalence to HTM execution semantics.
I make two major simplifications:
simulating transaction aborts as immediate failure injections,
and treating transaction atomicity as a global mutex during data-race analysis;
and provide corresponding equivalence proofs.

{\bf Notation.} Let $I = TN_1@L_1, TN_2@L_2, ... TN_n@L_n$,
with $N_i$ a thread ID and $L_i$ a code line number,
denote the execution sequence of a program as it runs according to the specified thread interleaving.
This serialization of concurrent execution is told from the perspective of all CPUs at once
and hence assumes sequential consistency.
For discussion of relaxed memory models refer to \cref{sec:tm-warpzone-relaxed}.

\subsection{Example}

Consider again the program in \Cref{fig:htm-example-reproduced}.
Note that the C-style {\tt x++} operations, when compiled into assembly, %\cite{printy},
become multiple memory accesses which can be interleaved with other threads,
\revisionminor{as shown below in \Cref{fig:htm-example-assemblified}.}

\revision{
\begin{figure}[h]
	\begin{center}
		\begin{tabular}{ll}
			1 & \texttt{\flow{if} ((status = \call{\_xbegin}()) == \const{\_XBEGIN\_STARTED}) \{} \\
			%2 & \texttt{~~~~x++;} \\
			$2a$ & \texttt{~~~~temp <- x;} \\
			$2b$ & \texttt{~~~~temp <- temp + 1;} \\
			$2c$ & \texttt{~~~~x <- temp;} \\
			3 & \texttt{~~~~\call{\_xend}();} \\
			4 & \texttt{\} \flow{else} \{} \\
			5 & \texttt{~~~~\call{mutex\_lock}(\&m);} \\
			$6a$ & \texttt{~~~~temp <- x;} \\
			$6b$ & \texttt{~~~~temp <- temp + 1;} \\
			$6c$ & \texttt{~~~~x <- temp;} \\
			7 & \texttt{~~~~\call{mutex\_unlock}(\&m);} \\
			8 & \texttt{\}} \\
		\end{tabular}
	\end{center}
	\caption{\Cref{fig:htm-example-reproduced} with non-atomic operations shown as pseudo-assembly.}
	\label{fig:htm-example-assemblified}
\end{figure}

% T1 T2 T3 but latex doesn't allow #s in cmd names
\newcommand\ti{\ensuremath{\hilight{lavender}{\mathbf{T1}}}\xspace}
\newcommand\tj{\ensuremath{\hilight{seafoam}{\mathbf{T2}}}\xspace}
\newcommand\tk{\ensuremath{\hilight{salmon}{\mathbf{T3}}}\xspace}

\newcommand\tiat[1]{\ensuremath{\hilight{lavender}{\mathbf{T1}@#1}}\xspace}
\newcommand\tjat[1]{\ensuremath{\hilight{seafoam}{\mathbf{T2}@#1}}\xspace}
\newcommand\tkat[1]{\ensuremath{\hilight{salmon}{\mathbf{T3}@#1}}\xspace}

\begin{figure}[t]
	\begin{center}
		\begin{tabular}{ll}
			\hilight{lavender}{\bf Thread 1} &
			\hilight{seafoam}{\bf Thread 2} \\
			\hline
			\texttt{status = \call{\_xbegin}()} & \\
			\texttt{temp <- x;} & \\
			\texttt{temp <- temp + 1;} & \\
			& \texttt{status = \call{\_xbegin}()} \\
			& \texttt{x++;} \\
			& \texttt{\call{\_xend}();} \\
			\texttt{x <- temp;} & \\
			\texttt{\call{\_xend}();} & \\
		\end{tabular}
	\end{center}
	\caption[Impossible hypothetical interleaving within transactions of \Cref{fig:htm-example-assemblified}.]
		{Impossible hypothetical interleaving within transactions of \Cref{fig:htm-example-assemblified}.
		The effects of $\tiat{2a-2b}$ are never globally visible, as its transasction must abort,
		and moreover $\tiat{2c-3}$ will not execute at all.}
	\label{fig:htm-example-interleaving}
\end{figure}
}

If these instructions from the {\tt x++} in the transaction are preempted,
with another thread's access to {\tt x} interleaved in between,
the transaction will abort.
So, a hypothetical interleaving such as shown in \Cref{fig:htm-example-interleaving},
or equivalently written sequentially as:
\[
	\tiat{1}, \tiat{2a}, \tiat{2b}, \tjat{1}, \tjat{2}, \tjat{3}, \tiat{2c}, \tiat{3}
\]
or, henceforth further \revisionminor{simplified} for clarity:
\[
	\tiat{1-2b}, \tjat{1-3}, \tiat{2c-3}
\]
is not possible; rather, \ti will fall into the backup path:
\[
	\tiat{1-2b}, \tjat{1-3}, \tiat{4-7}
\]
However, the {\tt x++} operation from the failure path (correspondingly $6a$, $6b$, $6c$)
{\em can} be thusly separated with conflicting accesses interleaved in between,
since the mutex only protects the failure path against other failure paths,
but not against the transaction itself.
So (assuming {\tt x} is intended to be a precise counter rather than a sloppy one),
%losing one of the increments to which constitutes a bug),
the following interleaving exposes a bug:%
\footnote{Note also that this bug requires either at least 3 threads or at least 2 iterations between 2 threads to expose;
this highlights MC's dependence on its test cases to produce meaningful state spaces in the first place.}
\[
	\tiat{1-2b}, \tjat{1-3}, \tiat{4-6b}, \tkat{1-3}, \tiat{6c-7}
\]
Prior work \cite{tm-benchmark-cmu} proposed the idiom shown in \Cref{fig:htm-fixed}
to exclude this family of interleavings,
which shows that correctly synchronizing even the simplest transactions may be surprisingly difficult or complex.
%further motivating the research.

\begin{figure}[h]
	\begin{center}
		\begin{tabular}{ll}
		%\texttt{void count() \{} \\
		%\texttt{~~~~for (int i = 0; i < 1000; i++) \{} \\
		& \texttt{\ctype{bool} prevent\_transactions = \const{false};} \\
		\\
		0 & \texttt{\flow{while} (prevent\_transactions) \flow{continue};} \\
		1 & \texttt{\flow{if} ((status = \call{\_xbegin}()) == \const{\_XBEGIN\_STARTED}) \{} \\
		2 & \texttt{~~~~\flow{if} (prevent\_transactions)} \\
		3 & \texttt{~~~~~~~~\call{\_xabort}();} \\
		4 & \texttt{~~~~x++;} \\
		5 & \texttt{~~~~\call{\_xend}();} \\
		6 & \texttt{\} \flow{else} \{} \\
		7 & \texttt{~~~~\call{mutex\_lock}(\&m);} \\
		8 & \texttt{~~~~prevent\_transactions = \const{true};} \\
		9 & \texttt{~~~~x++;} \\
		A & \texttt{~~~~prevent\_transactions = \const{false};} \\
		B & \texttt{~~~~\call{mutex\_unlock}(\&m);} \\
		C & \texttt{\}} \\
		%\texttt{~~~~\}} \\
		%\texttt{\}} \\
		\end{tabular}
	\end{center}
	\caption[Variant of the program in \Cref{fig:htm-example-reproduced}.]
		{Variant of the program in \Cref{fig:htm-example-reproduced},
		with additional synchronization to protect the failure path from the transactional path.
		The optional line 0 serves to prevent a cascade of failure paths
		for the sake of performance
		by allowing threads to wait until transacting is safe again.}
	\label{fig:htm-fixed}
\end{figure}

\subsection{Modeling transaction failure}
\label{sec:tm-design-failure}

In the previous section's examples,
the way I stated interleavings such as $\tiat{1-2c}, \tjat{1-3}, \tiat{4-7}$%
\footnote{For a clearer example to follow, I have reordered \ti's write to {\tt x} before \tj's part, compared to before.}
glossed over
%the HTM execution semantics underlying how such a sequence of operations would be carried out.
how such a sequence of operations would be carried out under HTM.
For example, \ti's write during $2c$ is not actually visible to \tj,
although it would be under a thread-scheduling-only concurrency model.

Intel's official TSX documentation \cite{intel-tsx-overview}
summarizes its interface and behaviour in prose.
Recent work has used proof assistants to formalize some of the execution semantics
of x86 in general \cite{x86-semantics}
and of transactions in particular (at both hardware- and language-level) \cite{relaxed-transactions-pldi}.
However, state-of-the-art advances in model checking algorithms
still state their theorems and proofs in prose \cite{bpor,optimal-dpor,tsopso,satcheck,mcr,quicksand,rcmc},
so this section's proofs will regrettably do likewise,
leaving the rigor of mechanization to future work.
The reader may at least rest assured that the proofs herein
rely on transactional semantics that have themselves been formally verified.

To summarize HTM's execution semantics:

\begin{enumerate}
	\item Any modifications to shared state (such as $2c$) by a transacting thread
		are not visible to any other during the transaction
		(such as $2c$ in this example, despite \tj executing afterwards).
	\item All local and global state changes during a transaction
		(such as \ti's lines $1-2c$ in this example)
		are discarded when returning an abort code from {\tt xbegin}
		(jumping to line $4$, in this example).
\end{enumerate}

% there shouldbe a better word than 'requires' here.. benefits from.. "the point of it is"..
While use of HTM in production requires the performance advantage
of temporarily staging such accesses in local CPU cache,
model checking such programs need be concerned only with the program's {\em observable} behaviours.
I claim that MCing the simpler interleaving $\tiat{1}, \tjat{1-3}, \tiat{4-7}$
is an equivalent verification \revisionminor{to} MCing the one above;
in fact, this interleaving suffices to check
all observable behaviours
of all interleavings
of all subsets of $\tjat{1-3}$
with all subsets of $\tiat{2a-2c}$,
whether they share a memory conflict or not.
Stated formally, let:

\newcommand\tii{\ensuremath{\hilight{lavender}{\mathbf{Ti}}}\xspace}
\newcommand\tjj{\ensuremath{\hilight{seafoam}{\mathbf{Tj}}}\xspace}
\newcommand\tkk{\ensuremath{\hilight{salmon}{\mathbf{Tk}}}\xspace}

\newcommand\tiiat[1]{\ensuremath{\hilight{lavender}{\mathbf{Ti}@#1}}\xspace}
\newcommand\tjjat[1]{\ensuremath{\hilight{seafoam} {\mathbf{Tj}@#1}}\xspace}
\newcommand\tkkat[1]{\ensuremath{\hilight{salmon}  {\mathbf{Tk}@#1}}\xspace}

\begin{itemize}
	\item $\tiiat{\alpha}$ be an HTM begin operation,
	\item $\tiiat{\beta_1}\dots\tiiat{\beta_n}$ be the transaction body (with $\beta_n$ the HTM end call),
	\item $\tiiat{\phi_1}\dots\tiiat{\phi_m}$ be the failure path, and
	\item $\tiiat{\omega_1}\dots\tiiat{\omega_l}$ be the subsequent code executed unconditionally.
\end{itemize}
Note that arbitrary code may not be structured to distinguish these as nicely as in the examples;
e.g., more code may exist in the success branch after {\tt \_xend()};
such would be considered part of $\omega$ here.

Then, without loss of generality (for any number of other threads \tjj/\tkk,
\revisionminor{WLOG presented here as two threads,}
and for any number of thread switches away from \tii during the transaction,
\revisionminor{WLOG presented here as one switch}):
\vspace{1em}

\begin{lemma}[Equivalence of Aborts]
	\label{lem:equiv}
	For any interleaving prefix
	\[
	\begin{tabular}{c}
		$\tiiat{\alpha},\tiiat{\beta_1}\dots\tiiat{\beta_b},$\\
		$\tjjat{\gamma_1}\dots\tjjat{\gamma_j},$ \\
		$\tkkat{\kappa_1}\dots\tkkat{\kappa_k},$ \\
		$\tiiat{\beta_{b+1}}$ % $\dots\tiiat{\beta_{n-1}}$ -- excluded bc might abort
	\end{tabular}
	\]
	with $b<n$, $j \ne i$, $k \ne i$, et cetera, either:
	\begin{enumerate}
		\item $\tiiat{\alpha},\tjjat{\gamma_1}\dots\tjjat{\gamma_j},\tkkat{\kappa_1}\dots\tkkat{\kappa_k},\tiiat{\phi_1}\dots$
			(conflicting case), or
		\item $\tiiat{\alpha},\tiiat{\beta_1}\dots\tiiat{\beta_b}\dots\tiiat{\beta_n},\tjjat{\gamma_1}\dots\tjjat{\gamma_j},\tkkat{\kappa_1}\dots\tkkat{\kappa_k}$
			(independent case)
	\end{enumerate}
	exists and is observationally equivalent.
\end{lemma}

\begin{proof}
	Case on whether the operations by \tjj and/or \tkk have any memory conflicts (read/\allowbreak{}write or write/write)
	with $\tiiat{\beta_1}\dots\tiiat{\beta_n}$.
	If so, then the hardware will abort \tii's transaction, discarding the effects of $\tiiat{\beta_1}\dots\tiiat{\beta_n}$
	and jumping to $\tiiat{\phi_1}$,
	satisfying case 1.
	Otherwise, by DPOR's definition of transition dependence (\cite{dpor}; see also \cref{sec:landslide-dpor-hb}),
	$\tiiat{\beta_{b+1}}\dots\tiiat{\beta_n}$ is independent with the transitions of \tjj and \tkk,
	may be successfully executed until transaction commit,
	and reordering them produces an equivalent interleaving,
	satisfying case~2.
\end{proof}

The claim's second part follows naturally.
\vspace{1em}

\begin{theorem}[Atomicity of Transactions]
	\label{thm:atom}
	For any state space $S$ of a transactionally-concurrent program,
	an equivalent state space exists in which all transactions are either executed atomically or aborted immediately.
\end{theorem}

\begin{proof}
	For every $I \in S$ with $\tiiat{\alpha},\tiiat{\beta_1}\dots\tiiat{\beta_b},$ $\tjjat{\dots},\tkkat{\dots},\tiiat{\beta_{b+1}} \in I$,
	apply \Cref{lem:equiv} to obtain an equivalent interleaving $I'$ satisfying the theorem condition.
	The resulting $S'$ can then be MCed without ever simulating HTM rollbacks.
\end{proof}

\subsection{Memory access analysis}
\label{sec:tm-design-locks}

Next comes the issue of memory accesses within transactions with regard to data-race analysis (\cref{sec:background-datarace}).
\Cref{thm:atom} provides that the body of all transactions may be executed atomically within the MC environment.
While they may interleave between other non-transactional sequences,
no other operations (whether transactional or not) will interrupt them.
I claim this level of atomicity is equivalent to that provided by a global lock,
and hence abstracting it as such in Landslide's data-race analysis is sound.

Let $\tiiat{\mu},\tjjat{\nu}$ be a pair of memory accesses to the same address, at least one a write,
in some transactional execution $I$ normalized under \Cref{lem:equiv}.
Then let $\mathsf{lockify}_m(\tkkat{L})$ denote a function over instructions in $I$,
which replaces $\tkkat{L}$ with $\tkkat{\mathsf{lock}(m)}$ if $L$ is a successful HTM begin,
with a no-op if $L$ is a transaction abort,
or with $\tkkat{\mathsf{unlock}(m)}$ if $L$ is an HTM end,
or no replacement otherwise.
Finally, let $I' = \exists m. \mathsf{lockify}_m(I)$,
the execution with the boundaries of all successful transactions replaced by an abstract global lock $m$.
\Cref{lem:equiv} guarantees mutual exclusion of $m$.
\vspace{1em}

\begin{theorem}[Transactions are a Global Lock]
	\label{thm:global-lock}
	$\tiiat{\mu},\tjjat{\nu}$ is a data race in $I$ iff it is a data race in $I'$.
\end{theorem}

\begin{proof}
I prove one case \revisionminor{for each of the two variant definitons of data races that Landslide supports
(\cref{sec:landslide-datarace})}.
%\cite{quicksand}. % wtf no
For each, I state below what it means to race in an execution with synchronizing HTM instructions.

\begin{itemize}
	\item {\bf Limited Happens-Before.}
		%Under this definition,
		To race in $I$ they must be reorderable at instruction granularity,
		at least one with a thread switch immediately before or after
		\cite{tsan,hybriddatarace}.
		\begin{itemize}
			\llitem $I \Rightarrow I'$:
				If $\tiiat{\mu},\tjjat{\nu}$ race in $I$,
				then they cannot both be in successful transactions,
				or else placing \tiiat{\mu} within the boundaries of \tjjat{\nu}'s transaction
				would cause the latter to abort, invalidating \tjjat{\nu}, or vice versa.
				Hence they will not both hold $m$ in $I'$.
				Otherwise their lock-sets and DPOR dependence relation remain unchanged.
			\llitem $I' \Rightarrow I$:
				If $\tiiat{\mu},\tjjat{\nu}$ race in $I'$,
				both corresponding threads cannot hold $m$;
				WLOG let $\tii$ not hold $m$ during $\tiiat{\mu}$.
				Then in $I$, $\tiiat{\mu}$ is not in a transaction.
				With the remainder of their lock-sets still disjoint,
				and still not DPOR-dependent, $\tjjat{\nu}$ (or its containing transaction)
				can then be reordered directly before or after $\tiiat{\mu}$.
		\end{itemize}
	\item {\bf Pure Happens-Before.}
		WLOG fix $\tiiat{\mu} \prec \tjjat{\nu} \in I$.
		Then to race in $I$ there must be no pair of synchronizing instructions
		$\tiiat{\epsilon}$ (a release edge) and $\tjjat{\chi}$ (an acquire edge) such that
		\[
			\tiiat{\mu} \prec \tiiat{\epsilon} \prec \tjjat{\chi} \prec \tjjat{\nu} \in I
		\]
		to update the vector clock epoch between \tiiat{\mu} and \tjjat{\nu} \cite{djit,fasttrack}.
		\begin{itemize}
			\llitem $I \Rightarrow I'$:
				If $\tiiat{\mu},\tjjat{\nu}$ race in $I$,
				then they cannot both be in successful transactions,
				or else \Cref{lem:equiv} normalization would provide
				the corresponding HTM end and begin for $\tiiat{\epsilon}$ and $\tjjat{\chi}$ respectively.
				Hence there will be no unlock/lock pair on $m$ in $I'$ to satisfy the above sequence.
			\llitem $I' \Rightarrow I$:
				If $\tiiat{\mu},\tjjat{\nu}$ race in $I'$,
				then they cannot both hold $m$,
				or else $\mathsf{lockify}_m$ would provide the corresponding
				unlock and lock for $\tiiat{\epsilon}$ and $\tjjat{\chi}$ respectively.
				Hence there will be no HTM end/begin pair in $I$ to satisfy the above sequence.
		\end{itemize}
		%In both cases $I$ and $I'$ otherwise have the same synchronizing instructions.
\end{itemize}
Hence, data-race analysis is sound when transaction boundaries are replaced by an abstract global lock.
\end{proof}

\subsection{Weak atomicity}
\label{sec:tm-weak-atomicity}

\revision{Prior work has highlighted the difference between
{\em strong} and {\em weak atomicity}
of transactions
\cite{htm-subtleties,stm-weak-atomicity}.
Thus far we have assumed transactions will abort from any intervening memory conflict,
whether or not that access is itself transactional.
This corresponds to strong atomicity, the execution semantics provided by TSX.
% \cite{bugtm}. % sec 3.1 of that paper, not sure of an earlier citation that also states it in these exact terms
However, weakly-atomic STM systems
do not abort transactions when
a non-transactional access conflicts,
and condition \#1 of the execution semantics from \cref{sec:tm-design-failure} does not hold.
To model weak atomicity, \Cref{lem:equiv} must be extended with a third possibility:
%as follows (updated parts underlined for emphasis).
\vspace{1em}

\begin{lemma}[Weak Equivalence of Aborts]
	\label{lem:equiv-weak}
	For any interleaving prefix such as defined in \Cref{lem:equiv}, either one of the original two conditions hold, or:
	\begin{enumerate}
		\setcounter{enumi}{2}
		\item All events in
			$\tjjat{\gamma_1}\dots\tjjat{\gamma_j},
			 \tkkat{\kappa_1}\dots\tkkat{\kappa_k}$
			 % ehh, bad parallel structure with below regarding the verb 'exists'
			are non-transactional, and the original interleaving itself is a legal execution.
	\end{enumerate}
\end{lemma}

\begin{proof}
	As before, except first case on
	whether any intervening operations by $\tjj$/$\tkk$ are transactional.%
	\footnote{If $\tjj$ or $\tkk$ already have a transaction active,
	their {\sf begin} must have occurred before $\tiiat{\alpha}$,
	and the theorem can apply starting from the earlier {\tt begin} with the thread IDs swapped.}
	If so, either the original case \#1 or \#2 applies.
	If not, $\tii$ may ultimately proceed with $\tiiat{\beta_{b+1}}$
	(its preempted next transactional instruction)
	without aborting,
	and the new case \#3 applies instead.
\end{proof}

\Cref{thm:atom} must be weakened correspondingly.
\vspace{1em}

\begin{theorem}[Weak Atomicity of Transactions]
	\label{thm:atom-weak}
	For any state space $S$ of a transactionally-concurrent program,
	an equivalent state space exists in which all transactions are either executed
	{\bf atomically
	with respect to other transactions only},
	or aborted immediately.
\end{theorem}

\begin{proof}
	As before.
\end{proof}

Conveniently,
\Cref{thm:atom-weak} still guarantees transaction-transaction atomicity,
so mutual exclusion of $m$ is preserved
in \cref{sec:tm-design-locks}'s formulation of
$I' = \exists m. \mathsf{lockify}_m(I)$,
%(a global lock does not protect against interleaving by arbitrary other code that never takes the lock either)
and
%the proof of
\Cref{thm:global-lock} applies unmodified.
\cref{sec:tm-impl-weak-atomicity}
discusses the implementation consequences of this weakening.}

%%%%%%%%%%%%%%%%%%%%%%%%%%%%%%%%%%%%%%%%%%%%%%%%%%%%%%%%%%%%%%%%%%%%%%%%%%%%%%%%

\section{Implementation}
\label{sec:tm-implementation}

Support for TSX programs in Landslide is implemented in \revisionminor{five} parts, broadly speaking:
the user interface,
plus one internal part corresponding to each proof above,
plus a bonus optimization for pruning equivalent interleavings under the new concurrency model.

\subsection{User interface}
\label{sec:tm-interface}

Landslide provides its own ``implementation'' of the TSX interface,
which matches GCC's interface exactly,
located in {\tt 410user/inc/htm.h}
under {\tt pebsim/p2-basecode/}.
The interface is implemented in {\tt 410user/libhtm/htm.c},
perhaps surprisingly, as totally empty functions;
Landslide hooks these functions' addresses during instrumentation
and inserts preemption points, failure injections,
et cetera as necessary whenever the execution encounters them.

Transactional test programs should be ported to the Pebbles userspace if not already,
then any use of compiler HTM intrinsics replaced with Landslide's interface.%
\footnote{Attempts to execute a real TSX instruction under Landslide, instead of using the custom interface,
will be reported as ``invalid opcode'' bugs,
as neither of its supported simulation platforms support the feature.
}
They should then be put in either {\tt 410user/progs/} or {\tt user/progs/}
and the {\tt 410TESTS} or {\tt STUDENTTESTS} line (respectively) of {\tt config-incomplete.mk} be edited to add the test name,
before running {\tt p2-setup.sh} (\cref{sec:landslide-setup})
% TODO(implementation): Provide an obfuscated refp2?
on a (hopefully) correct P2 as usual.
Several example HTM tests are provided as {\tt 410user/progs/htm*.c}.

Finally, Quicksand supports the following command-line options to enable various sets of HTM features within Landslide.
%Each option requires each of the previously-listed options to also be present.
\begin{itemize}
	\item {\tt -X} (for ``tsX'' or ``Xbegin'') enables the basic features:
		preemption points on each {\tt \_xbegin()} and {\tt \_xend()} call,
		failure injection on each of the former
		(always returning {\tt \_XABORT\_RETRY} as the failure code)
		(\revisionminor{\cref{sec:tm-design-failure}}, \cref{sec:txn-failure}),
		and treatment of transactional regions during data-race analysis (\cref{sec:txn-datarace}).
	\item {\tt -A} (for ``xAbort codes'')
		enables multiple xabort failure codes (\cref{sec:txn-failure}).
		Requires {\tt -X}.
	\item {\tt -S} (for ``Stm'' or ``Suppress retries'')
		disables the {\tt \_XABORT\_RETRY} failure reason,
		causing Landslide to emulate the semantics of STM rather than HTM.
		%(This feature was used in \cite{sigbovik-htm}~\sect{5}.)
		Requires {\tt -X -A}.
	\item {\tt -R} (for ``Retry sets'')
		enables state space reduction based on independences between individual {\tt xbegin} results
		(\cref{sec:tm-retrysets}).
		Requires {\tt -X} and {\em not} {\tt -A}.
	\item \revision{{\tt -W} (for ``Weak atomicity'')
		enables weak atomicity, %for transactions,
		allowing non-transactional code to interleave within transactions,
		corresponding to \cref{sec:tm-weak-atomicity}.
		Requires {\tt -S}.}
\end{itemize}

\subsection{Failure injection}
\label{sec:txn-failure}

Each preemption point ({\tt struct hax} as defined in {\tt tree.h})
includes three new fields.
The boolean {\tt h->xbegin} is set if the preemption point occurred at an {\tt \_xbegin()} call.
Then if set (as the tag of an option type),
the twin lists of integers
{\tt h->xabort\_codes\_ever}
and
{\tt h->xabort\_codes\_todo}
store possible error codes this {\tt \_xbegin()} call should be tested to possibly return.
The former list stores all such error codes, whether already tested or yet to be tested,
while the latter serves as a workqueue that indicates only those not already tested yet
(serving an analogous purpose as the {\tt all\_explored} flag for thread scheduling).
The state space estimators (\cref{sec:landslide-estimate})
check the length of these lists, in addition to DPOR's tagged threads,
when computing the number of marked children.

\subsubsection{Adding possible abort codes}
\label{sec:txn-abort-codes}

Up to four abort codes are considered depending on \cref{sec:tm-interface}'s testing options.
All abort codes are simultaneously added to both lists,
unless already present on {\tt xabort\_codes\_ever},
in which case not added to {\tt xabort\_codes\_todo} to avoid duplicate work.
\begin{itemize}
	\item {\tt \_XABORT\_RETRY}:
		When a {\tt xbegin} preemption point is created ({\tt save\_setjmp()}),
		if {\tt -S} is not set,
		both {\tt xabort\_codes} lists are initialized to contain this code.
		This represents the possibility for a hardware transaction to fail
		for reasons outside of the programmer's control, such as system interrupts or cache eviction.
		If {\tt -S} is set they are initialized to empty,
		representing either STM's policy of failing only when a true memory conflict arises,
		or a TSX user wrapping all her {\tt \_xbegin()}s in a retry loop such as in \cite{sigbovik-htm}.
	\item {\tt \_XABORT\_CONFLICT}:
		Whenever DPOR detects a memory conflict between two transitions ({\tt shimsham\_shm()}),
		if {\tt -A} is set,
		% for each of the two which are within a transaction, % no!!!!!!!
		%it finds the nearest ancestor {\tt xbegin} preemption point by the same thread
		if the later-executed transition is a transactional success path,
		it adds this code to both lists of the immediately preceding preemption point.
		Note that a transaction should suffer a conflict abort
		only when executed {\em after} a conflicting memory access
		to avoid circular causality
		(see {\tt 410user/progs/htm\_causality.c} for rationale).
		Should the transaction happen to be executed first,
		DPOR will first try to reorder it as normal,
		and then abort it as described if the conflict persists.
	\item {\tt \_XABORT\_EXPLICIT}:
		Whenever the program invokes {\tt \_xabort()} ({\tt sched\_update\_user\_\allowbreak{}state\_machine()}),
		if {\tt -A} is set,
		this code is added to both lists, bitwise-ored with the user-supplied argument code
		as specified in \cite{htm-gcc},
		and execution of the transactional path immediately stops.
	\item {\tt \_XABORT\_CAPACITY}:
		Whenever the program invokes a system call during a transaction,
		if {\tt -A} is set,
		this code is added to both lists,
		and execution of the transactional path immediately stops.
		\revision{So called because doing so would trigger a mode switch,
		which the specification allows to abort for any reason.%
		\footnote{\revision{As far as I know, the reason for syscall aborts is both hardware- and kernel-dependent,
		and in fact nondeterministic;
		on my Haswell Core i7 running Linux 4.19, I observed a majority of {\tt \_XABORT\_CONFLICT}s,
		a scant few {\tt \_XABORT\_CAPACITY}s,
		and surprisingly many with no reason specified whatever;
		however, not once did a transaction succeed with even a single {\tt gettid()} call
		\cite{tsx-syscall}.
		For simplicity of both implementation and client programs,
		Landslide always distinguishes system call aborts with {\tt \_XABORT\_CAPACITY}.}}
		See below for discussion of other cache capacity concerns.}
\end{itemize}

{\bf Limitations.}
Landslide does not yet check for false sharing,
i.e. read/write or write/\allowbreak{}write access pairs to different memory addresses that share a cache line.
On real hardware, these would produce {\tt \_XABORT\_CONFLICT} failures,
but to find such access pairs would require extending {\tt mem\_shm\_intersect()}'s set intersection algorithm
to consider a certain degree of ($N$-byte-aligned) fuzziness when comparing addresses,
as well as adding a command-line option to configure said $N$, the cache-line size to simulate.
For now, Landslide (wrongly) lumps false sharing in among the ``spurious'' {\tt \_XABORT\_RETRY} failure reasons.
%
On HTM, even if the user wraps all such spurious failure reasons in a retry loop,
false sharing (i.e., disjoint memory accesses that share a cache line)
should still produce a non-retryable abort,
begetting a discrepancy with STM,
which aborts only when the memory addresses match exactly.
Accordingly, the {\tt -S} option described above provides STM semantics as currently implemented.
Future work could extend Landslide with an option to configure false sharing conflicts
to remain faithful to HTM semantics and still abort even under {\tt -S}.
%
Likewise, Landslide does not check for a transaction's memory footprint exceeding the CPU's cache capacity,
which on real hardware would trigger a {\tt \_XABORT\_CAPACITY} abort.
Landslide's memory tracking could simulate this check as well,
perhaps with a configurable cache size,
but would likely be theoretically uninteresting
(all programs in the upcoming evaluation have trivial memory usage),
and so is left unimplemented for now.
%
Finally,
{\tt \_XABORT\_NESTED},
the last abort code specified by \cite{htm-gcc},
depends on currently-unsupported {\tt xbegin} nesting, which I discuss further in
\cref{sec:tm-warpzone-nested}.

\subsubsection{Injecting abort codes}
When traversing the state space (\cref{sec:landslide-statespace}),
in addition to performing DPOR to select non-independent thread interleavings (\cref{sec:landslide-dpor}),
the abort codes under each {\tt xbegin} preemption point are also considered ``marked'' paths which must be tested.
Hence {\tt explore()}, by way of {\tt any\_tagged\_child()},
will pop off the {\tt xabort\_codes\_todo} queue
when it's time to explore that preemption point in the usual depth-first manner.%
\footnote{The search order prioritizes abort codes before scheduling other threads
at such preemption points,
which is just an implementation detail, not theoretically necessary.}
The optional abort code is then passed through {\tt arbiter\_append\_choice()}/{\tt arbiter\_pop\_choice()}
to {\tt cause\_transaction\_failure()},
which edits the simulation state ({\tt \%eip} and {\tt \%eax})
to force {\tt \_xbegin()} to return the provided code.

\subsection{Data race analysis}
\label{sec:txn-datarace}

When a thread returns {\tt \_XBEGIN\_STARTED} from {\tt \_xbegin()}
(analogous to {\tt mutex\_trylock()}),
Landslide's scheduler sets the {\tt user\_txn} action flag for that thread (\cref{sec:landslide-scheduler}),
and if using Pure Happens-Before,
applies \textsc{FT acquire} (\cref{sec:landslide-phb}) using a dummy lock address to represent the abstract global lock.
When a thread reaches {\tt \_xend()},
the flag is cleared,
and under Pure Happens-Before,
\textsc{FT release} is applied.
Then when {\tt check\_locksets()} compares an access pair,
under Limited Happens-Before,
it is considered a data race only if at least one thread's {\tt user\_txn} was not set
in addition to the usual conditions; % in multiplication to..? && c.c
under Pure Happens-Before,
the vector clocks are simply checked as usual.

\subsection{Weak atomicity}
\label{sec:tm-impl-weak-atomicity}

\revision{If the user supplies the {\tt -W} option (\cref{sec:tm-interface}),
Landslide will emulate STM's weak atomicity semantics (\cref{sec:tm-weak-atomicity}).
%to allow non-transactional code
%to interleave between a {\tt \_xbegin()}/{\tt \_xend()} pair.
Concretely, it achieves this by
reporting both halves of data races that occur between transactional and non-transactional code (\cref{sec:landslide-datarace})
(previously, only the latter half would be reported),
allowing preemption points to be identified during transactions (\cref{sec:landslide-pps}),
and allowing system calls during transactions to proceed rather than forcing a {\tt CAPACITY} abort
(as programmers may now need to nest other synchronization within transactions to protect against non-transactional code).
%and traversing the interleaving arbitrarily far back from a {\tt \_xend()} to find the appropriate preemption
%(previously, the xbegin would be at the immediately previous PP)
Landslide requires the other STM options ({\tt -A -S}) also to be enabled to emulate weak atomicity,
partly because I know of no existing HTM system which implements it,
but mostly to keep the implementation simple.
As one might expect, it inflates state spaces considerably, % TODO: measure it with htm2 or some such?
and moreover,
every single HTM benchmark in the upcoming evaluation,
implemented to rely on strong atomicity,
exhibits bugs under weak atomicity.
Commit 68401fc implements this mode.}

\subsection{Retry independence}
\label{sec:tm-retrysets}

Finally, I identified a specific pattern of transactional code where existing state-space reduction algorithms
will fail to identify and prune equivalent thread interleavings.
\Cref{fig:retry-set-example} shows a minimal example program which exhibits this problem.
In this program, each thread's transactional path conflicts with the other thread's abort path,
while the two transactional paths (disjoint memory accesses) and the two abort paths (reads only) are
independent with each other.

\begin{figure}[h]
        \begin{center}
        \begin{tabular}{ll}
		\multicolumn{2}{c}{Initially {\tt \ctype{int} foo = \const{0}, bar = \const{0};}} \\
                \\
                \begin{tabular}{l}
                {\bf Thread 1} \\
                \hline
			\texttt{\flow{if} (\call{\_xbegin}() == \const{STARTED}) \{} \\
			\texttt{~~~~foo++;} \\
			\texttt{~~~~\call{\_xend}();} \\
			\texttt{\} \flow{else} \{} \\
			%\texttt{~~~~\call{mutex\_lock}(\&m);} \\
			\texttt{~~~~\flow{assert}{}(foo + bar < \const{2});} \\
			%\texttt{~~~~\call{mutex\_unlock}(\&m);} \\
			\texttt{\}} \\
                \end{tabular}
                &
                \begin{tabular}{l}
                {\bf Thread 2} \\
                \hline
			\texttt{\flow{if} (\call{\_xbegin}() == \const{STARTED}) \{} \\
			\texttt{~~~~bar++;} \\
			\texttt{~~~~\call{\_xend}();} \\
			\texttt{\} \flow{else} \{} \\
			%\texttt{~~~~\call{mutex\_lock}(\&m);} \\
			\texttt{~~~~\flow{assert}{}(foo + bar < \const{2});} \\
			%\texttt{~~~~\call{mutex\_unlock}(\&m);} \\
			\texttt{\}} \\
                \end{tabular}
        \end{tabular}
        \end{center}
	\caption{Motivating example for retry independence reduction.}
	\label{fig:retry-set-example}
\end{figure}

Ordinarily, in the state space subset which schedules thread 1 before thread 2,
there would be 4 combinations of each thread succeeding or aborting their respective transactions;
among those, at least one would show a memory conflict,
causing DPOR to explore into the other half of the state space which schedules thread 2 first,
where the same 4 combinations of transaction reuslts would be tested in the other order,
even though 2 are redundant under schedule reordering.
Consequently, na\"ively testing both success and retry aborts in both threads regardless of reordering
will unnecessarily execute % oops
both equivalent interleavings from each such pair;
to identify such equivalences,
DPOR must somehow remember which combination of {\tt xbegin} results led to the memory conflict in the first place.

While \Cref{fig:retry-set-example}'s example may seem contrived
(what program's transactions would just give up and do no work if aborted?),
it is easy to imagine a larger transactional data structure,
two insertions into which might operate on disjoint nodes or array indices,
allowing simultaneous transactions to usually succeed,
while the more uncommon abort paths might take the opportunity to assert a full consistency check of all elements,
ultimately resulting in a similar conflict pattern.
Also, the evaluation will later show
(\cref{sec:tm-eval-retrysets})
that even programs with fully conflicting success/abort paths
may still exhibit some equivalent thread interleavings of this pattern
after their failure paths are split apart into smaller transitions by data-race preemption points.

To address this,
I extended Landslide
with {\em retry set reduction} (commit 86657c7),
an experimental feature named after prior work's analogous {\em sleep sets}
(\cite{partial-order-methods,dpor,optimal-dpor}, \cref{sec:landslide-sleepsets}).
Whenever DPOR tags a new branch for exploration (\cref{sec:landslide-explore}),
if either or both sides of a memory conflict were part of a transaction,
it records a {\em retry set},
i.e., the pair of {\tt xbegin} results executed by the conflicting threads,
to accompany that branch.
Like the {\tt xabort\_codes} lists, the state space estimators (\cref{sec:landslide-estimate})
check the number of retry sets when counting the number of marked children.
% ah-- not compatible after all!
%in addition to tracking the set of abort codes possible at each preemption point.
Later, when traversing the new branch,
Landslide remembers those threads' expected {\tt xbegin} results
and refuses to test any others
(unless DPOR separately found them, too, to conflict),
thereby skipping over (i.e., pruning) reorderings of any other {\tt xbegin} results that would be independent.
Like sleep sets, it also considers the preempted thread ``retry-set blocked''
(like sleep-set blocking \cite{optimal-dpor}),
and refuses to run it until the conflicting thread runs its transaction first,
unless such would result in deadlock.%
\footnote{
	Before the search ordering update to Landslide's normal sleep sets implementation
	(\cref{sec:landslide-sleepsets}),
	I observed false-positive retry-set-blocked deadlocks fairly often;
	after the update,
	it took until 981773 interleavings (>9 hours) into {\tt htm2}(3,2) to find one
	and confirm the need to still explicitly avoid them by abandoning the retry set in such cases.
	\revisionminor{Commit 02b70c8 implements this fix.}
}
%The {\tt landslide} repository's {\tt retry-sets} branch implements this feature,
Upcoming in the evaluation, \Cref{fig:retry-reduction} will visualize the reduction achieved in two test cases.

% In contrast to Landslide's implementation of the sleep set reduction (\cref{sec:landslide-explore}),
% blah blah which is implicit and works in retrospect by pruning at tag-sibling time, which figures out what would have been in a sleep set by analyzing the past independences and explored siblings,
% this cannot work that way,
% because in the first place, the former works with the "keep running current thread" ordering heuristic,
% (which automatically takes care of the first half of the sleep set problem)
% whereas no such similar heuristic exists (can exist??) for xbegin PPs, we always just run the success path first

%%%%%%%%%%%%%%%%%%%%%%%%%%%%%%%%%%%%%%%%%%%%%%%%%%%%%%%%%%%%%%%%%%%%%%%%%%%%%%%%

\section{Evaluation}
\label{sec:tm-eval}

While prior work has focused on verifying implementations of transactional memory themselves
\cite{specifying-verifying-tm,tm-correctness,tm-completeness,mc-tm-with-spin},
Landslide is to the best of my knowledge the first model checker to support transactional client code.
% mumble smth new concurrency model
Accordingly, there is no baseline against which to compare Landslide's performance.
Likewise, since Landslide's HTM semantics emulation relies on the equivalences proved in \cref{sec:tm-design},
I did not actually implement an HTM-style speculative-execution-and-rollback simulation mode.
% so have no basis to measure how much it saves, or smth like that
On this count, at least,
I hope the reader finds it self-evident that
the equivalence proofs provide exponential state space reduction
compared to actually testing (and thereafter aborting) every combination of instructions within transactions.
Beyond those, this chapter's evaluation will take a relatively green-field and exploratory approach.
I pose the following evaluation questions.

% on one hand,
% comparing with \cite{mc-tm-with-spin}, succeeding for (2,3) then burning out with (2,4) is consistent with prior work
% on the other hand, they were verifying STM implementation, which is no doubt more operations per iteration?
% whatever, don't need to make this argument

\begin{enumerate}
	\item How quickly does Landslide find bugs in incorrect transactional programs? %of varying sizes?
	\item Does Landslide find any previously-unknown bugs in real-world transactional code?%
		\footnote{The savvy reader will realize that whether or not the author poses this evaluation question to begin with spoils its answer.}
	%\item How quickly does Landslide verify correct transactional programs of varying sizes?
	\item How does Iterative Deepening's (\cref{sec:quicksand-id})
		performance compare to Maximal State Space mode (\cref{sec:landslide-quicksand-options})?%
		\footnote{Not necessarily related to HTM,
		but the latter was implemented well after \Cref{chap:quicksand}'s conference paper was published,
		so this was the most convenient test suite to evaluate it on.}
	\item How well does Landslide's verification scale
		with increasing thread/iteration count for correct transactional programs?
	\item What further reduction can be achieved beyond the baseline provided by
		the global-lock/failure-injection equivalences?
	%\item By the way, should MC research papers quantify variance in their CPU-time performance experiments?
\end{enumerate}

\subsubsection{Experimental setup}
\label{sec:tm-eval-exp-setup}

The evaluation suite comprises several unit tests hand-written by yours truly,
microbenchmarks and transactional data structures from \cite{tm-benchmark-cmu},
a transactional spinlock from \cite{spinlock-rtm-github},
and various combinations thereof,
as follows.

% notes to self abour marios tests
% counter: tests k threads on n iterations, benchmark xchg-aways vs txn. removed benchmarking, added assertion that no counts are lost.

\begin{enumerate}
	\item Unit tests and microbenchmarks
	\begin{itemize}
		\item {\tt htm1}: The bug from \Cref{fig:htm-example,fig:htm-example-reproduced}.
			%(assertion failure on {\tt count}).
		\item {\tt htm2}: The fixed version as in \Cref{fig:htm-fixed}.
		\item {\tt counter}: Microbenchmark version of {\tt htm2} which replaces the complex locking failure path with an atomic {\tt xadd}, from \cite{tm-benchmark-cmu}.
		\item {\tt swap}: Microbenchmark that swaps values in an array, from \cite{tm-benchmark-cmu}.
		\item {\tt swapbug}: {\tt swap} modified to introduce circular locking in the failure path. %(deadlock).
		\item {\tt fig63}: Generalized version of \Cref{fig:retry-set-example},
			contrived to induce as much reduction as possible from the retry sets optimization.
	\end{itemize}
	\item Data structure tests
	\begin{itemize}
		\item {\tt avl\_insert}: AVL tree concurrent insertion test \cite{tm-benchmark-cmu}.
		\item {\tt avl\_fixed}: {\tt avl\_insert} with the AVL bug fixed (spoilers!!).
		\item {\tt map\_basic}: Separate-chaining hashmap concurrent insertion test \cite{tm-benchmark-cmu}.
		\item {\tt map\_basicer}: {\tt map\_basic} modified with a larger initial size to skip the resizing step.
		\item {\tt avl\_mutex}: {\tt avl\_fixed} with transactional sections simplified by abstraction into a mutex.
		\item {\tt map\_mutex}: {\tt map\_basic} simplified likewise.
		\item {\tt map\_basicer\_mutex}: {\tt map\_basicer} simplified likewise.
	\end{itemize}
	\item Lock abstraction tests
	\begin{itemize}
		\item {\tt lock}(): Checks that multiple threads using a transactional lock
			cannot access the critical section simultaneously.
			% htm_spinlock and htm_mutex
		\item {\tt lock\_fast}(): Checks that a transactional lock's fast path will not suffer conflict aborts
			if its client threads' critical sections are independent.
			% htm_spinlock_fastpath and htm_mutex_fastpath
	\end{itemize}
		These are each parameterized over implementations {\tt spinlock} (from \cite{spinlock-rtm-github}),
		{\tt spin\_\allowbreak{}fixed} (spoilers!!),
		and {\tt mutex} (replaces the spinlock with a Landslide-annotated P2 mutex to reduce state space size).
\end{enumerate}

The notation {\tt testname}$(K,N)$ %will denote
denotes a test configuration of $K$ threads, each running $N$ iterations of the test logic.
All tests were run on an 8-core 2.7GHz Core i7 with 32 GB RAM.
Reported CPU-times include time spent on all state spaces
Quicksand saw fit to run,
not just the maximal or the buggy state space;
for verification tests (run with {\tt -M}),
this still includes abandoned smaller jobs that were run to saturate the set of data-race preemption points
(\Cref{chap:quicksand}).
To minimize variance in CPU-time measurements,
I ensured the test machine was not loaded beyond normal web browser use,
and ran only one instance of Landslide at a time;
for further discussion of variance see \cite{sigbovik-htm}.
The number of interleavings in each state space \revisionminor{is}, of course, deterministic and do not vary across runs.

I investigated the popular transactional benchmark suite STAMP \cite{stamp}
to include in this test suite,
but found that all transactional code therein was written without failure paths,
so would likely not contribute any theoretical depth to the evaluation.
STAMP uses the OpenTM interface \cite{opentm},
which allows the programmer to specify transactional code regions to be implemented atomically,
however the underlying architecture may require
(whether HTM or STM; presumably with retry loops on HTM).
On one hand,
this lends credence to my upcoming conclusion
that HTM programs should be written at a higher abstraction level than calling {\tt \_xbegin()} directly;
on the other,
OpenTM's requirement of {\em virtualized transaction} semantics
(i.e., being unconstrained by memory footprint, able to make system calls, and able to nest arbitrarily)
is more suited to STM and glosses over many pitfalls of HTM programming.
%% #overlyhonestmethods
% I also wanted to test the Eunomia algorithm for transactional search trees \cite{eunomia},
% but its implementation is not open source,
% and the authors ignored my email requesting to use it.

Finally, the keen-eyed reader will notice the state space sizes reported here
differ from those reported in \cite{sigbovik-htm}.
All experiments have been re-run
on account of three updates to Landslide's exploration algorithm
implemented since then:
the sleep sets optimization for DPOR (\cref{sec:landslide-sleepsets}, commits 0447666 and 588687c),
the {\tt thrlib\_function} and {\tt TRUSTED\_THR\_JOIN} directives
to mark internal thread library logic as trusted
(\cref{sec:landslide-config-landslide}, commits 64a02e4 and a50d4ea),
and fixing a soundness bug in which Landslide could neglect to inject transaction failures
immediately after a thread switch (commit dcae85b).
On account of the former two updates, some state spaces may be smaller than before;
on account of the third, some may be larger.
These updates do not discredit the bug-finding results (a bug is a bug),
but the previously-published verification results should be considered outdated.

\newcommand\ETA[1]{\hilight{brownish}{{\em #1}}\xspace}
\newcommand\cpu[1]{\hilight{darkcyan}{{#1}}\xspace}
\newcommand\wtm[1]{\hilight{lime}{{#1}}\xspace}
\newcommand\ints[1]{\hilight{pinkish}{{#1}}\xspace}
\newcommand\ETAdag[1]{\ETA{\ensuremath{\dagger}#1}}

\subsection{Bug-finding}
\label{sec:tm-eval-bugs}

\begin{table*}[t]
	\begin{center}
		\footnotesize
	\begin{tabular}{cc||r|r|r||r|r|r|r}
			&	&\multicolumn{3}{c||}{\bf Quicksand mode}&\multicolumn{4}{c}{{\bf Maximal state space mode} ({\tt -M})} \\
		\bf buggy test	& \bf K,N &\cpu{\bf cpu (s)}&\wtm{\bf wall (s)}&\ints{\bf int's}&\cpu{\bf cpu (s)}&\wtm{\bf wall (s)}&\ints{\bf int's}& \ETA{\bf SS size (est.)} \\
		\hline
		\hline
		{\tt htm1}
			& 2,1	&\cpu{  32.73}&\wtm{8.02}&\ints{5}& \cpu{*9.24}& \wtm{*6.49}& \ints{5}& \ETA{12} \\
			% 5 interleavings tested; 1 preemptions; job time 3s; pldi time 8015847; new-fixed pldi cputime 32726660
			% 5 interleavings tested; 1 preemptions; job time 2s; pldi time 6492339; new-fixed pldi cputime 9242577 -M
		(assertion)
			& 2,2	&\cpu{  50.72}&\wtm{9.22}&\ints{9}& \cpu{*9.93}& \wtm{*7.03}& \ints{9}& \ETA{102} \\
			% 9 interleavings tested; 1 preemptions; job time 5s; pldi time 9220760; new-fixed pldi cputime 50719368
			% 9 interleavings tested; 1 preemptions; job time 3s; pldi time 7028111; new-fixed pldi cputime 9933516 -M
			& 2,3	&\cpu{ 88.35}&\wtm{14.63}&\ints{17}& \cpu{*10.03}& \wtm{*7.12}& \ints{17}& \ETA{819} \\
			% 17 interleavings tested; 1 preemptions; job time 7s; pldi time 14632911; new-fixed pldi cputime 88346856
			% 17 interleavings tested; 1 preemptions; job time 3s; pldi time 7119866; new-fixed pldi cputime 10025347 -M
			& 2,4	&\cpu{108.31}&\wtm{17.44}&\ints{33}& \cpu{*11.24}& \wtm{*8.29}& \ints{33}& \ETA{6553} \\
			% 33 interleavings tested; 1 preemptions; job time 9s; pldi time 17442317; new-fixed pldi cputime 108309778
			% 33 interleavings tested; 1 preemptions; job time 3s; pldi time 8294161; new-fixed pldi cputime 11235910 -M
		\cline{2-9}
			& 3,1	&\cpu{ 60.63}&\wtm{11.63}&\ints{5}& \cpu{*9.63}& \wtm{*6.70}& \ints{5}& \ETA{76} \\
			% 5 interleavings tested; 1 preemptions; job time 5s; pldi time 11631147; new-fixed pldi cputime 60628270
			% 5 interleavings tested; 1 preemptions; job time 2s; pldi time 6695336; new-fixed pldi cputime 9627076 -M
			& 3,2	&\cpu{ 78.69}&\wtm{13.17}&\ints{9}& \cpu{*9.43}& \wtm{*6.37}& \ints{9}& \ETA{3686} \\
			% 9 interleavings tested; 1 preemptions; job time 6s; pldi time 13173912; new-fixed pldi cputime 78690603
			% 9 interleavings tested; 1 preemptions; job time 2s; pldi time 6370046; new-fixed pldi cputime 9427268 -M
			& 3,3	&\cpu{ 70.50}&\wtm{13.05}&\ints{17}& \cpu{*10.43}& \wtm{*6.96}& \ints{17}& \ETA{176947} \\
			% 17 interleavings tested; 1 preemptions; job time 6s; pldi time 13045016; new-fixed pldi cputime 70504465
			% 17 interleavings tested; 1 preemptions; job time 3s; pldi time 6962930; new-fixed pldi cputime 10427482 -M
			& 3,4	&\cpu{ 70.69}&\wtm{12.96}&\ints{33}& \cpu{*11.74}& \wtm{*8.71}& \ints{33}& \ETA{8493465} \\
			% 33 interleavings tested; 1 preemptions; job time 7s; pldi time 12964843; new-fixed pldi cputime 70692866
			% 33 interleavings tested; 1 preemptions; job time 4s; pldi time 8710265; new-fixed pldi cputime 11743862 -M
		\cline{2-9}
			& 4,1	&\cpu{ 51.83}&\wtm{9.85}&\ints{5}& \cpu{*9.56}& \wtm{*6.63}& \ints{5}& \ETA{460} \\
			% 5 interleavings tested; 1 preemptions; job time 5s; pldi time 9851604; new-fixed pldi cputime 51833724
			% 5 interleavings tested; 1 preemptions; job time 2s; pldi time 6627275; new-fixed pldi cputime 9560320 -M
			& 4,2	&\cpu{ 44.83}&\wtm{8.94}&\ints{9}& \cpu{*9.83}& \wtm{*6.84}& \ints{9}& \ETA{132710} \\
			% 9 interleavings tested; 1 preemptions; job time 4s; pldi time 8936964; new-fixed pldi cputime 44826654
			% 9 interleavings tested; 1 preemptions; job time 2s; pldi time 6835025; new-fixed pldi cputime 9830626 -M
		\hline
		{\tt swapbug}
			& 2,1	&\cpu{*26.25}&\wtm{*6.42}&\ints{*6}& \cpu{47.80}& \wtm{13.43}& \ints{33}	& \ETA{73} \\
			% using fewer interleavings one for quicksand out of 2 close in cputime bugs
			% 6 interleavings tested; 1 preemptions; job time 4s; pldi time 6415943; new-fixed pldi cputime 26251146
			% 13 interleavings tested; 1 preemptions; job time 3s; pldi time 4654607; new-fixed pldi cputime 16966236
			% 33 interleavings tested; 1 preemptions; job time 4s; pldi time 13430843; new-fixed pldi cputime 47803093
		(deadlock)
			& 2,2	&\cpu{*18.08}&\wtm{*4.98}&\ints{*10}& \cpu{51.37}& \wtm{16.78}& \ints{85}	& \ETA{860} \\
			% 10 interleavings tested; 1 preemptions; job time 3s; pldi time 4980992; new-fixed pldi cputime 18078134
			% 85 interleavings tested; 1 preemptions; job time 7s; pldi time 16777253; new-fixed pldi cputime 51372309
			& 2,3	&\cpu{*20.93}&\wtm{*5.57}&\ints{*18}& \cpu{57.87}& \wtm{23.99}& \ints{217}	& \ETA{9120} \\
			% 18 interleavings tested; 1 preemptions; job time 4s; pldi time 5573337; new-fixed pldi cputime 20932814
			% 217 interleavings tested; 1 preemptions; job time 15s; pldi time 23987230; new-fixed pldi cputime 57868535
			& 2,4	&\cpu{*38.92}&\wtm{*8.61}&\ints{*34}& \cpu{82.95}& \wtm{48.31}& \ints{537}	& \ETA{91239} \\
			% 34 interleavings tested; 1 preemptions; job time 6s; pldi time 8614371; new-fixed pldi cputime 38917975
			% 537 interleavings tested; 1 preemptions; job time 38s; pldi time 48314592; new-fixed pldi cputime 82952352
		\cline{2-9}
			& 3,1	&\cpu{*38.59}&\wtm{*8.50}&\ints{*32}&\cpu{88.28}& \wtm{72.41}& \ints{1016}	& \ETA{3543} \\
			% 32 interleavings; 258 preemptions; job time 5s; pldi time 8506344; new-fixed pldi cputime 38594745
			% 1016 interleavings; 258 preemptions; job time 58s; pldi time 72406995; new-fixed pldi cputime 88281009
			& 3,2	&\cpu{*1572.83}&\wtm{*199.98}&\ints{*262}&\ETA{--}&\ETA{>1h}&\ETA{--}& \ETA{1683509} \\
			% 262 interleavings; 260 preemptions; job time 47s; pldi time 199982231; new-fixed pldi cputime 1572825801
			% [JOB 8] progress: 47047/1683509 brs (2.794578%), ETA 1y 237d 20h 41m 57s (elapsed 59m 44s)
		\hline
		{\tt avl\_insert}
			% note the somewhat chaotic behaviour between quicksand 2,2 and 2,3; the latter surprisingly FABs faster
			% this is likely because non-drpp SSes were not big enough to defer for longer in 2,2,
			% whereas in 2,3 it went straight to the buggy SS asap
			& 2,2	&\cpu{2494.77}&\wtm{315.46}&\ints{*29}	&\cpu{*95.53}&\wtm{*30.56}&\ints{79}&\ETA{158505} \\
			% 79 interleavings; 1 preemptions; job time 9s; pldi time 30558236; new-fixed pldi cputime 95530430
			% 29 interleavings; 1 preemptions; job time 14s; pldi time 315461192; new-fixed pldi cputime 2494772471
		(segfault+)
			& 2,3	&\cpu{308.10}&\wtm{*43.95}&\ints{*33}&\cpu{*249.42}&\wtm{144.33}&\ints{835}&\ETA{13664203} \\
			% 835 interleavings; 1 preemptions; job time 1m 50s; pldi time 144332259; new-fixed pldi cputime 249418666
			% 33 interleavings; 1 preemptions; job time 10s; pldi time 43954391; new-fixed pldi cputime 308100075
			& 2,4	&\cpu{*2979.29}&\wtm{*390.34}&\ints{*1457}&\ETA{--}&\ETA{>1h}&\ETA{--}& \ETA{61882736} \\
			% 1457 interleavings; 4 preemptions; job time 3m 11s; pldi time 390342652; new-fixed pldi cputime 2979291067
			% maximal: progress: 22854/61882736 brs (0.036931%), ETA 9y 307d 18h 4m 41s (elapsed 51m 54s)
			% update for below, post thr-bottom1 fix - ok it FABed after all;
			% not sure if higher ulimit still necessary but i still used it to be safe
			% (old comment:) this quicksand jerk required ulimit -n higher than 1024
			%
			% the jerk didn't even find the bug anymore, after dpor-preferred tid
			% heres the stats from how it did before (older comment:)
			% & 2,4	&\cpu{*7598.19}&\wtm{*970.67}&\ints{*2633}&\ETA{--}&\ETA{>1h}&\ETA{--}& \ETA{79020639} \\
		\cline{2-9}
			& 3,1	&\cpu{*87.10}&\wtm{*14.81}&\ints{*14}	&\cpu{94.08}&\wtm{23.60}&\ints{24}&\ETA{207575} \\
			% 14 interleavings tested; 3 preemptions; job time 8s; pldi time 14807247; new-fixed pldi cputime 87102600
			% 24 interleavings tested; 3 preemptions; job time 5s; pldi time 23600679; new-fixed pldi cputime 94082790
			& 3,2	&\cpu{*3672.84}&\wtm{*475.03}&\ints{*145}&\ETA{--}&\ETA{>1h}&\ETA{--}& \ETA{1635075071} \\
			% maximal: progress: 6571/1635075071 brs (0.000402%), ETA 489418866284y 68d 4h 41m 44s (elapsed 29m 54s)
			% nb. this quicksand one is actually the slower of 2 FABs, but so close in time to each other
			% that the one with substantially fewer interleavings. they were:
			% 145 interleavings; 1 preemptions; job time 23s; pldi time 475030016; new-fixed pldi cputime 3672836411
			% 225 interleavings; 2 preemptions; job time 32s; pldi time 465874053; new-fixed pldi cputime 3601564200
			% not like maximal found the bug at all so it's not like there's a comparison to make fair/unfair anyway
		\hline
		{\tt lock\_fast}
		%({\tt spinlock})
			& 2,1	&\cpu{18.33}&\wtm{5.19}&\ints{2}&\cpu{*3.12}&\wtm{*3.12}&\ints{2}& \ETA{4} \\
			% 2 interleavings; 0 preemptions; job time 4s; pldi time 5191170; new-fixed pldi cputime 18328224
			% 2 interleavings; 0 preemptions; job time 2s; pldi time 3119081; new-fixed pldi cputime 3118912
		(perf)
			& 9,9	&\cpu{22.43}&\wtm{6.24}&\ints{2}&\cpu{*4.71}&\wtm{*4.71}&\ints{2}& \ETA{inf} \\
			% 2 interleavings; 0 preemptions; job time 3s; pldi time 6238515; new-fixed pldi cputime 22431755
			% 2 interleavings; 0 preemptions; job time 3s; pldi time 4712012; new-fixed pldi cputime 4711884
	\end{tabular}
	\end{center}
	\caption[Landslide's bug-finding performance on various test configurations.]
		{Landslide's bug-finding performance on various test configurations.
		Iterative Deepening (\cref{sec:quicksand-id}), optimized for fast bug-finding,
		is compared against Maximal State Space mode (\cref{sec:landslide-quicksand-options}),
		optimized for fast verification.
		For each, I list the CPU-time and wall-clock time elapsed,
		plus the number of interleavings tested in the ultimately buggy state space
		until the bug was found.
		* marks the winning measurements between each series.
		Lastly, state space estimation (\cref{sec:landslide-estimate})
		%though approximate at best,
		confers a sense of the exponential explosion.
	}
	\label{tab:tm-buges}
\end{table*}

\Cref{tab:tm-buges}
presents the bug-finding results.
I configured Landslide to run Quicksand's Iterative Deepening algorithm on 8 cores, %(\cref{sec:quicksand-id}),
shown left,
as well as to prioritize the maximal state space, %(\cref{sec:landslide-quicksand-options}),
shown right,
each with a time limit of 1 hour.
Tests {\tt htm1}, {\tt swapbug}, and {\tt avl\_insert} were run with {\tt landslide -X}
(i.e., retry aborts enabled and different abort codes not distinguished);
{\tt lock\_fast} was run with {\tt landslide -X -A -S}
(i.e., suppressing retry aborts, due to the spinlock's use of a retry loop confirmed with manual inspection).

\subsubsection{Finding bugs quickly}

As the test parameters increase,
the multiplicative factor in bug-finding speed (2-4x, eyeballing) is generally smaller
than that of the total number of interleavings (10-1000x).
In other words,
should transactional bugs exist,
Landslide is likely to find them reasonably quickly %in these transactional programs
despite prohibitive exponential explosion in total state space size.
This corroborates the results from \Cref{chap:quicksand},
extending its good news to the world of HTM.

\subsubsection{Finding new bugs}

In addition to the bugs I intentionally wrote in {\tt htm1} and {\tt swapbug},
%to my pleasant surprise % only the former of these is a surprise now
Landslide also found two bugs in the ``real-world'' transactional algorithms I tested.

\begin{figure}[t]
	\begin{center}
		\begin{tabular}{l}
		\texttt{\flow{while} (\_retry);} \\
		\texttt{\flow{if} (\call{\_xbegin}() == \const{SUCCESS}) \{} \\
		\texttt{~~~~\hilight{brickred}{tie}(\_root,inserted) = \call{\_insert}(\_root,n); } \\
		\texttt{~~~~\call{\_xend}();} \\
		\texttt{\} \flow{else} \{} \\
		\texttt{~~~~\call{pthread\_mutex\_lock}(\&\_tree\_lock);} \\
		\texttt{~~~~\_retry = \const{true};} \\
		\texttt{~~~~\hilight{brickred}{tie}(\_root,inserted) = \call{\_insert}(\_root,n); } \\
		\texttt{~~~~\_retry = \const{false};} \\
		\texttt{~~~~\call{pthread\_mutex\_unlock}(\&\_tree\_lock);} \\
		\texttt{\}} \\
		\end{tabular}
	\end{center}
	%\caption{mario man had a very large house. or a castle.}
	\caption[Code from {\tt htmavl.hpp} showing the segfault bug Landslide found.]
		{Unmodified code from {\tt htmavl.hpp} showing the previously-unknown segfault bug
		Landslide found in {\tt avl\_insert}.
		The transaction path fails to check {\tt \_retry},
		leading to data races and corruption just as in {\tt htm1}.
		}
	\label{fig:avlbug}
\end{figure}

\begin{itemize}
	\item
	{\bf Atomicity violation.}
	{\tt avl\_insert} with any parameters higher than (2,1)
	exposed a previously-unknown bug in the transactional AVL tree.
	\Cref{fig:avlbug} shows the root cause, essentially the {\tt htm1} bug in disguise.
	This manifested
	alternately
	% landslide-trace-1533850116.527842.html
	% landslide-trace-1533850236.873032.html
	as a segfault (for test parameters (2,2) and (3,1))
	and
	% landslide-trace-1533854371.201589.html
	as a consistency-check assertion failure (for test parameters (2,3)).
	The presence of {\tt while (\_retry);} makes the necessary preemption window extremely small
	(between it and {\tt \_xbegin()}),
	making the bug extremely unlikely to manifest under stress testing,
	but Landslide is blind to such matters of chance.

	As a matter of full disclosure,
	I noted that the loop does not affect the test's possible behaviours,
	only its {\em likely} ones,
	and so removed it to make the test more Landslide-friendly.
	%keep the state space size manageable.
	%Lest any reader think I hacked the test after already knowing about the bug,
	To dispel any doubt about bias or test hacking,
	I confirmed that Landslide still finds the bug with the spin loop unmodified,
	on (3,1) in the same 53 interleavings,
	although it suffers resource exhaustion
	% from dealing with the loop
	on parameters of (2,2) or greater.
	% specifically, it runs out of file descriptors, since on bochs you need 1 fd per PP,
	% and the branch just gets 1000+ PPs deep
	\item
	{\bf Spurious spinlock abort.}
	{\tt lock\_fast} discovered a spurious transactional-path write conflict
	in the {\tt spinlock} HTM-lock implementation.%
	\footnote{{\tt lock\_fast}'s unusual (9,9) parameter shows that this state space size is constant: %in $(K,N)$:
	DPOR will always either deem all thread transitions independent and end exploration immediately,
	or the test's assertion will trip as soon as the first conflict is found.}
	This ``performance bug'' causes the lock to suffer slow-path spin-locking
	even in cases where the user's thread transitions are completely independent
	(for example, locking the root of an AVL tree,
	then traversing in different directions to make disjoint modifications).
	The test case detects this with {\tt \_xtest()}.
	\Cref{fig:spinlockbug} shows the root cause:
	the {\tt isfree()} routine (corresponding to the AVL's {\tt \_retry})
	used an atomic compare-and-swap that would always write to memory even without modifying it.
	I corrected this in
	% the
	{\tt spin\_fixed}
	%implementation
	by replacing it
	with a normal read
	(being used only in the transactional path, no barriers are required to protect it).
	A cursory search on Github found one user of this code,
	a transactional LevelDB implementation \cite{htm-leveldb-github},
	whose author had also noticed and corrected this problem in the same way.

	As another matter of full disclosure,
	I noticed this bug through manual inspection
	while adapting the spinlock's client code to be Landslide-friendly,
	then wrote {\tt lock\_fast} specifically to target this behaviour,
	so unlike {\tt avl\_insert}, it does not count as Landslide finding a previously-unknown bug.
	However, I feel in retrospect that how and when
	an HTM-backed concurrency abstraction
	will fall into its slow path
	is a reasonable performance property for a user to want to verify,
	so I consider Landslide confirming the bug (and later verifying its absence, in \cref{sec:tm-verif})
	a positive result anyway.
\end{itemize}

\begin{figure}[t]
	\begin{center}
		\begin{tabular}{l}
		\texttt{\ctype{bool} \call{hle\_spinlock\_isfree}(\ctype{spinlock\_t} *lock) \{} \\
		\texttt{~~~~\ccomment{// XXX: should be "return lock->v == 0;"}} \\
		\texttt{~~~~\flow{return} \call{\_\_sync\_bool\_compare\_and\_swap}(\&lock->v, \const{0}, \const{0});} \\
		\texttt{\}} \\
		\texttt{\ctype{void} \call{rtm\_spinlock\_acquire}(\ctype{spinlock\_t} *lock) \{} \\
		\texttt{~~~~\flow{if} ((tm\_status = \call{\_xbegin}()) == \const{\_XBEGIN\_STARTED}) \{} \\
		\texttt{~~~~~~~~\flow{if} (\call{hle\_spinlock\_isfree}(lock)) \flow{return};} \\
		\texttt{~~~~~~~~\call{\_xabort}(\const{0xff});} \\
		\texttt{~~~~\} \flow{else} \{} \\
		\texttt{~~~~~~~~\ccomment{// ... retrying \&c abbreviated for brevity ...}} \\
		\texttt{~~~~~~~~\call{hle\_spinlock\_acquire}(lock);} \\
		\texttt{~~~~\}} \\
		\texttt{\}} \\
		\texttt{\ctype{void} \call{rtm\_spinlock\_release}(\ctype{spinlock\_t} *lock)\{} \\
		\texttt{~~~~\flow{if} (\call{hle\_spinlock\_isfree}(lock)) \{} \\
		\texttt{~~~~~~~~\call{\_xend}();} \\
		\texttt{~~~~\} \flow{else} \{} \\
		\texttt{~~~~~~~~\call{hle\_spinlock\_release}(lock);} \\
		\texttt{~~~~\}} \\
		\texttt{\}} \\
		\end{tabular}
	\end{center}
	\caption[Code from {\tt spinlock-rtm.c} %, modified only to remove unrelated logic for brevity,
		showing the performance bug Landslide found.]
		{Code from {\tt spinlock-rtm.c}, modified only to remove unrelated logic for brevity,
		showing the performance bug Landslide found in {\tt lock\_fast}({\tt spinlock}).
		The {\tt isfree()} routine uses an atomic read-and-write operation where just a read would suffice,
		which leads to superfluous memory conflicts in the transactional path
		(seen at both of its callsites below).
	}
	\label{fig:spinlockbug}
\end{figure}

Note that in the AVL tree bug,
the code's author was the very same person who proposed the protocol in \Cref{fig:htm-fixed},
yet still got it wrong once, having to write it out by hand throughout both data structures.
This motivates the need for model checking such programs,
no matter how much of a concurrency expert the author may be.
It also suggests HTM primitives should be encapsulated behind higher-level abstractions,
such as lock elision \cite{lock-elision}
or a simple spinlock \cite{spinlock-rtm-github},
which can be verified in isolation with smaller state spaces
then trusted in turn when checking their client programs \cite{dbug-phdthesis}.
\cref{sec:tm-verif} explores this further.

Regarding the spinlock bug,
as HTM is fundamentally a performance-minded concurrency extension,
the user may also care about more probabilistic properties of her code,
such as requiring a transaction abort rate below a certain threshhold
owing to the nature of its workload.
Landslide cannot in general test for performance degradation bugs,
because all interleavings are equal in Landslide's eyes, and probability is no object.
However, {\tt lock\_fast} illustrates that model checking can still check some interesting performance properties
as long as the element of probability can be removed.
Future work may attempt to verify a wider range of performance properties,
but with a hybrid approach between model checking and what other technique is not yet known.

\subsubsection{Performance}

Quicksand's ability to find bugs in fewer distinct interleavings (i.e., overall smaller state spaces)
does not necessarily correlate with better performance in terms of CPU-time.
Comparing \Cref{tab:tm-buges}'s trends
to the break-even point in Quicksand's evaluation (\cref{sec:quicksand-eval}),
most of these tests are too small for its approach to pay off,
with {\tt swapbug} and {\tt avl\_insert} as its notable wins.
While plenty more wins were observed in \cref{sec:quicksand-eval},
%an upper bound also exists:
%in {\tt swapbug}(3,2),
%not even the {\em minimal} state space was completed in time.
this suggests
future MCs could prioritize state spaces using not just size estimation
but a hybrid approach also conisdering state space maximality and preemption bounds \cite{chess-icb}
to soften the trade-off
both for smaller tests and for verification.

\subsection{Verification}
\label{sec:tm-verif}

For the test cases with no bugs found, I sought to provide Landslide's verification guarantee
(\cref{sec:quicksand-soundness})
for up to as many threads and test iterations as possible under a reasonable time limit.
The results, obtained using the same {\tt -X -M} configuration options as
in the previous section,
%were used in Table~\ref{tab:tm-buges},
are shown in \Cref{tab:tm-verifs} in the ``Baseline DPOR'' column.
%Of course, because stateless model checking is a perpetual existence
%of staring up the sheer cliff face that is the exponential curve,
%many tests were only able to achieve verification for only modest values of $(K, N)$
%before becoming intractable.
%After a timeout of 10 wall-clock hours,
%I measured these tests'
%estimated remaining time and state space size and list those instead,
%\ETAdag{emphasized thusly}.
For test configurations which could not be verified
within 10 wall-clock hours,
I report their estimated state space size and runtime measured after that timeout instead,
\ETAdag{typeset thusly}.

\begin{table}[p]
	\begin{center}
		\footnotesize
		\begin{tabular}{cc||r|r||r|r||r|r}
			& & \multicolumn{2}{c||}{\bf Baseline DPOR}
			  & \multicolumn{2}{c||}{\bf Retry sets ({\tt -R})}
			  & \multicolumn{2}{c}{\bf STM ({\tt -A -S})} \\
			& & \cpu{\bf cpu (s)} & \ints{\bf SS size}
			  & \cpu{\bf cpu (s)} & \ints{\bf SS size}
			  & \cpu{\bf cpu (s)} & \ints{\bf SS size} \\
			\bf test & \bf K,N
			  & \em (or \ETAdag{\bf \em est.}) & \em (or \ETAdag{\bf \em est.})
			  & \em (or \ETAdag{\bf \em est.}) & \em (or \ETAdag{\bf \em est.})
			  & \em (or \ETAdag{\bf \em est.}) & \em (or \ETAdag{\bf \em est.}) \\
			\hline
			\hline
			{\tt htm2}
			& 2,1 & \cpu{18.26}	& \ints{22}	& \cpu{17.99}	& \ints{15}	& \cpu{3.25}	& \ints{4} \\
			& 2,2 & \cpu{63.61}	& \ints{1446}	& \cpu{26.40}	& \ints{334}	& \cpu{28.04}	& \ints{286} \\
			& 2,3 & \cpu{3429.04}	& \ints{86536}	& \cpu{196.18}	& \ints{6366}	& \cpu{997.91}	& \ints{24740} \\
			& 2,4 & \ETAdag{29d 8h}	&\ETAdag{2710056}&\cpu{4771.17}	& \ints{123140}	&\ETAdag{5d 1h}	&\ETAdag{1792330} \\
			% [JOB 4] progress: 683895/2710056 brs (25.235453%), ETA 28d 21h 54m 32s (elapsed 9h 59m 54s)
			% [JOB 4] progress: 714585/1792330 brs (39.869037%), ETA 4d 14h 46m 38s (elapsed 9h 59m 54s)
			& 3,1 & \cpu{43.24}	& \ints{774}	& \cpu{24.02}	& \ints{224}	& \cpu{22.79}	& \ints{140} \\
			& 3,2 &\ETAdag{1y 299d}&\ETAdag{4510472}&\ETAdag{294d 5h}&\ETAdag{8185793}&\ETAdag{10d 16h}&\ETAdag{2322150}\\
			% argh.. since re running the retry set i feel like i have to re run the others...
			% the small state spaces though all end up similar perf numbers though;
			% i wonder why this one alone broke 1mil?
			% [JOB 4] progress: 759740/4510472 brs (16.843911%), ETA 1y 298d 22h 14m 13s (elapsed 9h 59m 55s)
			% [JOB 4] progress: 1133447/8185793 brs (13.846513%), ETA 293d 19h 13m 36s (elapsed 9h 59m 55s) % good retry
			% [JOB 4] progress: 646606/2322150 brs (27.845142%), ETA 10d 5h 50m 15s (elapsed 9h 59m 55s)
			% old, owned retry set --
			% ls-id-log-20180817141549.log.ByMMKx.gz
			% [JOB 4] BUG FOUND: landslide-trace-1534529755.156051.html (981774 interleavings; 132 preemptions; ...)
			% it's almost certainly getting retry set blocked
			% so you just have to detect a retry set block deadlock, and un deadlock it
			% this should not affect other tests (not even in cpu time (beyond icache considerations))
			% when you fix it and re run it, remember to use leave-logs, and put a ALWAYS prince in landslide
			% to see when it happens
			& 4,1 & \cpu{9225.52}	& \ints{212146}	& \cpu{376.70}	& \ints{11973}	& \cpu{2158.93}	& \ints{44995} \\
			%% i don't even remember which test case this commented stuff was for anymore
			% measured with !CONSIDER_ONLY_MOST_RECENT_DPOR_PREFERRED_TID
			%& 2,2 & \cpu{72.40}	& \ints{1970}	& \cpu{27.28}	& \ints{356}	& \cpu{26.84}	& \ints{288} \\
			%& 3,1 &\cpu{38.76}	& \ints{708}	& \cpu{23.84}	& \ints{224}	& \cpu{23.38}	& \ints{140} \\
			%& 3,2 &\ETAdag{vvv}	& \ETAdag{vvv}	& \ETAdag{vvv}	& \ETAdag{vvv}	&\ETAdag{4d 21h}&\ETAdag{3456001}\\
			%[JOB 4] progress: 709997/3456001 brs (20.543887%), ETA 4d 10h 36m 25s (elapsed 9h 59m 54s)
			\hline
			{\tt counter}
			& 2,1 & \cpu{6.92}	& \ints{10}	& \cpu{6.81}	& \ints{8}	& \cpu{3.28}	& \ints{4} \\
			& 2,2 & \cpu{13.27}	& \ints{190}	& \cpu{10.59}	& \ints{102}	& \cpu{8.73}	& \ints{48} \\
			& 2,3 & \cpu{155.09}	& \ints{3970}	& \cpu{67.82}	& \ints{1558}	& \cpu{40.08}	& \ints{904} \\
			& 2,4 & \cpu{3664.66}	& \ints{86950}	& \cpu{1150.34}	& \ints{25398}	& \cpu{805.85}	& \ints{19128} \\
			& 3,1 & \cpu{11.26}	& \ints{120}	& \cpu{9.30}	& \ints{64}	& \cpu{8.26}	& \ints{40} \\
			& 3,2 & \cpu{2572.13}	& \ints{60606}	& \cpu{2363.41}	& \ints{44862}	& \cpu{639.62}	& \ints{14304} \\
			& 4,1 & \cpu{129.25}	& \ints{3006}	& \cpu{64.91}	& \ints{1296}	& \cpu{40.78}	& \ints{848} \\
			\hline
			{\tt swap}
			& 2,1 & \cpu{65.72}	& \ints{99}	& \cpu{66.05}	& \ints{59}	& \cpu{3.40}	& \ints{4} \\
			& 2,2 & \cpu{18124.06}	& \ints{277824}	& \cpu{1030.23}	& \ints{19542}	& \cpu{703.82}	& \ints{11600} \\
			& 3,1 & \cpu{3820.64}	& \ints{60912}	& \cpu{608.48}	& \ints{10706}	& \cpu{89.69}	& \ints{1014} \\
			\hline
			{\tt fig63}
			& 2,1 & \cpu{7.12}	& \ints{10}	& \cpu{6.92}	& \ints{6}	& \cpu{3.40}	& \ints{1} \\
			& 2,2 & \cpu{9.24}	& \ints{108}	& \cpu{8.61}	& \ints{76}	& \cpu{3.46}	& \ints{1} \\
			& 2,3 & \cpu{54.68}	& \ints{1934}	& \cpu{31.76}	& \ints{977}	& \cpu{3.39}	& \ints{1} \\
			& 2,4 & \cpu{1054.80} 	& \ints{36600} 	& \cpu{417.99}	& \ints{14512}	& \cpu{3.58}	& \ints{1} \\
			% nb this 22 would be a 21 with !CONSIDER_ONLY... ..but, if it's the biggest loss, really n b d
			& 3,1 & \cpu{11.15}	& \ints{148}	& \cpu{7.76}	& \ints{22}	& \cpu{3.53}	& \ints{1} \\
			& 3,2 & \cpu{717.53}	& \ints{21642}	& \cpu{217.30}	& \ints{6467}	& \cpu{3.52}	& \ints{1} \\
			& 4,1 & \cpu{111.83}	& \ints{3064}	& \cpu{11.32}	& \ints{130}	& \cpu{3.49}	& \ints{1} \\
			\hline
			{\tt avl\_insert}
			& 2,1 & \cpu{672.68}	& \ints{15125}	& \cpu{307.94}	& \ints{6287}	& \cpu{136.60}	& \ints{2774} \\
			\hline
			{\tt avl\_fixed}
			& 2,1 & \cpu{739.60} 	& \ints{20459} 	& \cpu{332.32}	& \ints{9675}	& \cpu{122.40}	& \ints{2774} \\
			& 2,2 &\ETAdag{58880y}&\ETAdag{4.62$\times\mathit{10^7}$}&\ETAdag{51191y}&\ETAdag{3.33$\times\mathit{10^7}$}&\ETAdag{36393y}&\ETAdag{3.51$\times\mathit{10^7}$}\\
			% [JOB 19] progress: 366191/46150466 brs (0.793472%), ETA 58879y 219d 6h 4m 16s (elapsed 9h 59m 28s)
			% [JOB 19] progress: 351391/33318009 brs (1.054658%), ETA 51191y 111d 17h 37m 55s (elapsed 9h 59m 28s)
			% progress: 302022/35094477 brs (0.860597%), ETA 36392y 304d 6h 53m 45s (elapsed 9h 58m 49s)
			& 3,1 & \ETAdag{5y 292d}&\ETAdag{2.87$\times\mathit{10^9}$}&\ETAdag{735y 212d}&\ETAdag{3.35$\times\mathit{10^8}$}&\ETAdag{97y 100d}&\ETAdag{1.70$\times\mathit{10^9}$} \\
			% [JOB 18] progress: 447687/2873664642 brs (0.015579%), ETA 5y 291d 11h 13m 30s (elapsed 9h 59m 41s)
			% [JOB 25] progress: 18151/335288683 brs (0.005414%), ETA 735y 211d 9h 1m 8s (elapsed 22m 25s)
			% [JOB 18] progress: 357586/1698185036 brs (0.021057%), ETA 97y 99d 8h 44m 33s (elapsed 9h 59m 41s)
			\hline
			{\tt map\_basic}
			& 2,1 & \cpu{1950.48}	& \ints{30719}	& \cpu{867.68}	& \ints{13237}	& \cpu{367.91}	& \ints{5446} \\
			& 2,2 & \ETAdag{1y 176d}&\ETAdag{7.55$\times\mathit{10^7}$}&\ETAdag{11y 94d}&\ETAdag{9.57$\times\mathit{10^7}$}&\ETAdag{13d 11h}&\ETAdag{565334} \\
			% [JOB 22] progress: 87777/75516602 brs (0.116235%), ETA 1y 176d 16h 38m 50s (elapsed 3h 1m 48s)
			% [JOB 22] progress: 128715/95718464 brs (0.134472%), ETA 11y 93d 8h 25m 39s (elapsed 4h 52m 39s)
			% [JOB 22] progress: 18380/565334 brs (3.251171%), ETA 13d 1h 5m 3s (elapsed 38m 33s)
			& 3,1 &\ETAdag{6y 133d}&\ETAdag{1.85$\times\mathit{10^9}$}&\ETAdag{8y 282d}&\ETAdag{7.90$\times\mathit{10^8}$}&\ETAdag{4y 156d}&\ETAdag{8.88$\times\mathit{10^8}$} \\
			% [JOB 22] progress: 284530/1847957714 brs (0.015397%), ETA 6y 132d 6h 20m 26s (elapsed 9h 49m 5s)
			% [JOB 22] progress: 282358/789973638 brs (0.035743%), ETA 8y 281d 14h 50m 3s (elapsed 9h 49m 54s)
			% [JOB 22] progress: 274532/888242178 brs (0.030907%), ETA 4y 155d 23h 50m 51s (elapsed 9h 47m 5s)
			\hline
			{\tt map\_basicer}
			& 2,1 & \cpu{28.40} 	& \ints{150} 	& \cpu{26.55}	& \ints{94}	& \cpu{14.88}	& \ints{9} \\
			& 2,2 &\ETAdag{11h 17m}	&\ETAdag{727759}&\cpu{16455.24}	& \ints{283756}	& \cpu{1285.06}	& \ints{21684} \\
			% [JOB 7] progress: 634229/727759 brs (87.148190%), ETA 1h 16m 33s (elapsed 9h 36m 22s)
			% total CPU time consumed: 10h 11s (36011250596 usecs) (core saturation: 12%)
			& 3,1 &\ETAdag{26h 15m}	&\ETAdag{1451708}&\cpu{21153.64}& \ints{366030}	& \cpu{705.04}	& \ints{12707} \\
			% [JOB 5] progress: 655856/1451708 brs (45.178208%), ETA 16h 14m 59s (elapsed 9h 59m 50s)
			% total CPU time consumed: 10h 10s (36010204482 usecs) (core saturation: 12%)

			%%%% nb. below is a fake version of the results, basically measured with abort set reduction,
			%%%% after implementing abort set coalescing/merging, but not with either "dpor preferred tid"
			%%%% or "trust thr join" optimizations.
			%%%% this causes both (1) baseline to look worse than it should, and (2) abort sets to be
			%%%% pretty inconsistent about whether they actually provide reduction, or accidentally inflate
			%%%% the state space instead.
			%%%% nevertheless, i'm keeping it around for posterity so you can compare, e.g., this fake baseline
			%%%% against real baseline, to evaluate the impact of the non-abort-set dpor optimizations alone.
			% {\tt htm2}
			% & 2,1 & \cpu{16.30}	& \ints{22}	& \cpu{17.93}	& \ints{15}	& \cpu{3.32}	& \ints{4} \\
			% & 2,2 & \cpu{69.80}	& \ints{1970}	& \cpu{27.87}	& \ints{408}	& \cpu{26.59}	& \ints{288} \\
			% & 2,3 &\cpu{3510.00}	& \ints{104914}	& \cpu{}	& \ints{}	& \cpu{}	& \ints{} \\
			% %& 2,4 &	\ETAdag{...}	& \ETAdag{...}	\\
			% & 3,1 &\cpu{41.07}	& \ints{941}	& \cpu{23.86}	& \ints{273}	& \cpu{23.03}	& \ints{162} \\
			% & 3,2 &	\ETAdag{9d 19h}	& \ETAdag{5058149}	\\ % add 10h ok
			% % [JOB 4] progress: 923822/5058149 brs (18.264030%), ETA 9d 8h 30m 38s (elapsed 9h 59m 55s)
			% & 4,1 & \cpu{12699.65}	& \ints{344240}	& \cpu{}	& \ints{}	& \cpu{}	& \ints{} \\
			% \hline
			% {\tt counter}
			% & 2,1 & \cpu{5.40}	& \ints{10}	& \cpu{7.00}	& \ints{8}	& \cpu{3.22}	& \ints{4} \\
			% & 2,2 & \cpu{10.91}	& \ints{190}	& \cpu{10.57}	& \ints{122}	& \cpu{7.90}	& \ints{48} \\
			% & 2,3 & \cpu{134.21}	& \ints{3970}	& \cpu{85.29}	& \ints{2376}	& \cpu{35.00}	& \ints{904} \\
			% & 2,4 & \cpu{3009.40}	& \ints{86950}	& \cpu{}	& \ints{}	& \cpu{}	& \ints{} \\
			% & 3,1 & \cpu{8.30}	& \ints{120}	& \cpu{9.40}	& \ints{74}	& \cpu{7.55}	& \ints{40} \\
			% & 3,2 & \cpu{2038.90}	& \ints{60606}	& \cpu{}	& \ints{60140}	& \cpu{486.56}	& \ints{14304} \\
			% & 4,1 & \cpu{101.95}	& \ints{3006}	& \cpu{82.16}	& \ints{2148}	& \cpu{32.95}	& \ints{848} \\
			% \hline
			% {\tt swap}
			% % these definitely have a different retry sets data point, run with some coalescing updates,
			% % but i think before the dpor-preferred-tid search ordering fix
			% % i think fig63 might have been updated this way as well? but not sure
			% %& 2,1 & \cpu{38.79}	& \ints{99}	& \cpu{37.33}	& \ints{59}	& \cpu{3.25}	& \ints{4} \\
			% %& 2,2 & \cpu{10299.27}	& \ints{233396}	& \cpu{433.01}	& \ints{11588}	& \cpu{382.28}	& \ints{9046} \\
			% %& 3,1 & \cpu{2359.85}	& \ints{59850}	& \cpu{382.27}	& \ints{9712}	& \cpu{71.41}	& \ints{1014} \\
			% & 2,1 & \cpu{38.79}	& \ints{99}	& \cpu{}	& \ints{}	& \cpu{3.25}	& \ints{4} \\
			% & 2,2 & \cpu{10299.27}	& \ints{233396}	& \cpu{}	& \ints{}	& \cpu{382.28}	& \ints{9046} \\
			% & 3,1 & \cpu{2359.85}	& \ints{59850}	& \cpu{}	& \ints{}	& \cpu{71.41}	& \ints{1014} \\
			% \hline
			% {\tt fig63}
			% & 2,1 & \cpu{6.66}	& \ints{18}	& \cpu{6.97}	& \ints{6}	& \cpu{3.10}	& \ints{1} \\
			% & 2,2 & \cpu{11.52}	& \ints{244}	& \cpu{8.75}	& \ints{76}	& \cpu{3.12}	& \ints{1} \\
			% & 2,3 & \cpu{107.04}	& \ints{4428}	& \cpu{32.26}	& \ints{1060}	& \cpu{3.13}	& \ints{1} \\
			% & 3,1 & \cpu{13.68}	& \ints{296}	& \cpu{7.55}	& \ints{21}	& \cpu{3.26}	& \ints{1} \\
			% & 3,2 & \ETAdag{}	& \ETAdag{}	& \cpu{126.65}	& \ints{4042}	& \cpu{3.35}	& \ints{1} \\
			% & 4,1 & \cpu{220.27}	& \ints{7382}	& \cpu{11.13}	& \ints{130}	& \cpu{3.27}	& \ints{1} \\
			% \hline
			% {\tt avl\_insert}
			% & 2,1 & \cpu{491.98}	& \ints{15056}	& \cpu{228.23}	& \ints{6296}	& \cpu{98.77}	& \ints{2776} \\
			% % (2,2) - bug found on regular w/ 127 ints, on retry sets w/ 201 ints, and codes w/ 45 ints
			% \hline
			% {\tt avl\_fixed}
			% & 2,1 & \cpu{572.49} 	& \ints{20420} 	& \cpu{284.37}	& \ints{9685}	& \cpu{99.16}	& \ints{2776} \\
			% & 2,2 & \ETAdag{133y}	& \ETAdag{51484570} % added 10hr but dosent matter
			% % the order of magnitude of these time estimates is basically completely random
			% % and fluctuates wildly from branch to branch
			% % [JOB 19] progress: 408805/51484570 brs (0.794034%), ETA 132y 356d 8h 55m 16s (elapsed 9h 59m 19s)
			% 	& \ETAdag{}	& \ETAdag{}
			% 	& \ETAdag{1060y}	& \ETAdag{31302769} \\
			% % for example if this one had timed out just 7 branches earlier, it woulda been 6y 246d instead
			% % [JOB 19] progress: 321936/31302769 brs (1.028459%), ETA 1059y 260d 9h 16m 8s (elapsed 9h 59m 0s)
			% & 3,1 & \ETAdag{30y 320d}	& \ETAdag{93493027}	\\
			% % [JOB 18] progress: 492976/93493027 brs (0.011233%), ETA 30y 319d 14h 48m 32s (elapsed 9h 59m 35s)
		\end{tabular}
	\end{center}
	\caption[Transactional tests verified (or not) by Landslide.]
		{Transactional tests verified (or not) by Landslide.
		Run with {\tt -M -X}, plus any additional reduction options listed.
		``Baseline DPOR'' always tests every abort path, without distinguishing among failure reasons
		(i.e., injecting only {\tt \_XABORT\_RETRY}).
		``Retry sets'' skips equivalent success path and/or retry abort reorderings (\cref{sec:tm-retrysets});
		``STM'' suppresses retry aborts and dynamically detects when to inject conflict aborts and so on
		(\cref{sec:txn-abort-codes}).
		State space estimates measured after a timeout of 10 hours (and include those 10 hours in the predicted total).
		}
	\label{tab:tm-verifs}
\end{table}

\subsubsection{Interpreting the verification guarantee}

Landslide was able to verify most of these tests for a fair range of thread and iteration counts,
often reaching up to 2 threads with 3-4 iterations each, 3 with 2 each, or 4 with 1.
In the case of {\tt htm2}, for example,
verifying up to $(K, N)$ represents a guarantee that,
even repeating \Cref{fig:htm-fixed}'s atomicity protocol $N$ times
in any scheduling sequence or combination of transaction aborts,
it is impossible for $K$ threads to violate the intended atomicity property
(i.e., get 2 threads in the critical section simultaneously).

It is difficult to discern from prior work a concrete standard
for what values of $(K, N)$ constitute a ``good'' degree of verification.
One recent paper \cite{tsopso},
which likewise extended DPOR with a new dimension of concurrency (weak memory orderings),
reported verifying programs with up to 10 concurrent events,
presumably shared memory accesses.
Another \cite{optimal-dpor}
reported test cases with as many as 19 threads,
although with what must be very little synchronization or memory conflicts,
%although with what nature of synchronization or memory conflicts is not clear,
as even their baseline DPOR checked only 4096 interleavings on that test.
%
\Cref{tab:tm-test-case-evence} confers a sense of the complexity of this evaluation's test suite,
with the final column showing the approximate maximum number of preemption points
reached among Landslide's verifications
(i.e., $({\sf txn}+{\sf sync}+{\sf race}) \times K \times N$ for the highest $K \times N$ completed).
%%%% on second thought, i'm not gonna do this -- any branch that xend shows up in, is gonna lose any of the pps coming
%%%% from the failure path's mxes or data races, so xend will never be an *additional* pp
%%%% and of course it will scare the committee
% \footnote{The 1.5 multiplier on {\sf txn} represents the fact that for each transaction,
% the {\tt xend} preemption point marking the end of its success path shows up about half the time,
% the other half being failure injections.
% % in reality it's usually less than half, because (1) drpps can show up in failure paths, which will get interleaved around,
% % thereby giving the abort path more interleavings, and (2) under -A -S, you might get multiple codes for duplicate abort paths
% % abort sets could theoretically push them above half but probably not in any test case here
% }).

% TODO: check figure placement
\begin{table}[t]
	\begin{center}
		\small
		\begin{tabular}{c||c|c|c||c}
			{\bf test} & {\bf \#txn} & {\bf \#sync} & {\bf \#race} & {\bf max events verified} \\
			\hline
			{\tt htm2}		& 1	& 2	& 4	& 42	\\ % 2,3
			{\tt counter}		& 1	& 0	& 1	& 16	\\ % 2,4
			{\tt swap}		& 1	& 4	& 8	& 52	\\ % 2,2
			{\tt fig63}		& 1	& 0	& $K$-1	& 18	\\ % 3,2 k=3 -> 18; 2,4 -> 16
			{\tt avl\_insert}	& 1	& 2	& 7	& 20	\\ % 2,1 nb. only 5 drpps under -A -S but w/e
			% nb less than 9 on stm
			{\tt avl\_fixed}	& 1	& 2	& 9	& 24	\\ % 2,1 nb. again, only 5 under -A -S
			% nb stm takes it down from 13 to 10 races
			{\tt map\_basic}	& 1	& 4	& 13	& 36	\\ % 2,1
			% mysteriously had this at 6 at one point but ive no idea where that came from
			{\tt map\_basicer}	& 1	& 2	& 5	& 32	\\ % 2,2

			{\tt lock}({\tt spinlock})	& 1	& 2	& 4	& 28	\\ % 2,2
			{\tt lock}({\tt spin\_fixed})	& 1	& 2	& 4	& 28	\\ % 2,2 or 4,1
			{\tt lock}({\tt mutex})		& 1	& 2	& 2	& 40	\\ % 2,4
			% not putting lock_fast because it like, doesn't have a "state space" per se
		\end{tabular}
	\end{center}
	\caption[Number of concurrency events per iteration of each test case.]
		{Number of concurrency events per iteration of each test case.
	Note that no test used any synchronization besides mutexes % ;
	%while they do use the P2 thread API, %in the first place,
	(the P2 thread API was annotated as trusted (\cref{sec:landslide-config-landslide})
	and so does not contribute to state space size).
	Also note that ``\#race'' means the number of unique accesses identified as racy,
	rather than racing pairs
	(the other half of a pair might well be within a transaction, which cannot be preempted on).
	}
	\label{tab:tm-test-case-evence}
\end{table}

In the case of {\tt htm2},
it is easy to look at the program with human intuition and judge that,
because the protocol's only state is stored in a single boolean,
with no unbounded-capacity data structures or
contention-dependent exponential backoff loops,
it would be unimaginable that adding a 5th thread to the system
could make any difference in correctness where 4 threads could not,
or that a 4th repetition between 2 threads could make a difference where 3 could not,
and ultimately that it must safely generalize to all $(K, N)$.
% under certain invisible assumptions like the system runs with SC memory! (spoilers!!)
Generalizing the verification is not so straightforward
for more complicated algorithms,
which may involve complex conflict patterns such as tree rebalancing or map resizing.
Other formal verification approaches aside,
the user must ultimately be content with the probabilistic assurance that
as verified $K$ and $N$ increase,
the likelihood that a bug exists which requires more threads or iterations to expose grows ever lower
% (in the author's fancy, it may as well drop exponentially in the same way as the state spaces grow)
(for example, none of the bugs in \cref{sec:tm-eval-bugs} required any higher parameters than (2,2) to expose).
% and by analogy with ICB, none of the bugs in quicksand evaluation took more than a preemption bound of 3 to expose
% but it's not clear how to state this "analogy" concisely in text

Nevertheless, pushing $K$ and $N$ higher is obviously desirable,
even if it means applying reductions that require human intuition to trust \revisionminor{their soundness}. %are sound.
Moreover, much as verifying {\tt htm2}'s soundness is a positive result,
the attempts at larger data structures quickly suffered exponential explosion
for even small thread/iteration counts,
in one case
%failing to make any guarantee about even the minimum possible test configuration,
%and in another
failing to verify whether or not a previously-found bug had actually been fixed.
In the following subsections I explore three possible mitigation approaches.

\subsubsection{Retry set reduction}
\label{sec:tm-eval-retrysets}

Firstly, the middle column of \Cref{tab:tm-verifs}
shows the impact of Landslide's experimental retry set reduction (\cref{sec:tm-retrysets}).
Despite its conservative implementation,
it provides roughly 2-6x reduction in most tests, with up to 17x in extreme cases.
The biggest win is apparent in {\tt fig63},
the test contrived to induce as much reduction as possible
using transactional paths that conflict only with aborts and not with each other, and vice versa.
In {\tt fig63}(2,1), corresponding to \Cref{fig:retry-set-example},
retry-set-enabled DPOR correctly prunes down to the optimal 6 interleavings,
while the baseline treats it identically to the fully-conflicting {\tt counter}(2,1).
\Cref{fig:retry-reduction} depicts the difference between the state spaces explored by the two approaches.

\begin{figure}[t]
	\begin{center}
		\includegraphics[width=\textwidth]{retry-sets.pdf}
	\end{center}
	\caption[Number of concurrency events per iteration of each test case.]
		{Visualization of retry set state space reduction.
	On both {\tt counter}(2,1) and {\tt fig63}(2,1), baseline DPOR tests % oops
	all 10 interleavings pictured,
	with the middle 2 arising from the data-race preemption point within the abort path.
	%On {\tt counter}, all transitions marked as memory conflicts conflict with each other;
	%on {\tt fig63}, only pairs of one {\tt xend} and one {\tt data race} each conflict.
	With the reduction enabled,
	after the 4th and 6th branches
	(i.e., when preempting to reorder threads),
	Landslide activates the retry set
	indicated at the top of the next upcoming subtree,
	allowing it to identify and skip 2 redundant branches in {\tt counter}
	and 4 in {\tt fig63}.
	%to remember which parts of it to skip.
	}
	\label{fig:retry-reduction}
\end{figure}

Perhaps surprisingly, retry sets also provide reduction
even when transactional success and abort paths are fully conflicting,
(i.e., all tests besides {\tt fig63}).
With just synchronization preemption points (including {\tt \_xbegin()} and {\tt \_xend()}),
both the baseline and retry sets would explore exactly all 8 permutations of success and abort between two threads.
However, in the presence of data-race preemption points,
even for example on {\tt counter}(2,1)
(whose abort path is just 1 {\tt xadd} operation),
and whose optimal state space size should be 8 regardless of data-race preemption points),
baseline DPOR tests reorderings of one thread's transaction both with the other's failure path,
and with the other's data race therein.%
\footnote{
Not pictured in \Cref{fig:retry-reduction} is the symmetric subtree of branches 5-6 in the right half of the state space,
which would occur after branch 10, reordering the blue thread before the pink's data race.
Such would be equivalent to branches 2 and 4,
and is pruned by the normal sleep set algorithm (\cref{sec:landslide-sleepsets}, {\tt equiv\_already\_explored()})
even in baseline DPOR, with no need for retry sets.
}
Retry sets on the other hand identify and skip that equivalence
(technically speaking, retry set DPOR reorders with the data race first during depth-first search,
then skips generating a retry set for reordering the full abort path).

% --- counter-retry       2018-08-03 13:50:37.324021777 -0400
% +++ counter-baseline    2018-08-03 13:50:53.059021559 -0400
%  0--1@105c41 1-4@105c41 2-5@xbeg-delay 3-5@xbegin 4-5@xend 5-5@105c41 6-4@1001042 7-4@105c41 8-6@xbeg-delay 9-6@xbegin 10-6@xend 11-6@105c41 12-4@1001042 13-4@105c41 14-3@1028e9
% #[EXPLORE]       dbg #2/tid5 retry set: {TID6 (1,0), TID5 (1,0)}
% #[EXPLORE]       from #14/tid3, chose tid 6 (xabort injection), child of #9/tid6
%  0--1@105c41 1-4@105c41 2-5@xbeg-delay 3-5@xbegin 4-5@xend 5-5@105c41 6-4@1001042 7-4@105c41 8-6@xbeg-delay 9-6@xbegin 10-6@datarace-a 11-6@105c41 12-4@1001042 13-4@105c41 14-3@1028e9
% #[EXPLORE]       dbg #2/tid5 retry set: {TID6 (1,0), TID5 (1,0)}
% #[EXPLORE]       dbg #2/tid5 retry set: {TID6 (0,1), TID5 (1,0)}
% #[EXPLORE]       from #14/tid3, chose tid 5 (xabort injection), child of #3/tid5
%  0--1@105c41 1-4@105c41 2-5@xbeg-delay 3-5@xbegin 4-5@datarace-a 5-5@105c41 6-4@1001042 7-4@105c41 8-6@xbeg-delay 9-6@xbegin 10-6@xend 11-6@105c41 12-4@1001042 13-4@105c41 14-3@1028e9
% #[EXPLORE]       from #14/tid3, chose tid 6 (xabort injection), child of #9/tid6
%  0--1@105c41 1-4@105c41 2-5@xbeg-delay 3-5@xbegin 4-5@datarace-a 5-5@105c41 6-4@1001042 7-4@105c41 8-6@xbeg-delay 9-6@xbegin 10-6@datarace-a 11-6@105c41 12-4@1001042 13-4@105c41 14-3@1028e9
% #[EXPLORE]       from #14/tid3, chose tid 6, child of #4/tid5
%  0--1@105c41 1-4@105c41 2-5@xbeg-delay 3-5@xbegin 4-5@datarace-a 5-6@xbeg-delay 6-6@xbegin 7-6@xend 8-6@105c41 9-5@105c41 10-4@1001042 11-4@1001042 12-4@105c41 13-3@1028e9
% #[EXPLORE]       from #13/tid3, chose tid 6 (xabort injection), child of #6/tid6
%  0--1@105c41 1-4@105c41 2-5@xbeg-delay 3-5@xbegin 4-5@datarace-a 5-6@xbeg-delay 6-6@xbegin 7-6@datarace-a 8-6@105c41 9-5@105c41 10-4@1001042 11-4@1001042 12-4@105c41 13-3@1028e9
% #[EXPLORE]       from #13/tid3, chose tid 6, child of #2/tid5
% #[D-MAIL]        tt'ing abort set: {TID6 (1,0), TID5 (1,0)}
%  0--1@105c41 1-4@105c41 2-5@xbeg-delay 3-6@xbeg-delay 4-6@xbegin 5-6@xend 6-6@105c41 7-5@xbegin 8-5@xend 9-5@105c41 10-4@1001042 11-4@1001042 12-4@105c41 13-3@1028e9
% #[EXPLORE]       from #13/tid3, chose tid 6, child of #2/tid5
% #[D-MAIL]        tt'ing abort set: {TID6 (0,1), TID5 (1,0)}
% +0--1@105c41 1-4@105c41 2-5@xbeg-delay 3-6@xbeg-delay 4-6@xbegin 5-6@xend 6-6@105c41 7-5@xbegin 8-5@datarace-a 9-5@105c41 10-4@1001042 11-4@1001042 12-4@105c41 13-3@1028e9
%  0--1@105c41 1-4@105c41 2-5@xbeg-delay 3-6@xbeg-delay 4-6@xbegin 5-6@datarace-a 6-6@105c41 7-5@xbegin 8-5@xend 9-5@105c41 10-4@1001042 11-4@1001042 12-4@105c41 13-3@1028e9
% +0--1@105c41 1-4@105c41 2-5@xbeg-delay 3-6@xbeg-delay 4-6@xbegin 5-6@datarace-a 6-6@105c41 7-5@xbegin 8-5@datarace-a 9-5@105c41 10-4@1001042 11-4@1001042 12-4@105c41 13-3@1028e9

% --- fig-retry   2018-08-03 16:41:45.910879780 -0400
% +++ fig-baseline        2018-08-03 16:42:02.822879546 -0400
% @@ -16 +110 @@
%  7-5@xbegin 8-5@xend 9-5@105c41 10-6@105c41 11-4@105c41 12-6@105c41 13-4@105c41 14-6@1008178 15-6@xbegin 16-6@xend 17-6@105c41 18-3@1028e9
%  7-5@xbegin 8-5@xend 9-5@105c41 10-6@105c41 11-4@105c41 12-6@105c41 13-4@105c41 14-6@1008178 15-6@xbegin 16-6@datarace-a 17-6@datarace 18-6@105c41 19-3@1028e9
% #[EXPLORE]       dbg #6/tid5 retry set: {TID6 (0,1), TID5 (1,0)}
% #[EXPLORE]       from #19/tid3, chose tid 5 (xbegin failure injection), child of #7/tid5
%  7-5@xbegin 8-5@datarace-a 9-5@datarace 10-5@105c41 11-6@105c41 12-4@105c41 13-6@105c41 14-4@105c41 15-6@1008178 16-6@xbegin 17-6@xend 18-6@105c41 19-3@1028e9
% #[EXPLORE]       dbg #9/tid5 retry set: {TID6 (1,0), TID5 (1,1)}
% #[EXPLORE]       dbg #6/tid5 retry set: {TID6 (0,1), TID5 (1,0)}
% #[EXPLORE]       from #19/tid3, chose tid 6 (xbegin failure injection), child of #16/tid6
%  7-5@xbegin 8-5@datarace-a 9-5@datarace 10-5@105c41 11-6@105c41 12-4@105c41 13-6@105c41 14-4@105c41 15-6@1008178 16-6@xbegin 17-6@datarace-a 18-6@datarace 19-6@105c41 20-3@1028e9
% #[EXPLORE]       dbg #9/tid5 retry set: {TID6 (1,0), TID5 (1,1)}
% #[EXPLORE]       dbg #6/tid5 retry set: {TID6 (0,1), TID5 (1,0)}
% #[EXPLORE]       from #20/tid3, chose tid 6, child of #9/tid5
% #[D-MAIL]        tt'ing abort set: {TID6 (1,0), TID5 (1,1)}
%  7-5@xbegin 8-5@datarace-a 9-5@datarace 10-6@105c41 11-4@105c41 12-6@105c41 13-4@105c41 14-6@1008178 15-6@xbegin 16-6@xend 17-6@105c41 18-5@105c41 19-3@1028e9
% #[EXPLORE]       dbg #6/tid5 retry set: {TID6 (0,1), TID5 (1,0)}
% #[EXPLORE]       from #19/tid3, chose tid 6, child of #6/tid5
% #[D-MAIL]        tt'ing abort set: {TID6 (0,1), TID5 (1,0)}
% +7-5@xbegin 8-5@datarace-a 9-5@datarace 10-6@105c41 11-4@105c41 12-6@105c41 13-4@105c41 14-6@1008178 15-6@xbegin 16-6@datarace-a 17-6@datarace 18-6@105c41 19-5@105c41 20-3@1028e9
% +7-6@105c41 8-4@105c41 9-6@105c41 10-4@105c41 11-6@1008178 12-6@xbegin 13-6@xend 14-6@105c41 15-5@xbegin 16-5@xend 17-5@105c41 18-3@1028e9
% +7-6@105c41 8-4@105c41 9-6@105c41 10-4@105c41 11-6@1008178 12-6@xbegin 13-6@xend 14-6@105c41 15-5@xbegin 16-5@datarace-a 17-5@datarace 18-5@105c41 19-3@1028e9
%  7-6@105c41 8-4@105c41 9-6@105c41 10-4@105c41 11-6@1008178 12-6@xbegin 13-6@datarace-a 14-6@datarace 15-6@105c41 16-5@xbegin 17-5@xend 18-5@105c41 19-3@1028e9
% +7-6@105c41 8-4@105c41 9-6@105c41 10-4@105c41 11-6@1008178 12-6@xbegin 13-6@datarace-a 14-6@datarace 15-6@105c41 16-5@xbegin 17-5@datarace-a 18-5@datarace 19-5@105c41 20-3@1028e9


\subsubsection{STM (abort code) reduction}
\label{sec:tm-eval-stm}

Secondly, the state spaces could be reduced
simply by restricting the concurrency model
to only a subset of nondeterministic {\tt xbegin} outcomes possible under HTM.
Concretely speaking,
the {\tt -X -A -S} combination of Landslide options
suppresses retry aborts
(\cref{sec:txn-abort-codes}),
which must be checked at every transactional preemption point,
replacing them with explicit and conflict aborts,
which Landslide injects only after identifying memory conflicts through DPOR
or encountering an {\tt xabort},
respectively,
ultimately simulating the semantics of STM rather than HTM.
This can allow for state space reduction %result in smaller state spaces compared to the baseline
when transactions happen to be non-conflicting,
but more impactfully,
conflict aborts can occur only {\em after} the other thread's conflicting access,
so between a pair of transactions
% only the first two during execution, not true for subsequent pairs, if they are in conflict with previous pairs at all
only the success, success and success, abort sequences need be tested;
abort, success and abort, abort may (in fact, must) be skipped.
The final column in \Cref{tab:tm-verifs} shows the result of STM semantics verification,
which always results in at least 2x reduction compared to the baseline,
although retry sets can make up some of the lost ground in some cases.

\subsubsection{Abstraction reduction}
\label{sec:tm-abstraction}

Visual inspection of the AVL tree and separate-chaining map implementations \cite{tm-benchmark-cmu},
after correcting the former's atomicity bug (\Cref{fig:avlbug}),
reveals that every use of HTM followed exactly the same pattern:
running identical data structure logic in both the transactional and abort paths,
as though HTM were merely a mutual exclusion lock with fancy performance characteristics.
% this citation is a lil problematic
Prior work \cite{dbug-phdthesis} proposed {\em abstraction reduction},
in which the user identifies program components that can be separated by a well-understood API,
then tests each one against the API individually,
effectively turning multiplicative state space size factors into additive ones.

In this case, I split the lock-like HTM use
and the mutually-exclusive data structure code
into separate tests,
{\tt lock}, which checks that the use of HTM guarantees mutual exclusion,
and {\tt avl\_mutex}/{\tt map\_mutex},
which replace the open-coded HTM use with an already-trusted P2 mutex.
{\tt lock\_fast},
a bonus test,
checks the transactional lock's performance by asserting that
its internal logic won't trigger conflict aborts
even when the client's accesses are independent.
\Cref{fig:htm-lock} shows their core logic.
Finally, I parameterized them over how the lock was implemented:
a real-world {\tt spinlock} implementation from \cite{spinlock-rtm-github},
{\tt spin\_fixed}, the same with the performance bug from \cref{sec:tm-eval-bugs} fixed,
and {\tt mutex}, using Landslide-annotated P2 mutexes (as the AVL and map implementations do).

\begin{figure}[h]
	\begin{center}
	\begin{tabular}{p{0.47\textwidth}p{0.5\textwidth}}
		\begin{tabular}{l}
			\texttt{\ctype{static int} num\_in\_section = 0;} \\
			\texttt{\flow{for} (\ctype{int} i = \const{0}; i < \const{NITERS}; i++) \{} \\
			\texttt{~~~~\call{rtm\_spinlock\_acquire}(\&lock);} \\
			\texttt{~~~~num\_in\_section++;} \\
			\texttt{~~~~\flow{if} (!\call{\_xtest()})} \\
			\texttt{~~~~~~~~\call{thr\_yield}(\const{-1});} \\
			\texttt{~~~~\flow{assert}{}(num\_in\_section == \const{1}); } \\
			\texttt{~~~~num\_in\_section--;} \\
			\texttt{~~~~\call{rtm\_spinlock\_release}(\&lock);} \\
			\texttt{\}} \\
		\end{tabular}
		&
		\begin{tabular}{ll}
			\texttt{\flow{for} (\ctype{int} i = \const{0}; i < \const{NITERS}; i++) \{} \\
			\texttt{~~~~\call{rtm\_spinlock\_acquire}(\&lock);} \\
			\texttt{~~~~\flow{assert}{}(\call{\_xtest}());} \\
			\texttt{~~~~\call{rtm\_spinlock\_release}(\&lock);} \\
			\texttt{\}} \\
		\end{tabular}
		\\
		\\
		(a) {\tt lock}(), tests mutual exclusion.
		&
		(b) {\tt lock\_fast}(), tests for no spurious aborts.
	\end{tabular}
	\end{center}
	\caption{Abstraction reduction test cases.}
	\label{fig:htm-lock}
\end{figure}

\Cref{tab:tm-verifs2} shows the new resulting levels of verification
Landslide reached before the same 10-hour timeout.
Provided that one trusts the {\tt lock} tests correctly check the desired properties,
and that open-coding hadn't introduced any new bugs (such as \Cref{fig:avlbug}'s),
the benefit is clear:
the data structure tests' state spaces become much more tractable,
%the data structure tests suddenly reach unprecedented thread/iteration counts with ease,
% 'so on' here includes free/re-malloc mem conflicts; see comment in avl_mutex.cpp
their state space growth now defined only by internal conflicts from tree rebalancing, map resizing, and so on.
In total,
summing the testing times of {\tt lock}({\tt mutex})$(K,N)$ and {\tt avl\_mutex}$(K,N)$
produces the same verification as {\tt avl\_fixed}$(K,N)$ far more cheaply.
Furthermore, {\tt lock}'s verification can be reused,
whereas {\tt avl\_fixed} and {\tt map\_basic} effectively duplicated the mutex verification between them.

\begin{table}[t]
	\begin{center}
		\footnotesize
		\begin{tabular}{cc}
		\begin{tabular}{cc||r|r}
			& & \multicolumn{2}{c}{\bf STM ({\tt -A -S})} \\
			& & \cpu{\bf cpu (s)} & \ints{\bf SS size} \\
			\bf test & \bf K,N  & \em (or \ETAdag{\bf \em est.})
			                      & \em (or \ETAdag{\bf \em est.}) \\
			\hline
			\hline
			{\tt lock}
			& 2,1 & \cpu{3.41} & \ints{4} \\
			({\tt spinlock})
			& 2,2 & \cpu{198.77} & \ints{1702} \\
			% if lock(spin_fixed) will time out on 2,3 and 3,2, this is guaranteed to aswell
			& 3,1 & \cpu{35.28} & \ints{246} \\
			\hline
			{\tt lock}
			& 2,1 & \cpu{3.55} & \ints{4} \\
			({\tt spin\_fixed})
			& 2,2 & \cpu{105.57} & \ints{998} \\
			& 2,3 & \ETAdag{33h 13m} & \ETAdag{321553} \\
			% [JOB 4] progress: 138911/321553 brs (43.200015%), ETA 23h 12m 50s (elapsed 9h 59m 48s)
			% total CPU time consumed: 10h 4s (36004577784 usecs) (core saturation: 12%)
			& 3,1 & \cpu{28.27} & \ints{186} \\
			& 3,2 & \ETAdag{13y 281d} & \ETAdag{1443676} \\
			% [JOB 4] progress: 169893/1443676 brs (11.768080%), ETA 13y 280d 17h 30m 54s (elapsed 9h 59m 48s)
			& 4,1 & \ETAdag{16h 26m} & \ETAdag{432628} \\ % was: 30408, 282084 pre mkrun
			% [JOB 4] progress: 286910/432628 brs (66.317896%), ETA 6h 26m 25s (elapsed 9h 59m 48s)
			% total CPU time consumed: 10h 4s (36004288982 usecs) (core saturation: 12%)
			\hline
			{\tt lock}
			& 2,1 & \cpu{3.44} & \ints{4}		\\
			({\tt mutex})
			& 2,2 & \cpu{16.83} & \ints{180}	\\
			& 2,3 & \cpu{405.21} & \ints{9372}	\\
			& 2,4 & \cpu{24999.68} & \ints{489480}	\\ % was: 15266 sec 384512 pre mkrun
			% total CPU time consumed: 6h 56m 39s (24999680839 usecs) (core saturation: 12%)
			& 3,1 & \cpu{15.00} & \ints{132}	\\
			%& 3,2 & \cpu{29432.87} & \ints{707086}	\\ % pre mkrun fix (let alone bottom1)
			& 3,2 & \ETAdag{26h 42m} & \ETAdag{1223955} \\
			% [JOB 2] progress: 714882/1223955 brs (58.407517%), ETA 16h 41m 58s (elapsed 9h 59m 55s)
			% total CPU time consumed: 10h 3s (36003621109 usecs) (core saturation: 12%)
			& 4,1 & \cpu{665.89} & \ints{15064}	\\
			\hline
			{\tt lock\_fast}
			& 2,1 & \cpu{3.25} & \ints{1}	\\
			({\tt spin\_fixed})
			& 9,9 & \cpu{4.61} & \ints{1}	\\
			\hline
			{\tt lock\_fast}
			& 2,1 & \cpu{3.19} & \ints{1}	\\
			({\tt mutex})
			& 9,9 & \cpu{4.62} & \ints{1}	\\
			\end{tabular}
			&
			\begin{tabular}{cc||r|r}
			& & \multicolumn{2}{c}{\bf Non-transactional} \\
			& & \cpu{\bf cpu (s)} & \ints{\bf SS size} \\
			\bf test & \bf K,N & \em (or \ETAdag{\bf \em ETA}) & \em (or \ETAdag{\bf \em est.}) \\
			\hline
			\hline
			{\tt avl\_mutex}
			& 2,1 & \cpu{3.48} & \ints{7} \\
			& 2,2 & \cpu{6.25} & \ints{85} \\
			& 2,3 & \cpu{24.69} & \ints{561} \\
			& 2,4 & \cpu{217.98} & \ints{4984} \\
			& 3,1 & \cpu{8.26} & \ints{129} \\
			& 3,2 & \cpu{1403.46} & \ints{30653} \\
			& 3,3 & \ETAdag{qqq} & \ETAdag{qqq} \\
			& 4,1 & \cpu{199.96} & \ints{4488} \\
			& 4,2 & \ETAdag{qqq} & \ETAdag{qqq} \\
			\cline{1-4}
			{\tt map\_mutex}
			& 2,1 & \cpu{39.81} & \ints{83} \\
			& 2,2 & \ETAdag{26h 21m} & \ETAdag{1085126} \\
			% [JOB 9] progress: 496530/1085126 brs (45.757812%), ETA 16h 21m 0s (elapsed 9h 59m 38s)
			% total CPU time consumed: 10h 19s (36019507965 usecs) (core saturation: 12%)
			& 3,1 & \ETAdag{173d 17h} & \ETAdag{12572187} \\
			% [JOB 9] progress: 496281/12572187 brs (3.947451%), ETA 173d 6h 38m 55s (elapsed 9h 59m 43s)
			% total CPU time consumed: 10h 20s (36020148608 usecs) (core saturation: 12%)
			\cline{1-4}
			{\tt map\_basicer}
			& 2,1 & \cpu{14.34} & \ints{9} \\
			{\tt \_mutex}
			& 2,2 & \cpu{147.27} & \ints{2953} \\
			& 2,3 & \cpu{12946.21} & \ints{244691} \\
			& 2,4 & \ETAdag{2d 17h} & \ETAdag{1130184} \\
			% [JOB 3] progress: 572848/1130184 brs (50.686221%), ETA 2d 6h 52m 20s (elapsed 9h 59m 55s)
			& 3,1 & \cpu{637.81} & \ints{12707} \\
			& 3,2 & \ETAdag{102d 13h} & \ETAdag{13220616} \\
			% [JOB 3] progress: 560711/13220616 brs (4.241187%), ETA 102d 2h 56m 44s (elapsed 9h 59m 55s)
			& 4,1 & \ETAdag{33h 33m} & \ETAdag{1778661} \\
			% [JOB 3] progress: 601402/1778661 brs (33.812059%), ETA 23h 33m 16s (elapsed 9h 59m 55s)
			\\
		\end{tabular}
		\\
		\\
		(a) Verifying HTM locks alone.
		&
		(b) Verifying the lock's client code.
		\end{tabular}
	\end{center}
	\caption[Continuation of \Cref{tab:tm-verifs}, demonstrating abstraction reduction.]
		{Continuation of \Cref{tab:tm-verifs},
		% XXX: if you put this citation, it will cause "(document)" to appear in the list of back-references
		% in the corresponding bibliography entry, rather than an appropriate section number
		demonstrating abstraction reduction %\cite{dbug-phdthesis}
		on the {\tt avl\_fixed} and {\tt map\_basic} tests
		by verifying HTM mutex implementations separately.
		Tested with STM semantics, as both lock implementations include a retry loop.}
	\label{tab:tm-verifs2}
\end{table}

Note two curiosities:
firstly,
the impact that fixing \Cref{fig:spinlockbug}'s performance bug
(changing a read+write to a read only)
had on even the correctness tests:
{\tt lock}({\tt spin\_fixed})'s state spaces were reduced by nearly half compared to {\tt lock}({\tt spinlock}),
on account of DPOR no longer needing to reorder the (now) read-read access pairs.
Secondly, {\tt spinlock} and {\tt spin\_fixed} take longer to test per interleaving
%(CPU time divided by state space size)
than {\tt mutex}
% old ver.
%(roughly 9 interleavings per second for the former, 22 for the latter),
% lock spin fixed:
% (4+998+138911+186+169893+286910) / (355+10557+3600000+2827+3600000+3600000) = 5.51
% lock mutex:
% (4+180+9372+489480+132+714882+15064) / (344+1683+40521+2499968+1500+3600000+66589) = 19.79
(roughly 6 interleavings per second for the former, 20 for the latter),
because while {\tt mutex} abstracts away threads needing to wait their turn for the critical section
behind an API Landslide understands,
the spinlock's wait loop is open-coded, and Landslide must fall back on its costlier heuristic synchronization detection
(\cref{sec:landslide-blocking}).
In this way (and also, of course, because {\tt mutex}'s state spaces are smaller overall),
{\tt mutex} can be seen in turn as a further abstraction reduction of {\tt spinlock}.

% commentary -- writing this test case well requires expert knowledge of noobs (of writing landslide friendly tests)
% - placing the asserts well, using sync_test_and_set and asserting on the result rather than having a separate read,
%   has an impact on the state space, as seen above
% - design of critical_section and how it uses xtest to yield only in failure path, and why that's justified

% other one to test -
% https://github.com/vamsikc/leveldb-tsx/blob/master/ext/xsync/include/scope.hpp

% "Although I skipped implementing it to begin with, the sheer size of this reduction also
% reflects on the numbers one might see betwee the original reduction in theorem so-and-so that lets us
% treat transactional success paths as mutually-exclusive (and, once started, abort-proof) to begin with"
% this sentence is a complete mess....

%%%%%%%%%%%%%%%%%%%%%%%%%%%%%%%%%%%%%%%%%%%%%%%%%%%%%%%%%%%%%%%%%%%%%%%%%%%%%%%%

\section{Discussion}
\label{sec:tm-discussion}

In this section I review some of the evaluation's results in a broader context,
list the current limitations of Landslide's implementation,
and discuss open problems for future work.

\subsubsection{Retry set optimality}

For all the reduction retry sets demonstrated in \Cref{tab:tm-verifs},
some inefficiencies remain in its strategy.
For example, it is not clear how to prune soundly
when three or more threads must be reorderd around one transactional preemption point,
or when a second pair of partially-independent transactions interleaves while an existing retry set is already active.
Accordingly, I implemented the optimization as conservatively as possible in these cases,
``saturating'' the retry sets to fall back to no pruning
({\tt update\_hax\_abort\_set()} and
{\tt update\_hax\_abandon\_abort\_set()}, respectively).

Likewise, the cases of {\tt htm2}(2,3), (2,4), and (4,1),
in which retry sets achieved better reduction than STM mode,
\revisionminor{suggest} that the latter does not necessarily subsume the former,
and that combining the two could in theory achieve further reduction still,
However, {\tt xbegin} results other than {\tt \_XABORT\_RETRY}
may depend on the execution logic
(explicit aborts may be conditional on some change by a conflicting thread,
and conflict aborts cannot occur before their conflicting transaction to begin with%
\footnote{See {\tt 410user/progs/htm\_causality.c}; that was a fun realization to have
already halfway into implementing retry sets the wrong way at first.}%
),
so \revisionminor{it remains an open problem}
to soundly prune either success paths or retry aborts
while other abort codes are in play,
\revisionminor{and preserve the reduction achieved with the baseline configuration}.

While
%the basic idea is
motivated by straightforward analogy to the known-sound sleep sets,
the intersection of retry sets with Landslide's other exploration features may cause unforeseen problems.
For now, its use is prohibited in conjunction with other state-space-affecting features
such as ICB (\cref{sec:landslide-icb})
as well as multiple abort codes.
I personally believe retry set reduction to be sound under these restrictions,
having carefully scrutinized its behaviour
while constructing \Cref{fig:retry-set-example}
and from inspecting state spaces arising from larger test parameters as well; %while implementing it;
nevertheless, this falls well short of formal proof, which I must leave to future work.

\subsubsection{STM reduction soundness}

In \cref{sec:tm-eval-stm} I showed that state spaces could be reduced even further than with retry sets
by assuming an HTM interface which abstracts away {\tt \_XABORT\_RETRY} behind a loop.
However, suppressing retry aborts is not guaranteed to faithfully test all possible behaviours
observable under HTM.
As an example, note in \Cref{tab:tm-verifs}
how STM semantics reduced % oops
{\tt fig63}'s state space on all $(K,N)$ configurations to 1.
Because its transactional paths are all mutually independent,
DPOR identifies no need either to inject conflict aborts or to reorder threads.
However, this skips the slow-path consistency assertion completely.
If the programmer had intended it to run ``every so often'' at the whim of the timer interrupt,
applying this reduction would be unsound.
Also note the state space size of 4 for many (2,1) test configurations,
corresponding exactly to the aforementioned success, success and success, abort sequences
(times two ways to interleave the two threads).
Because of the scheduling dependency for conflict aborts,
Landslide cannot recognize the failure path's data races
without a third freely-reorderable iteration;
STM mode must be run with $K \times N \ge 3$ to meaningfully test conflicts between failure and success paths at all.
A user wishing to distinguish conflict aborts, retry aborts, and so on during testing
without glossing over any of HTM's peculiarities
could supply the {\tt -X -A} options without {\tt -S};
however, this will inevitably result in state spaces at least as large as the baseline.

On the other hand, some programs may clearly annotate their intention
for abort paths to be executed only in case of actual memory conflicts.
{\tt swap}, {\tt avl\_insert}, and {\tt map\_basic} abstract their {\tt \_xbegin()} calls
behind an interface which can be implemented either with or without retry loops,
while {\tt lock}({\tt spinlock}) and {\tt lock}({\tt mutex}) implement the retry loop directly.
In these cases, the user can assure herself of STM reduction's soundness by visual inspection.
In fact, Landslide's current implementation gets stuck in infinitely deep interleavings
whenever it encounters a retry loop
% not actually sure why, as it should be around-a-pp and not on the 1st branch...
(bypassing even its heuristic infinite loop detection),
so for now the user must inspect the test case to determine which testing mode to use.
Future work could automatically identify a program's retry loops
and give up on {\tt \_XABORT\_RETRY} by switching to STM mode on-the-fly,
much like Landslide's heuristic synchronization detection does for yield loops (\cref{sec:landslide-blocking}).

\subsubsection{Nested transactions}
\label{sec:tm-warpzone-nested}

Whenever {\tt \_xbegin()} is called with a transaction already active, or {\tt \_xend()} while not,
Landslide's current implementation immediately stops and reports a bug.
However, just as concurrent programs often hold multiple locks simultaneously,
one may wish to conduct multiple transactional routines simultaneously,
especially as they may be abstracted across different code modules as a project grows in scale.
Recent work \cite{hybrid-htm-stm,relaxed-transactions-popl} has developed
both implementations and formal semantics for executing nested transactions,
so future work should extend the verification concurrency model to permit such programs.
For now, the best Landslide can offer is to check the transactional components and their client code
separately against their APIs
with abstraction reduction (\cref{sec:tm-abstraction}),
then check the rest of the program with (for example) traditional mutexes that can nest safely.

\subsubsection{Relaxed memory orderings}
\label{sec:tm-warpzone-relaxed}

\cref{sec:tm-design}'s formalization of thread interleavings does not account for read/write reorderings
possible on relaxed consistency architectures \cite{memory-consistency-models}.
In fact,
even after \cite{tm-benchmark-cmu}'s proposed fix to the atomicity protocol in \Cref{fig:htm-fixed},
it is still incorrect on Total Store Order (TSO) architectures such as x86,
let alone on weaker memory models.
Despite stores being totally-ordered, x86 may still reorder stores after subsequent loads
\cite{sully-thesis}.
Accordingly, an execution of lines $8,9a,9b,9c$
may be locally visible to another thread as $9a,8,9b,9c$,
and hence an apparent interleaving of
\[
        \tiat{1},\tjat{1-5},\tiat{7},\underline{\tiat{9a}},\tkat{1-5},\underline{\tiat{8}},\tiat{9b-B}
\]
is possible
(reordered accesses underlined for emphasis).
An acquire barrier is needed between lines 8 and 9 to solve this problem on TSO \cite{tsx-need-barrier}
(on x86, either {\tt mfence} or {\tt xchg}/{\tt xadd}).
% gcc generates the latter for its __sync_lock_test_and_set intrinsic
Recent work \cite{relaxed-transactions-pldi} also demonstrated unsoundness
in a similar lock elision implementation on ARMv8 (PSO),
in which the transactional path reads the lock's internal state directly
rather than using a separate flag.
In \Cref{fig:htm-fixed}, a release barrier before line A is also necessary under PSO.
% On Partial Store Order (PSO) architectures,
% %even more barriers may be necessary.
% a release barrier is necessary before line A as well.
% Recent work \cite{relaxed-transactions-pldi} also discovered similar unsoundness
% in lock elision \cite{speculative-lock-elision}
% on ARMv8 (PSO).
% Recent work \cite{relaxed-transactions-pldi} also discovered similar unsoundness

Because Landslide's concurrency model includes only instruction-level thread nondeterminism,
not per-CPU memory buffer reorderings,
its current HTM implementation cannot find this bug.
In fact, it erroneously verifies the corresponding test {\tt htm2}(3,1)
%% sigbovik version, probably pre-thrlib-function and definitely pre-dpor-preferred-tid and pre-trusted-join
% in 3 CPU-hours,
% with 467730 distinct interleavings in total,
in 40 CPU-seconds,
with 774 interleavings in total,
none of which include the above-listed sequence.
Recent work has extended DPOR to support TSO and PSO memory nondeterminism \cite{tsopso},
as well as proposed formal execution semantics for HTM on these architectures
\cite{relaxed-transactions-popl,relaxed-transactions-pldi};
if both of these advances were incorporated into Landslide's concurrency model,
it could find or verify the absence of such bugs.
Visual inspection of the transactional AVL tree and hashmap
\cite{tm-benchmark-cmu} found no barriers used in this implementation pattern;
I would urge any reader interested in using those to add them in by hand first.
The test case {\tt lock}({\tt mutex}) ({\tt 410user/progs/htm\_mutex.c} in the repository)
provides an example of how to use compiler intrinsics to emit the necessary barriers.

%%%%%%%%%%%%%%%%%%%%%%%%%%%%%%%%%%%%%%%%%%%%%%%%%%%%%%%%%%%%%%%%%%%%%%%%%%%%%%%%

\section{Summary}

\revision{This chapter has extended the concurrency model of stateless model checking to support transactional memory,
and proposed several new reduction strategies for coping with the concomitant % hi william
state space explosion.
I have made the following contributions:

\begin{itemize}
	\item Three equivalence proofs,
		demonstrating that model checkers can soundly emulate strongly-atomic HTM execution semantics by
		treating transactions as globally-uninterruptible atomic sections
		and injecting failure without simulating rollback
		(\cref{sec:tm-design-failure})
		and applying conventional data race analysis with all transactions treated as a global lock
		(\cref{sec:tm-design-locks}),
		and can emulate weakly-atomic STM systems by relaxing the above atomicity to apply to other transactions only
		(\cref{sec:tm-weak-atomicity}).
	\item An implementation of both forms of transactional memory in Landslide,
		with further options to configure the variety of abort codes injected
		(\cref{sec:tm-implementation}).
	\item A suite of transactional benchmark programs,
		designed to produce interesting state spaces of varying size, bugginess, and reducibility
		(\cref{sec:tm-eval-exp-setup}).
		This test suite is more appropriate for evaluating model checkers
		% get rekt STAMP :triumph:
		than prior work, which to date has focused on testing the raw performance of TM backend implementations,
		with little focus on the interactions of abort paths.
	\item Two previously-unknown bugs found in transactional programs from prior work
		(\cref{sec:tm-eval-bugs}).
	\item Retry set reduction,
		a new state space reduction algorithm
		(\cref{sec:tm-retrysets})
		which the evaluation shows identifies and skips equivalent interleavings
		compared to standard DPOR, even when transactional and abort paths are fully mutually conflicting
		(\cref{sec:tm-eval-retrysets}).
	\item Demonstration of two further heuristic reduction strategies on HTM programs,
		abort code reduction (\cref{sec:tm-eval-stm})
		and abstraction reduction (\cref{sec:tm-abstraction})
		which rely on human intuition to ensure their soundness
		to tackle even larger state spaces still.
	\item Verifications of \Cref{fig:htm-fixed}'s atomicity protocol %(with retry set reduction)
		and the {\tt lock}({\tt mutex}) transactional lock %(with abort code reduction)
		up to 2 threads, 4 iterations and 4 threads, 1 iteration for each
		(\cref{sec:tm-verif}).
		%$(K,N)$=$(2,4)$ and $(4,1)$ for each.
		% verified htm atomicity protocol (htm2) and mutex (lock mutex) up to (2,4), (4,1)
\end{itemize}

As concurrent systems grow in complexity to meet modern performance demands,
they grow ever more fraught with opportunity for subtle bugs,
and the programmer likewise faces ever more difficulty in verifying her programs.
Following in the footsteps of related work for relaxed memory concurrency,
this chapter has enabled stateless model checking to address yet another new technique,
and concludes this thesis's contributions to the field.
The upcoming chapters will offer a tour of related research problems,
solved and unsolved,
practical and philosophical,
past and future.

}

\chapter{Related Work}
\label{chap:related}

In this chapter I provide a brief tour of the related work.

\section{History of Stateless Model Checking}

{\bf Hall of Fame.}
Stateless model checking dates back to Verisoft \cite{verisoft}, the 1997 tool which first attempted to exhaustively explore the possibile ways to interleave threads.
Since then, researchers have built many tools along the same lines to test many kinds of programs.
%list of mcs follows
The most popular model checker, according to citation count, is Microsoft Research's CHESS \cite{chess},
a checker for userspace C++ programs which preempts at synchronization APIs.
%which introduced the Iterative Context Bounding technique \cite{chess-icb} for heuristically ordering the search to uncover bugs faster.
%
Other checkers include MaceMC \cite{macemc}, MoDist \cite{modist}, SAMC \cite{samc}, ETA \cite{dbug-retreat},
each of which limit thread communication to the boundaries of a message-passing API;
%
dBug \cite{dbug-ssv}, another CMU original, similar to CHESS;
%
and finally SPIN \cite{spin} and Inspect \cite{inspect} which can preempt at any shared variable access.
%
Perhaps by now, Landslide itself deserves a spot on this list: I first introduced it in my MS thesis \cite{landslide},
and it is most notable for testing kernel-level code through the use of a simulator.

{\bf Search Strategies}
To date a number of techniques have been proposed to mitigate exponential explosion, the Sisyphean rock of stateless model checkers.
Dynamic Partial Order Reduction (DPOR) \cite{dpor}, the baseline approach, prunes redundant interleavings by identifying independent thread transitions which can commute without changing the program state.
It is a {\em sound} reduction algorithm, meaning it will never fail to test a possible program behavior, despite skipping many execution sequences.
DPOR has since been extended in several ways:
a provably optimal version which guarantees to explore no more than one interleaving from each equivalence class \cite{optimal-dpor},
% TODO: a 3rd example?
and an extension of the algorithm to include nondeterminism arising from weak memory models \cite{tsopso}.
Recently, SATCheck and Maximal Causality Reduction have emerged as better-performing alternatives to DPOR.
These algorithms use SMT solvers \cite{z3} to identify additional equivalences by analyzing the values read and written by each memory access.

{\bf Other theoretical advances.}
A number of other techniques address the problem that even with DPOR, a state space may be too large.
Iterative Context Bounding (ICB) \cite{chess-icb} is a search ordering heuristic which prioritizes interleavings with fewer preemptions,
according to the insight that most bugs require fewer preemptions to expose.
Bounded Partial Order Reduction \cite{bpor} adapts DPOR to be soundly compatible with ICB.
Preemption Sealing \cite{sealing} allows programmers to exclude preemption points arising from trusted source code modules,
reducing the state space by limiting the scope of the test.
Probabilistic Concurrency Testing \cite{randomized-scheduler} eschews DPOR's depth-first approach entirely (as well as any potential for completing the state space), randomly sampling a broad cross-section of the state space and providing probabilistic bounds on uncovering bugs.

\section{Educational Testing Tools}

\section{Real-World Techniques}

\chapter{Future Work}
\label{chap:warpzone}
% TODO: ask for permission for this one
% https://twitter.com/Poem4your_sprog/status/878747092788838400
\inspirationalquote{
\begin{tabular}{p{0.51\textwidth}}
A little bit of being kind,
a tiny open door; \\
A nicer slice of peace of mind
begets a little more. \\
It won't correct the world tonight,
nor change tomorrow too; \\
But maybe if you do it right,
you'll find it changes you.
\end{tabular}}
{Sam "Poem\_for\_your\_sprog" Garland}

Each of the previous three chapters concluded by
discussing specific limitations,
listing concrete and immediate ways to address them
with existing techniques
(\sect{\ref{sec:quicksand-discussion}},
\sect{\ref{sec:education-discussion}},
\sect{\ref{sec:tm-discussion}}).
Meanwhile, this chapter takes a broader interpretation of ``future work'',
namely,
how might future Landslides solve research problems I didn't even pose to begin with.
%The reader looking for their own thesis topic would most likely find it here.
%The reader should beware of wild speculation ahead
%which may assume solutions to extremely difficult problems with a wave of the hand.

\subsubsection{User-friendliness}

In the user study survey (\sect{\ref{sec:education-eval-survey}}),
students most often complained
that interpreting Landslide's preemption traces to diagnose and understand their bugs was too difficult.
While the understanding step of course requires human intuition,
there is certainly room to improve the diagnosis step beyond just showing the user one static HTML table.
Related works like Symbiosis \cite{symbiosis} can find a minimal difference between the failing trace and a non-failing one,
which would allow the user to effortlessly pinpoint which preemptions are closest to the true root cause.
Further, using ICB \cite{chess-icb} to show the user a minimum set of preemptions necessary to expose the bug
could help her narrow down possible diagnoses more quickly.
Finally, the preemption trace itself could be interactive,
allowing the user to click and drag to reorder thread transitions and see how the interleaving would change,
or to click and drag preemption points across the source code to figure out how much need to be atomic.

State space size management remains an issue, as ever.
While Quicksand's professed selling point is that the user need input only a CPU budget,
at the same time,
pruning out the overall state space would allow Quicksand to achieve more meaningful results in that same budget.
Theoretical advances in state space reduction come out every year (\sect{\ref{sec:related-algs}}),
but the user's human intuition can also contribute if properly encouraged.
For the P2 tests (\sect{\ref{sec:education-pebbles-tests}}),
I configured Quicksand by hand to issue appropriate {\tt without\_function} commands to Landslide
(\sect{\ref{sec:landslide-pps}}),
and even more still for the HTM tests (\sect{\ref{sec:tm-eval-exp-setup}}),
but a user writing her own tests would have to configure DPOR by hand.
A more mature version would coach the user to decide which functions focus the test on,
using state space estimation to give an idea of expected testing time,
and listing DPOR's memory conflicts to help her identify more candidate functions to potentially ignore.
% there are a LOT of "hmm"s here for future research to address
% most basically like, a modification could have a side effect that would lead to changed behaviour
% arising from a preemption in a *subsequent* function
% so really this can't be limited to just PPs in the changed functions, but dataflow arising therefrom as well
Finally, Landslide could integrate with a version control system to do incremental testing,
automatically analyzing the functions touched by each patch,
and heuristically prioritizing preemption points therein
to quickly check small updates on top of an already-verified codebase.

\subsubsection{Test cases}

In its current form, Landslide is limited to testing only those behaviours that the test case it's hooked up to can generate.
The most obvious limitation of this is resource exhaustion scenarios:
stateless model checking simply cannot handle tests long-running enough to exhaust system memory
(succumbing to exponential explosion long before),
% "involve" here, for subjunctive?
so cannot exercise any flow control that involves {\tt malloc()} failing, for example.
This specific issue could be solved using by-hand annotations to block all preemption points until just before exhaustion,
%(requires model-checking expertise on the programmer's part)
or by extending Landslide to inject allocation failures at {\tt malloc()} callsites (akin to {\tt \_xbegin()}).
%(assumes certain behavioural properties of the allocation library).
However,
these require the user to realize in advance that she should worry about allocation-failure bugs,
and to configure Landslide specifically to target them.
In general, a mature testing tool should not require the user to
know in advance where her bugs might be
before being able to conduct an effective test.

In future work, a stateless model checking framework could collaborate with its user to semi-automatically design better tests.
% i hate the way maple actually does it, in principle, but i need to cite smth here
Concurrency coverage metrics such as proposed by Maple \cite{maple}
% wild speculation
could be extended to capture possible program behaviours under any test input,
not just within the fixed state space under one given test,
and symbolic execution frameworks such as Contessa \cite{contessa}
could search for possible test inputs to suggest changes to the test case which might expose them.
In cases where not-yet-covered conditional branches require
certain function return values,
such as {\tt malloc()} failing,
the tool could offer to add failure injection with the user's approval,
or at least prompt her to write a new test with that as a subgoal.

In the case study of submitted P2 bugs (\sect{\ref{sec:education-eval-bug-case-study}}),
I noticed several submissions of a common pattern
in which Landslide overlooked the {\tt thr\_exit()} use-after-free bug:
they all manually recycled the exited thread's stack using an internal free-list,
rather than calling {\tt free()} or {\tt remove\_pages},
so Landslide's heap-checker failed to see anything out of the ordinary about the subsequent writes to it.
Explicit annotations about the custom free-list's invariants could make Landslide treat it like {\tt free()}
and catch the errant write,
but that begs the question of knowing about the bug in advance.
Catching this bug would require extending the test case to {\tt thr\_fork()} a new thread
to reuse the old stack and suffer data corruption from the stray write.
In this case, the user might happen to find this just by increasing the number of threads/iterations,
which is just good verification principles,
but in general, automatically inferring such data structure invariants
to suggest better assertions
% is this true??
remains an open problem.

% test case design, writing code that actually checks the properties you're interested in
% for example the mutex test, with the if (!xtest) yield check
% then even after you got the right properties, optimizing the state space
% with without_function, ignore_sym, etc, while preserving the guarantee

\subsubsection{Verification}

% soundness concerns - the xbegin/delay-xbegin pp soundness bug popped up in the implementation, despite the soundness thm
% current SoA for this is to manually inspect the state spaces that arize for small K,N and just hope...
% as concurrency model becomes more complex, more opportunities for the implementation to deviate from what
% the soundness theorem actually proves
% how to prove landslide is actually doing what you think? analogous to sel4 spec/impl verif...

% extending landslide's verification for a limited nr of K N to infinite with induction/proof assistance

% - this should prob go in quicksand chapater
% identifying & skipping redundancies in smaller ID state spaces via memoization (same applies to ICB!)

% arugment why ML is NOT appropriate for smc in general -- if you try to train it to choose thread interleavings / search ordering,
% even if you're already in not-enough-cpu-time heuristics land, you're biasing it to find only the types of bugs you've already seen
% and destroying its ability to find novelties
% if your ML provides a partial verification, it is impossible for a human to understand what that verification means

\subsubsection{Heuristics}

In the author's experience,
the WBE algorithm (used in Landslide for estimating number of interleavings) consistently underestimates,
being often seen to count (nearly) monotonically up toward the true value as exploration progresses,
while the RE algorithm (used in Landslide for estimating total execution time) can be unstable and erratic,
bouncing wildly among orders of magnitude even in the space of ten or fewer adjacent interleavings.
These flaws are especially visible in the {\tt avl\_fixed} transactional verification results (Table~\ref{tab:tm-verifs}).
% TODO: simpler approach would be to account for like, "T1->T1" prefix subtree likely to be smaller than "T1->T2" prefix subtree
% because T1->T1->... means T1 has 2 fewer events left to schedule and (M+N-2 choose M-2) tends to be smaller than (M+N-1 choose M-1)
% for roughly similar M/N;
% estimator could account for this somehow... not sure how but
% likewise, estimator could account for retry sets by checking the size of the retry set and multiplying by 1/4 or 1/2 or w/e
% TODO: cite smth about machine learning
% TODO: publish the htm evaluation logs
% state space estimation could probably be drastically improved with machine learning
% accordingly, all the state space estimate logs (ie, quicksand logs) from the tsx chapter are made open-source on so-and-so

% ICB incorporations
% into quicksand - when SS estimates are too high, fall back on a preemption bounded version to find bugs - improve the partial verif
% into education - provide more partial verifs to students (but is it clear they'll understand the subtleties of it?)
% into txn - incorporate failure injection into the model

%%%%

% not gonna bother with this

% in the maximal state space, the "tag all evil ancestors" thing justified in my MS thesis figure 31337
% is not necessary, bc if it's the maximal state space, you know that data race PPs would separate the accesses in C1/C2/C3.
% but, actually figuring out if you're in the maximal state space on-the-fly to justify not tagging?
% (but but, does this even matter, is this an optimization at all?)
% i'm actually not even sure this would be sound even on maximal, because of RMW instructions, their 2 evence are inseparable



\chapter{Conclusion}
\label{chap:conclusion}
%\inspirationalquote{
%\begin{tabular}{p{0.6\textwidth}}
%"This is the end, right? The end of this story?"
%
%"We were always at the end. It's a free play, buddy. Clock's all zeroes."
%\end{tabular}}
%%{Nine \& Ten, 17776}
%{17776}
\inspirationalquote{
	\begin{tabular}{p{0.84\textwidth}}
	"If you are a god, Zeus, as the stories claim, then why did you create evolution?
	Why did you make a world that can only grow through cruelty and pain?"
	\\
	For a moment\inspirationalhyphen{}surely this meant death was near\inspirationalhyphen{}she thought she heard him answer:
	"My child, I was shaped by the gods that came before me, as you were shaped by me.
	The choice I had was between creation and oblivion, life and death.
	And I chose life, because any life is better than no life,
	because as long as there is life, there is hope\inspirationalhyphen{}if not for us, then for some generation to come."
	\\
	"Then how are you a god, if you can offer so little?" she whispered, feeling death creep closer.
	\\
	"I am a god because I take upon myself the burden of creation," the statue replied.
	\\
	"Then we are all gods," Alexandra said, and pushed the button.
	\end{tabular}}
%{Galatea, The Talos Principle}
{The Talos Principle}

Concurrency testing is, in a way, a microcosm for research in general:
an exponentially large problem
that invariably thwarts they whoever try to solve it ``perfectly'';
it bears fruit only to one who can find their own human-oriented meaning in partial results and compromises.
I speak of course literally on one side of this analogy, figuratively on the other.
This thesis has demonstrated such meaning for stateless model checking,
choosing perhaps one of the most challenging demographics of programmers in need of concurrency testing,
the undergraduate operating systems student,
as its target audience.
Along the way, it justified its techniques' applicability to problems beyond CMU's walls as well
by showing positive results at other universities
and by addressing new concurrency models developed by industry in just the last decade.
% also by comparing the preemption bounds of quicksand eval's bugs to prior work in ICB
% and finding them to actually be skewed to DEEPER bugs than in the "real" "world", so take that
As I write this I am also collaborating with 15-410 staff
to integrate Landslide as an official part of the course grading infrastructure,
so in some sense the na\"{i}ve wish to make an ``automatic TA'' that drew my attention to grad school to begin with
has also been fulfilled.

For as long as people have written science fiction books,
we have fantasized about robots that can make any complex intuitive decision a human can
in addition to the superhuman arithmetic ability that comes standard with silicon.
Recent trends in machine learning have pursued this aesthetic,
using statistics and pattern recognition to provide convincing output even on never-before-seen inputs,
but few such systems are designed to also alert the user when their output confidence is low.
% appropriate for serving ads (penalty for being wrong -- no click)
% appropriate for drug synthesis (clinical trial vets the output)
% not appropriate for... govt surveillance o.o;
% not appropriate for program verification
As tempting as it is to fantasize an oracle capable of telling for certain whether your program is free of concurrency bugs,
that would be a disservice to both user and machine:
the human skill that artificial intelligence should really try to emulate is
knowing its limits and asking for help when it gets confused.
% ways humans cooperate to make landslide better:
% - test case design incl. preemption sealing
% - abstraction reduction
% - manual inspection of data races

It's become fashionable among my peers lately,
whenever a new exploit is announced in the news,
to declare that security is doomed,
that computers were a mistake,
that we should all retire and become llama farmers,
and so on.
There is no doubt that we are burdened with massive technical debt from the mistakes made
during computer science's infancy.
Nevertheless,
safety-oriented programming languages and type systems are growing ever more popular,
and meanwhile formal verification techniques to check existing code,
written in the old unsafe ways, grow ever more powerful.
Landslide's ability to reach impressionable student minds is my little contribution in this long battle.
I have hope that we may one day live in peace with our silicon creations,
understanding and respecting their limitations just as they complement ours.
This thesis has been a profoundly meaningful journey for me and I hope you enjoyed reading it.

% TODO TODO TODO
% update this before you turn in duh
% log in to office machine and be like, ./graph.py --out graph.svg logshort.csv
% where logshort is log.csv but you did 'uniq' on it so the output datapoint set is not humongous
\newpage
\thispagestyle{empty}
\begin{center}
\begin{tabular}{c}
\vspace{12em} \\
\includegraphics[width=\textwidth]{pages.pdf}
\end{tabular}
\end{center}

%\appendix
%\include{appendix}

\backmatter

%\renewcommand{\baselinestretch}{1.0}\normalsize

% By default \bibsection is \chapter*, but we really want this to show
% up in the table of contents and pdf bookmarks.
\renewcommand{\bibsection}{\chapter{\bibname}
\inspirationalquote{Only a doctor of philosophy, Darth.}{Robert Marsh}
}
%\renewcommand{\bibpreamble}{This text goes between the ``Bibliography''
%  header and the actual list of references}
\bibliographystyle{alpha}
\bibliography{citations}

\end{document}
